% Useful variables
\newcommand{\DocMainTitle}{SmallFloats.jl}
\newcommand{\DocVersion}{0.1.0}
\newcommand{\DocAuthors}{Jeffrey Sarnoff}
\newcommand{\JuliaVersion}{1.12.6}

% ---- Insert preamble
\input{preamble.tex}


\part{Home}


\chapter{SmallFloats.jl}



\label{10656431308105265837}{}


\emph{A conforming, performance-oriented Julia implementation of the IEEE P3109 draft standard — arithmetic formats for machine learning at bitwidths 3–8.}



Bit-exact defined results on every default path; one projection engine as the single write path into a code point; approximation only behind an explicitly named, exhaustively measured κ registry; ≈ 8.9 million test assertions; table-gather kernels at fractions of a nanosecond per element.




\begin{minted}{jlcon}
julia> using SmallFloats

julia> Binary8p4se(1.6) + Binary8p4se(0.25)
Binary8p4se(1.875 ≡ 0x47)

julia> Binary5p3sf(0x08) == Binary5p3sf(1.0)     # UInt8 constructs from a code point
true
\end{minted}



\section{Documentation map}



\label{558805016199603468}{}


\begin{itemize}
\item \textbf{\hyperlinkref{16879307210612674453}{Introduction}} — what the package is, the design pillars, a thirty-second tour, installation.


\item \textbf{\hyperlinkref{17671780375889550826}{Cheat Sheet}} — one-page lookup for format names, conversion, projections, operations, arrays, blocks, packed storage, and common traps.


\item \textbf{\hyperlinkref{10272851679537904033}{User Guide}} — the complete public API in usage order: formats, values, projection specifications, the two operation registers, arrays and sorting, blocks, packed storage, conformance and κ, performance guidance.


\item \textbf{\hyperlinkref{7866567007412100191}{User Examples}} — runnable examples in three tiers: basic, machine learning, deep learning.


\item \textbf{\hyperlinkref{10344462847155476790}{Julia Compatibility}} — the Base register: which Base functions work on \texttt{Binary} values, what they map to, and what deliberately isn{\textquotesingle}t mapped.


\item \textbf{\hyperlinkref{4463396622184713734}{Technical Guide}} — internals: the encoding and projection engine, the oracle{\textquotesingle}s rigor classes, tables and kernels, the block layer{\textquotesingle}s exactness filters, the verification and benchmark doctrines.


\item \textbf{\hyperlinkref{4476403263922386643}{Technical Examples}} — internals-level recipes: pipeline introspection, exhaustive verification of custom code, κ measurement, doctrine-compliant benchmarking.


\item \textbf{\hyperlinkref{5861982930757625708}{Adding Operations}} — how to add new scalar (unary through quaternary) and block (scaled, elementwise, reductive) operations with every guarantee intact: registry mechanics, ω-semantics duties, worked examples.


\item \textbf{\hyperlinkref{5597696625397351643}{External Reference}} — docstrings for the documented public surface.


\item \textbf{\hyperlinkref{2804742620609579687}{Internal Reference}} — docstrings for the documented unexported machinery.

\end{itemize}


\part{Introduction}


\chapter{Introduction}



\label{16879307210612674453}{}


SmallFloats.jl is a conforming, performance-oriented Julia implementation of the \textbf{IEEE P3109 draft standard} — \emph{Arithmetic Formats for Machine Learning} — covering every binary format the draft defines at bitwidths 3 through 8, every projection (rounding × saturation) mode including the three stochastic families, the full scalar operation catalog, block-scaled ({\textquotedbl}MX-style{\textquotedbl}) operations, and the draft{\textquotesingle}s conformance and κ-approximation machinery.



\section{Why this package exists}



\label{13660426995343544750}{}


Small floating-point formats are easy to implement \emph{approximately} and surprisingly hard to implement \emph{exactly}. An FP8 \texttt{Exp} differs from a bit-exact one only in a handful of code points near rounding boundaries — invisible in a demo, decisive in a conformance suite, and quietly corrosive in research that compares number formats. SmallFloats.jl takes the position that for value sets this small there is no excuse for approximation:



\begin{itemize}
\item \textbf{Bit-exact defined results on every default path.} Every operation{\textquotesingle}s result is the correct projection of the mathematically exact value, established through a rigorous oracle (exact arithmetic, correctly-rounded \texttt{Float128} where IEEE mandates it, and MPFR directed-rounding enclosures with automatic precision escalation everywhere else).


\item \textbf{One projection engine.} \texttt{RoundToPrecision → Saturate → Encode} is the \emph{single} write path into a code point. There is no second, {\textquotedbl}fast but slightly different{\textquotedbl} rounding routine anywhere in the package.


\item \textbf{Approximation is opt-in, named, and measured.} Faster-but-inexact implementations live behind an explicit registry (\hyperlinkref{3878213895374713375}{\texttt{register\_approx!}}); each registration{\textquotesingle}s deviation bound κ is \emph{measured by exhaustive enumeration} at registration time, and understated declarations are rejected. Nothing approximate is reachable from the default API.


\item \textbf{Exhaustively tested.} The value sets are tiny (≤ 256 points per format), so the test suite does not sample — it enumerates: every operation on every input, every ordering of every pair, every stochastic draw at boundary budgets. The shipped suite carries ≈ 8.9 million assertions.


\item \textbf{Fast where it matters.} Pure-mode elementwise work runs through precomputed lookup tables (sub-nanosecond per element); the scalar path is fully specialized and allocation-free (≈ 18 ns for a complete \texttt{Add} including projection); sub-byte packed storage, integer-keyed ordering with an O(n) counting sort, and a mask-based rounding core round out the performance story. A reproducible Chairmarks benchmark suite ships in \texttt{benchmarking/}.

\end{itemize}


\section{Thirty-second tour}



\label{18023860074456033423}{}



\begin{minted}{jlcon}
julia> using SmallFloats

julia> x = Binary8p4se(1.6)          # construct = project under the default spec
Binary8p4se(1.625 ≡ 0x45)

julia> x + Binary8p4se(0.25)         # Base operators use the format's default spec
Binary8p4se(1.875 ≡ 0x47)

julia> Add(Binary8p4se, RTP_SatNone, x, Binary8p4se(0.25))
Binary8p4se(1.875 ≡ 0x47)            # spec-named register: any mode, explicitly

julia> σ = RSA_SatNone();            # stochastic rounding, default N = 8 bits

julia> Add(Binary8p4se, σ, Binary8p4se(2.0), Binary8p4se(0.03125); R = 255)
Binary8p4se(2.25 ≡ 0x49)             # explicit draw R makes it reproducible

julia> Exp(Binary8p4se, RNE_SatNone, Binary8p4se.([-1.0, 0.5, 2.0]))
3-element Vector{Binary8p4se}:       # array methods route through 256-byte tables
 Binary8p4se(0.375 ≡ 0x34)
 Binary8p4se(1.625 ≡ 0x45)
 Binary8p4se(7.5 ≡ 0x57)
\end{minted}



\section{The formats}



\label{13054678588546481008}{}


A format is the type \texttt{Binary\{K,P,SGN,EXT\}}: bitwidth \texttt{K ∈ 3:8}, precision \texttt{P ∈ 1:K} (with the signedness constraint), signed/unsigned \texttt{SGN}, and extended (has ±Inf) or finite \texttt{EXT}. All 120 draft formats carry their draft names as exported aliases — \texttt{Binary8p4se} is \texttt{Binary\{8,4,true,true\}} ({\textquotedbl}K = 8, P = 4, signed, extended{\textquotedbl}); the trailing letters are \texttt{s}/\texttt{u} for signedness and \texttt{e}/\texttt{f} for extended/finite. Values are 1-byte immutable wrappers around their code point; every format has a single NaN, no negative zero, and \texttt{Float64} serves as the exact interchange carrier for all datums.



\section{Where to go next}



\label{9854230807696872180}{}


\begin{itemize}
\item \textbf{\hyperlinkref{10272851679537904033}{User Guide}} — the complete tour of the public API: formats and values, projection specifications, the two operation registers, arrays and sorting, blocks and scaled operations, packed storage, conversion, conformance, and performance guidance.


\item \textbf{\hyperlinkref{7866567007412100191}{User Examples}} — worked, runnable examples in three tiers: basic usage, machine-learning quantization workflows, and deep-learning block/kernel patterns.


\item \textbf{\hyperlinkref{4463396622184713734}{Technical Guide}} — the internals: the encoding and projection engine, the oracle{\textquotesingle}s evaluation protocol and its two rigor classes, tables and kernels, the block layer{\textquotesingle}s exactness filters, and the testing and benchmarking doctrine.


\item \textbf{\hyperlinkref{4476403263922386643}{Technical Examples}} — internals-level recipes: pipeline introspection, verifying custom kernels against the oracle, κ measurement, differential builds, and doctrine-compliant benchmarking.

\end{itemize}


\section{Installation}



\label{16848947351613073531}{}


The package is not yet registered, so install it by URL:




\begin{minted}{julia}
pkg> add https://github.com/JeffreySarnoff/SmallFloats.jl
\end{minted}



Append \texttt{\#main} (or a tag/commit) to pin a revision. For an editable install, \texttt{pkg> dev https://github.com/JeffreySarnoff/SmallFloats.jl} clones into \texttt{{\textasciitilde}/.julia/dev}, or \texttt{pkg> dev path/to/SmallFloats.jl} uses an existing clone.



Requires Julia 1.12. \texttt{Quadmath} (libquadmath{\textquotesingle}s \texttt{Float128}) is a hard dependency used only inside the oracle and exact-fallback paths; on platforms where it misbehaves, \texttt{ENV[{\textquotedbl}SmallFloats\_Float128{\textquotedbl}] = {\textquotedbl}disable{\textquotedbl}} before loading selects the pure-MPFR configuration with \textbf{identical results} (this equivalence is itself part of the test suite).



\section{Status and conformance}



\label{11983925649891201722}{}


\texttt{conformance()} returns the live, serializable conformance declaration (formats, operations, modes, instantiated table specializations, and κ registrations); \texttt{SmallFloats.draft\_revision()} names the implemented draft revision. Interpretations made where the draft text under-determined behavior are marked \texttt{[interp]} in the source comments.



\part{Cheat Sheet}


\chapter{Cheat Sheet}



\label{17671780375889550826}{}


A one-page reference for the SmallFloats.jl operations you are most likely to write. For explanations and semantics, see the \hyperlinkref{10272851679537904033}{User Guide}; for implementation details, see the \hyperlinkref{4463396622184713734}{Technical Guide}.



\section{Start here}



\label{11995297219376789591}{}



\begin{minted}{julia}
using SmallFloats
using Random: Xoshiro

T = Binary8p4se
ρ = RNE_SatNone

x = T(1.6)
y = T(0.25)
z = Add(T, ρ, x, y)
\end{minted}



The explicit operation shape is always:




\begin{minted}{julia}
Op(result_format, projection_spec, operands...; rng, R)
\end{minted}



For same-format operands under the default projection, ordinary Julia syntax is available:




\begin{minted}{julia}
z == x + y
exp(x)
fma(x, y, z)
\end{minted}



\section{Format names}



\label{4610047011167023553}{}



\begin{minted}{text}
Binary K p P (s|u) (e|f)
       │   │  │     └─ extended (Inf) or finite domain
       │   │  └────── signed or unsigned
       │   └───────── significand precision, including the implicit bit
       └───────────── total bitwidth, 3 through 8
\end{minted}



Examples:




\begin{table}[h]
\centering
\begin{tabulary}{\linewidth}{R R}
\toprule
Name & Meaning \\
\toprule
\texttt{Binary8p4se} & 8-bit, precision 4, signed, extended \\
\texttt{Binary8p4sf} & 8-bit, precision 4, signed, finite \\
\texttt{Binary6p3ue} & 6-bit, precision 3, unsigned, extended \\
\texttt{Binary5p5uf} & 5-bit, precision 5, unsigned, finite \\
\bottomrule
\end{tabulary}

\end{table}



The parametric spelling is equivalent:




\begin{minted}{julia}
Binary8p4se === Binary{8, 4, true, true}
\end{minted}



Useful format queries:




\begin{minted}{julia}
bitwidth(T)       # K
precision(T)      # P
issigned(T)
isextended(T)
expbias(T)
expbitwidth(T)
trailingsigbits(T)

MaxFiniteOf(T)
MinFiniteOf(T)
MinPositiveOf(T)
MinNormalOf(T)
MaxSubnormalOf(T)
\end{minted}



Every query above also takes a \textbf{value} instead of a type — \texttt{bitwidth(x)} ≡ \texttt{bitwidth(typeof(x))} — and folds to the same constant:




\begin{minted}{julia}
x = Binary8p4se(1.6)
bitwidth(x)       # 8
MaxFiniteOf(x)    # Binary8p4se(224.0 ≡ 0x7e)
\end{minted}



\section{Values and code points}



\label{1762024146591479678}{}



\begin{minted}{julia}
x = T(1.6)                    # numeric value: project into T
x = Convert(T, ρ, 1.6)       # explicit projection

c = codepoint(x)              # UInt8 code point
x = T(c)                      # UInt8 means validated CODE POINT
x = rawvalue(T, c)            # unchecked code-point construction

d = decode(x)                 # exact Float64 datum
\end{minted}



\begin{tcolorbox}[toptitle=-1mm,bottomtitle=1mm,colback=admonition-warning!50!white,colframe=admonition-warning,title=\textbf{`UInt8` means code point}]
\texttt{T(0x02)} constructs code point \texttt{0x02}; \texttt{T(2)} projects the numeric value two. Use \texttt{Convert} when intent should be unmistakable.

\end{tcolorbox}


Random values (uniform [0,1) floor-projected; normal round-to-nearest, tails clamped to MaxFinite; \texttt{randn} signed formats only):




\begin{minted}{julia}
rand(T); rand(T, dims); rand!(A)      # uniform on [0, 1)
randn(T); randn(T, dims); randn!(A)   # standard normal, always finite
rand(Xoshiro(1), T)                   # explicit rng as usual

rand(rng, T;  projection=RTP_SatNone)   # scalar forms: land under any ProjSpec
randn(rng, T; projection=RTZ_SatFinite) # (stochastic ρ draws R from same rng)
\end{minted}



Classification and stepping:




\begin{minted}{julia}
isnan(x); isinf(x); isfinite(x); iszero(x)
signbit(x); issubnormal(x)
Class(x)
NextGreaterThan(x); NextLessThan(x)
nextfloat(x); prevfloat(x)
TotalOrder(x, y)
\end{minted}



\section{Projection specifications}



\label{714366806163910976}{}


A projection specification is \texttt{(rounding mode, saturation mode)}:




\begin{minted}{julia}
ρ = ProjSpec(TowardPositive(), SatFinite())
roundingmode(ρ)
saturationmode(ρ)
\end{minted}



\subsection{Deterministic projections}



\label{18424161794924729217}{}



\begin{table}[h]
\centering
\begin{tabulary}{\linewidth}{R R R R}
\toprule
Rounding & \texttt{SatFinite} & \texttt{SatPropagate} & \texttt{SatNone} \\
\toprule
nearest, ties even & \texttt{RNE\_SatFinite} & \texttt{RNE\_SatPropagate} & \texttt{RNE\_SatNone} \\
nearest, ties away & \texttt{RNA\_SatFinite} & \texttt{RNA\_SatPropagate} & \texttt{RNA\_SatNone} \\
toward +∞ & \texttt{RTP\_SatFinite} & \texttt{RTP\_SatPropagate} & \texttt{RTP\_SatNone} \\
toward −∞ & \texttt{RTN\_SatFinite} & \texttt{RTN\_SatPropagate} & \texttt{RTN\_SatNone} \\
toward zero & \texttt{RTZ\_SatFinite} & \texttt{RTZ\_SatPropagate} & \texttt{RTZ\_SatNone} \\
round to odd & \texttt{RTO\_SatFinite} & \texttt{RTO\_SatPropagate} & \texttt{RTO\_SatNone} \\
\bottomrule
\end{tabulary}

\end{table}



\texttt{RNE\_SatNone} is the package-wide default. Saturation modes mean:




\begin{table}[h]
\centering
\begin{tabulary}{\linewidth}{R R}
\toprule
Mode & Out-of-range behavior \\
\toprule
\texttt{SatFinite} & clamp everything to the finite range \\
\texttt{SatPropagate} & preserve representable infinities; clamp other overflow \\
\texttt{SatNone} & apply the draft{\textquotesingle}s domain-, signedness-, and direction-dependent rows \\
\bottomrule
\end{tabulary}

\end{table}



\subsection{Stochastic projections}



\label{13387950688615589261}{}


Constructors are grouped by stochastic variant:




\begin{table}[h]
\centering
\begin{tabulary}{\linewidth}{R R R R}
\toprule
Variant & \texttt{SatFinite} & \texttt{SatPropagate} & \texttt{SatNone} \\
\toprule
A & \texttt{RSA\_SatFinite(N)} & \texttt{RSA\_SatPropagate(N)} & \texttt{RSA\_SatNone(N)} \\
B & \texttt{RSB\_SatFinite(N)} & \texttt{RSB\_SatPropagate(N)} & \texttt{RSB\_SatNone(N)} \\
C & \texttt{RSC\_SatFinite(N)} & \texttt{RSC\_SatPropagate(N)} & \texttt{RSC\_SatNone(N)} \\
\bottomrule
\end{tabulary}

\end{table}



\texttt{N} is the random-bit budget, \texttt{1 ≤ N ≤ 60}. Omitting it uses \texttt{N = 8}.




\begin{minted}{julia}
σ = RSA_SatNone(8)

Add(T, σ, x, y; rng=Xoshiro(1))  # reproducible stream
Add(T, σ, x, y; R=17)            # exact draw, ideal for tests

isstochastic(σ)                   # true
nrandbits(σ)                      # 8
\end{minted}



For an \texttt{N}-bit mode, explicit \texttt{R} must be in \texttt{0:(2{\textasciicircum}N - 1)}.



\subsection{Session defaults}



\label{3970345369238056532}{}


Read with \texttt{DefaultX()}, set with \texttt{DefaultX!(v)}:




\begin{minted}{julia}
DefaultType()            # Binary8p2se     DefaultType!(Binary8p4se)
DefaultReturnType()      # Binary8p2se     DefaultReturnType!(Binary8p3se)
DefaultAccumulatorType() # binary32        DefaultAccumulatorType!(binary64)
DefaultRoundingMode()    # NearestTiesToEven()
DefaultSaturationMode()  # SatNone()
DefaultProjection()      # RNE_SatNone
DefaultRNG()             # Xoshiro         DefaultRNG!(Xoshiro(42))
DefaultRbits()           # 8               DefaultRbits!(16)
\end{minted}



Setting a rounding/saturation component rebuilds \texttt{DefaultProjection}; setting \texttt{DefaultProjection!} directly updates both components. Always: \texttt{DefaultProjection() === ProjSpec(DefaultRoundingMode(), DefaultSaturationMode())}.



The convenience methods (\texttt{a + b}, \texttt{Exp(x)}, \texttt{T(2.1)}) \textbf{do} consult \texttt{DefaultProjection}, so changing it changes their results globally. They read it through a speculation guard: allocation-free while the default holds its initial value, one dispatch per call after you change it. Explicit forms (\texttt{Add(T, ρ, x, y)}) are unaffected by the session default.



Consume a default via the combinators — never by computing on a bare \texttt{DefaultX()} read. No dispatch while the default is unchanged, one barrier dispatch after a change; zero-alloc when \texttt{f}{\textquotesingle}s result type doesn{\textquotesingle}t depend on the default (a default-typed result boxes once at escape):




\begin{minted}{julia}
with_default_type((T, x) -> T(x), 1.5)          # Binary8p2se(1.5 ≡ 0x41)
with_default_projection((ρ, x, y) -> Add(T, ρ, x, y), x, y)
with_default_returntype(f, args...)              # f(DefaultReturnType(), args...)
with_default_accumulatortype(f, args...)
\end{minted}



\section{Scalar operation catalog}



\label{9067118353150552449}{}



\begin{minted}{julia}
# explicit result format and projection
Add(T, ρ, x, y)
FMA(T, ρ, x, y, z)
Convert(T, ρ, external_value)

# same-format default-projection convenience
Add(x, y)
FMA(x, y, z)
\end{minted}




\begin{table}[h]
\centering
\begin{tabulary}{\linewidth}{R R}
\toprule
Arity & Operations \\
\toprule
Unary & \texttt{Abs}, \texttt{Negate}, \texttt{Sqrt}, \texttt{RSqrt}, \texttt{Recip}, \texttt{Exp}, \texttt{Exp2}, \texttt{ExpMinusOne}, \texttt{Log}, \texttt{Log2}, \texttt{LogOnePlus}, \texttt{Softplus}, \texttt{Sin}, \texttt{Cos}, \texttt{Tan}, \texttt{ArcSin}, \texttt{ArcCos}, \texttt{ArcTan}, \texttt{Sinh}, \texttt{Cosh}, \texttt{Tanh}, \texttt{ArcSinh}, \texttt{ArcCosh}, \texttt{ArcTanh}, \texttt{SinPi}, \texttt{CosPi}, \texttt{TanPi}, \texttt{ArcSinPi}, \texttt{ArcCosPi}, \texttt{ArcTanPi} \\
Binary & \texttt{CopySign}, \texttt{Add}, \texttt{Subtract}, \texttt{Multiply}, \texttt{Divide}, \texttt{Hypot}, \texttt{ArcTan2}, \texttt{ArcTan2Pi}, \texttt{Maximum}, \texttt{Minimum}, \texttt{MaximumNumber}, \texttt{MinimumNumber}, \texttt{MaximumMagnitude}, \texttt{MinimumMagnitude}, \texttt{MaximumMagnitudeNumber}, \texttt{MinimumMagnitudeNumber}, \texttt{MinimumFinite}, \texttt{MaximumFinite} \\
Ternary & \texttt{FMA}, \texttt{FAA}, \texttt{Clamp} \\
Conversion & \texttt{Convert} \\
\bottomrule
\end{tabulary}

\end{table}



Common Base spellings under \texttt{RNE\_SatNone}:




\begin{minted}{julia}
x + y; x - y; x * y; x / y
-x; abs(x); inv(x); sqrt(x)
exp(x); exp2(x); expm1(x)
log(x); log2(x); log1p(x)
sin(x); cos(x); tan(x)
asin(x); acos(x); atan(x); atan(y, x)
sinh(x); cosh(x); tanh(x)
min(x, y); max(x, y); clamp(x, y, z)
fma(x, y, z); muladd(x, y, z)
\end{minted}



\section{Arrays}



\label{13119822296912292823}{}


Every registered operation has elementwise array methods:




\begin{minted}{julia}
A = T.([1.0, 1.5, 2.0])
B = T.([0.25, 0.5, 0.75])

C = Add(T, ρ, A, B)
E = Exp(T, ρ, A)
F = FMA(T, ρ, A, B, C)

C = vmap(:Add, T, ρ, A, B)
vmap!(C, Val(:Add), T, ρ, A, B)
\end{minted}



Operands and destination must have matching axes. Deterministic unary/binary operations use cached result tables; affordable ternary signatures may also use tables. Stochastic operations compute each element and consume one draw per projection.




\begin{minted}{julia}
table_bytes()
empty_tables!()
\end{minted}



\section{Blocks and scaled operations}



\label{17754733778172498420}{}



\begin{minted}{julia}
FS = Binary8p1uf
FE = Binary8p4se

b = Block(FS(4.0), FE(1.5), FE(-0.75), FE(2.0), FE(0.5))

blocksize(b)
scaleformat(b)
elemformat(b)

ConvertFromBlock(T, ρ, b)
BlockReduceAdd(T, ρ, b)
BlockReduceMultiply(T, ρ, b)
BlockDotProduct(T, ρ, b, b)
\end{minted}



Every scalar operation except \texttt{Convert} has generated block and scaled forms:




\begin{minted}{julia}
BlockAdd(T, ρ, b1, b2, result_scale)
BlockExp(T, ρ, b, result_scale)

ScaledAdd(T, ρ, scale1, x1, scale2, x2)
ScaledExp(T, ρ, scale, x)
\end{minted}



Quantize against a supplied or automatically selected scale:




\begin{minted}{julia}
ConvertToBlock(FS, FE, ρ, values_tuple, scale)
ConvertToBlockMaxAbsFinite(FS, FE, scale_ρ, element_ρ, values_tuple)
\end{minted}



\texttt{BlockVector} stores many equal-shape blocks in a structure-of-arrays layout.



\section{Packed storage}



\label{8567984038512272605}{}



\begin{minted}{julia}
v = Binary5p2se.([0.5, 1.0, 1.5, 2.0])
pv = PackedVector(v)

pv[2]
pv[2] = Binary5p2se(0.75)
collect(pv)

out = vmap(:Exp, Binary5p2se, ρ, pv)
\end{minted}



\texttt{PackedVector} stores each code point in exactly \texttt{bitwidth(F)} bits. Computation unpacks tiles internally; packed arithmetic is deliberately not in-place.



\section{Conformance and approximations}



\label{9146403741965371237}{}



\begin{minted}{julia}
conformance()
conformance_dict()
conformance_report()
draft_revision()

κ, exhaustive = measure_kappa(fn, :Exp, T, (T,), ρ)
register_approx!(:fast_exp, :Exp, T, (T,), ρ, fn; κ)
impl = approx(:fast_exp)
kappa(impl)
kappa_measured(impl)
list_approx()
unregister_approx!(:fast_exp)
\end{minted}



Approximate implementations are never substituted into the default API.



\section{Common traps}



\label{455411289700343066}{}



\begin{table}[h]
\centering
\begin{tabulary}{\linewidth}{R R}
\toprule
Trap & Correct pattern \\
\toprule
Treating \texttt{UInt8} as a number & \texttt{T(Int(c))} for a numeric integer; \texttt{T(c::UInt8)} for a code point \\
Silently mixing byte-float formats & Convert explicitly: \texttt{Convert(T, ρ, x)} \\
Assuming \texttt{SatNone} always clamps & Choose \texttt{SatFinite} when clamping is required \\
Expecting IEEE division-by-zero & P3109 semantics here define \texttt{x / 0 → NaN} \\
Expecting negative zero & Every format has one zero; \texttt{T(-0.0) === zero(T)} \\
Passing a \texttt{Rational} & Convert explicitly to an exact supported carrier or knowingly to \texttt{Float64} \\
Reproducibility with stochastic rounding & Supply a seeded \texttt{rng}, or an explicit \texttt{R} in tests \\
Using \texttt{rawvalue} on unchecked input & Prefer validated \texttt{T(c::UInt8)} outside kernels \\
\bottomrule
\end{tabulary}

\end{table}



\section{Performance checklist}



\label{15543690777485812783}{}


\begin{itemize}
\item Keep format types and projection specs in \texttt{const} bindings, function arguments, or type parameters.


\item Put code using runtime-selected formats behind a function barrier.


\item Expect the first deterministic array call for a specialization to build a table; benchmark warm calls separately.


\item Use \texttt{PackedVector} when storage bandwidth matters more than direct byte access.


\item Use \texttt{table\_bytes()} to inspect cache footprint and \texttt{empty\_tables!()} to reset it.


\item Do not replace explicit projection semantics with \texttt{@fastmath}.

\end{itemize}


For worked applications, continue to \hyperlinkref{7866567007412100191}{User Examples}. For correctness and benchmark methodology, continue to \hyperlinkref{4476403263922386643}{Technical Examples}.



\part{User Guide}


\chapter{User Guide}



\label{10272851679537904033}{}


Everything you need to use SmallFloats.jl, in the order you{\textquotesingle}ll need it. All transcripts in this guide are captured from real sessions against the shipped package.



\section{Formats and format queries}



\label{15797967065273794352}{}


The parametric type \texttt{Binary\{K,P,SGN,EXT\}} defines a format: bitwidth \texttt{K ∈ 3:8}, precision \texttt{P} (significand bits including the implicit bit), \texttt{SGN::Bool} for signed/unsigned, \texttt{EXT::Bool} for extended (with ±Inf) versus finite. Every legal combination — 120 formats — has its draft name exported:




\begin{minted}{jlcon}
julia> Binary8p4se === Binary{8,4,true,true}
true

julia> formatname(Binary{6,3,true,false})
:Binary6p3sf
\end{minted}



The suffix reads: \texttt{p4} precision 4, \texttt{s}/\texttt{u} signed/unsigned, \texttt{e}/\texttt{f} extended/finite.



Format introspection follows the draft{\textquotesingle}s Group M, with both Julia-style and draft-named accessors:




\begin{minted}{jlcon}
julia> bitwidth(Binary8p4se), precision(Binary8p4se), expbias(Binary8p4se)
(8, 4, 8)

julia> MaxFiniteOf(Binary8p4se)
Binary8p4se(224.0 ≡ 0x7e)

julia> MinPositiveOf(Binary8p4se)
Binary8p4se(0.0009765625 ≡ 0x01)
\end{minted}



Also available: \texttt{MinFiniteOf}, \texttt{MinNormalOf}, \texttt{MaxSubnormalOf}, \texttt{expbitwidth}, \texttt{trailingsigbits}, \texttt{issigned}, \texttt{isextended}, and the draft-named forms (\texttt{BitwidthOf}, \texttt{PrecisionOf}, \texttt{ExponentBiasOf}, …).



Each of these accepts a \textbf{value} as well as a type — the answer is a pure function of the format{\textquotesingle}s type parameters, so \texttt{bitwidth(x)} and \texttt{bitwidth(typeof(x))} are the same query and both fold to a literal:




\begin{minted}{jlcon}
julia> x = Binary8p4se(1.6);

julia> BitwidthOf(x), PrecisionOf(x), SignednessOf(x), DomainOf(x)
(8, 4, true, true)

julia> MaxFiniteOf(x)
Binary8p4se(224.0 ≡ 0x7e)
\end{minted}



\section{Values}



\label{18067454253769209299}{}


A value is one byte: an immutable wrapper around its code point. Four ways in, two ways out:




\begin{minted}{jlcon}
julia> x = Binary8p4se(1.6)              # construct from a Real: projects (rounds)
Binary8p4se(1.625 ≡ 0x45)

julia> Binary8p4se(0x45)                 # construct from a UInt8 CODE POINT (validated)
Binary8p4se(1.625 ≡ 0x45)

julia> rawvalue(Binary8p4se, 0x45)       # same, unchecked — the kernel-internal route
Binary8p4se(1.625 ≡ 0x45)

julia> Convert(Binary8p4se, RNE_SatNone, 3)   # explicit conversion, any mode
Binary8p4se(3.0 ≡ 0x4c)

julia> decode(x)                          # exact datum as Float64 (always exact)
1.625

julia> codepoint(x)                       # the raw byte (extends Base.codepoint)
0x45
\end{minted}



\begin{tcolorbox}[toptitle=-1mm,bottomtitle=1mm,colback=admonition-note!50!white,colframe=admonition-note,title=\textbf{UInt8 means code point; every other number means value}]
\texttt{UInt8} is the \emph{only} argument type with code-point semantics — mirroring \texttt{Char(0x41)}. Every other \texttt{Real} (including other \texttt{Integer}s) constructs by projecting the numeric value, and \texttt{Convert} is numeric for \textbf{all} integers:


\begin{minted}{jlcon}
julia> Binary8p4se(0x02), Binary8p4se(2)
(Binary8p4se(0.001953125 ≡ 0x02), Binary8p4se(2.0 ≡ 0x48))
\end{minted}

Out-of-range codes throw for K < 8 (\texttt{Binary5p3sf(0x20)} is an error); the range check costs nothing measurable — 2.1 ns, identical to unchecked \texttt{rawvalue}. Round-tripping is \texttt{T(codepoint(x)) === x}.

\end{tcolorbox}


\texttt{Convert} accepts \texttt{Binary} values (any format), \texttt{Float16/32/64}, \texttt{Float128}, \texttt{Integer}, and \texttt{BigFloat} — types whose values it can project \emph{exactly}. \texttt{Float128} projects directly, preserving all 113 significand bits: a value that sits a hair above a rounding midpoint converts correctly where staging through \texttt{Float64} would round it onto the midpoint first and then break the tie the wrong way. For other \texttt{Real}s, convert explicitly first (e.g. \texttt{Binary8p4se(Float64(π))}) and own the double rounding; \texttt{Rational} inputs throw with that guidance rather than double-round silently.



Every format has exactly one NaN and no negative zero; extended formats add ±Inf:




\begin{minted}{jlcon}
julia> Binary8p4se(1e9)                   # overflow under the default spec → +Inf
Binary8p4se(Inf ≡ 0x7f)

julia> Binary8p4se(-0.0)                  # no −0: projects to the single zero
Binary8p4se(0.0 ≡ 0x00)

julia> isnan(Binary8p4se(NaN)), isfinite(Binary8p4se(2.0))
(true, true)
\end{minted}



Standard predicates (\texttt{isnan}, \texttt{isinf}, \texttt{isfinite}, \texttt{iszero}, \texttt{signbit}, \texttt{issubnormal}), \texttt{zero}/\texttt{one}/\texttt{eps}/\texttt{typemin}/\texttt{typemax}, and stepping via \texttt{NextGreaterThan} / \texttt{NextLessThan} (the draft{\textquotesingle}s Next operations; \texttt{nextfloat}/\texttt{prevfloat} alias them) all work:




\begin{minted}{jlcon}
julia> NextGreaterThan(Binary8p4se(1.0))
Binary8p4se(1.125 ≡ 0x41)

julia> Class(Binary8p4se(0.01))           # draft classification
ClassPosNormal::FPClass = 0x06
\end{minted}



\section{Random values}



\label{13027969696381432085}{}


The Random API works on every format, in every setup Julia programmers expect.



\textbf{What the draws mean.} \texttt{rand} draws the real uniform on [0, 1) at Float64 and floor-projects it onto the format{\textquotesingle}s grid, so each code point receives exactly the real measure of its interval; results are always in [0, 1). \texttt{randn} projects a standard-normal Float64 draw round-to-nearest with \texttt{SatFinite}, so tail draws beyond \texttt{MaxFiniteOf(T)} clamp to the extremal finite datum — \texttt{randn} never returns ±Inf or NaN, which matters for tiny-range formats (\texttt{Binary3p1se} has \texttt{MaxFinite = 1.0}). \texttt{randn} requires a signed format; an unsigned format throws:




\begin{minted}{jlcon}
julia> randn(Binary6p3ue)
ERROR: ArgumentError: randn requires a signed format; Binary6p3ue cannot represent negative draws
\end{minted}



\textbf{Setup 1 — quick and implicit.} No rng argument: draws come from Julia{\textquotesingle}s task-local default generator (a \texttt{Xoshiro}), and \texttt{Random.seed!} controls them exactly as for any other type:




\begin{minted}{jlcon}
julia> Random.seed!(1234); rand(Binary8p4se)
Binary8p4se(0.3125 ≡ 0x32)

julia> Random.seed!(1234); rand(Binary8p4se)      # same seed, same draw
Binary8p4se(0.3125 ≡ 0x32)
\end{minted}



\textbf{Setup 2 — an explicit rng stream.} Pass any \texttt{AbstractRNG} first, for reproducible streams independent of global state:




\begin{minted}{jlcon}
julia> rand(Xoshiro(1), Binary8p4se)
Binary8p4se(0.0703125 ≡ 0x21)

julia> rand(Xoshiro(8), Binary8p4se, 4)
4-element Vector{Binary8p4se}:
   Binary8p4se(0.40625 ≡ 0x35)
 Binary8p4se(0.021484375 ≡ 0x13)
      Binary8p4se(0.75 ≡ 0x3c)
   Binary8p4se(0.28125 ≡ 0x31)
\end{minted}



\textbf{Setup 3 — arrays and in-place.} \texttt{rand(T, dims...)} / \texttt{randn(T, dims...)} build \texttt{Array\{T\}}; \texttt{rand!(A)} / \texttt{randn!(A)} fill an existing array (one byte per element — cheap to preallocate and reuse):




\begin{minted}{jlcon}
julia> A = Vector{Binary8p4se}(undef, 3); randn!(Xoshiro(2), A)
3-element Vector{Binary8p4se}:
 Binary8p4se(-0.005859375 ≡ 0x86)
        Binary8p4se(1.75 ≡ 0x46)
        Binary8p4se(-1.0 ≡ 0xc0)

julia> randn(Xoshiro(3), Binary8p4se, 2, 3) |> typeof
Matrix{Binary8p4se}
\end{minted}



Feed array draws straight into the storage layers when that is the goal: \texttt{PackedVector(rand(Binary4p2se, n))}.



\textbf{Setup 4 — choosing the projection.} The scalar \texttt{::Type} forms take a \texttt{projection} keyword to land the draw under any \texttt{ProjSpec}:




\begin{minted}{jlcon}
julia> rand(Xoshiro(2), Binary8p4se; projection = RTP_SatNone)    # ceiling
Binary8p4se(0.0029296875 ≡ 0x03)

julia> randn(Xoshiro(6), Binary8p4se; projection = RTZ_SatFinite) # toward zero
Binary8p4se(-1.875 ≡ 0xc7)

julia> rand(Xoshiro(4), Binary8p4se; projection = RSA_SatNone(8)) # stochastic
Binary8p4se(0.8125 ≡ 0x3d)
\end{minted}



A stochastic projection draws its random bits from the \emph{same} rng, so seeded streams stay reproducible. The defaults are the contract-keepers (floor for \texttt{rand}, nearest + \texttt{SatFinite} for \texttt{randn}); opting out can produce \texttt{1.0} from \texttt{rand} or ±Inf/NaN from \texttt{randn}. The array and \texttt{!} forms always use the defaults — for arrays under another projection, draw scalars: \texttt{[rand(rng, T; projection = ρ) for \_ in 1:n]}. Worked transcripts, including the K = 3 tail-clamp comparison: \hyperlinkref{7866567007412100191}{User Examples}.



\section{Projection specifications}



\label{714366806163910976}{}


Every rounding decision in the package is governed by a \texttt{ProjSpec}, the pair of a \textbf{rounding mode} and a \textbf{saturation mode}. Both are zero-size singleton types, so a \texttt{ProjSpec} costs nothing and fully specializes every call it reaches.



Rounding modes: \texttt{NearestTiesToEven}, \texttt{NearestTiesToAway}, \texttt{TowardPositive}, \texttt{TowardNegative}, \texttt{TowardZero}, \texttt{ToOdd}, and the stochastic families \texttt{StochasticA\{N\}}, \texttt{StochasticB\{N\}}, \texttt{StochasticC\{N\}} with a random-bit budget \texttt{1 ≤ N ≤ 60}. Saturation modes: \texttt{SatFinite}, \texttt{SatPropagate}, \texttt{SatNone}.




\begin{minted}{jlcon}
julia> ρ = ProjSpec(TowardPositive(), SatFinite())
(TowardPositive, SatFinite)

julia> RTP_SatFinite === ρ                # every pairing is also predefined
true
\end{minted}



\subsection{Predefined specs}



\label{13218842448245266752}{}


Every deterministic (rounding, saturation) pairing is exported as a constant, named \texttt{R<mode>\_Sat<mode>}:




\begin{table}[h]
\centering
\begin{tabulary}{\linewidth}{R R R R}
\toprule
 & \texttt{SatFinite} & \texttt{SatPropagate} & \texttt{SatNone} \\
\toprule
\texttt{NearestTiesToEven} & \texttt{RNE\_SatFinite} & \texttt{RNE\_SatPropagate} & \texttt{RNE\_SatNone} \\
\texttt{NearestTiesToAway} & \texttt{RNA\_SatFinite} & \texttt{RNA\_SatPropagate} & \texttt{RNA\_SatNone} \\
\texttt{TowardPositive} & \texttt{RTP\_SatFinite} & \texttt{RTP\_SatPropagate} & \texttt{RTP\_SatNone} \\
\texttt{TowardNegative} & \texttt{RTN\_SatFinite} & \texttt{RTN\_SatPropagate} & \texttt{RTN\_SatNone} \\
\texttt{TowardZero} & \texttt{RTZ\_SatFinite} & \texttt{RTZ\_SatPropagate} & \texttt{RTZ\_SatNone} \\
\texttt{ToOdd} & \texttt{RTO\_SatFinite} & \texttt{RTO\_SatPropagate} & \texttt{RTO\_SatNone} \\
\bottomrule
\end{tabulary}

\end{table}



The stochastic families are parameterized by the random-bit budget \texttt{N}, so their predefined forms are \emph{constructors} rather than constants: \texttt{RSA\_*} for \texttt{StochasticA}, \texttt{RSB\_*} for \texttt{StochasticB}, \texttt{RSC\_*} for \texttt{StochasticC}, each crossed with the three saturation modes. Call them with the budget — or without, for the default \texttt{N = 8} (\texttt{SmallFloats.DEFAULT\_RBITS}):




\begin{minted}{jlcon}
julia> RSA_SatNone()                      # StochasticA with the default N = 8
(StochasticA[8], SatNone)

julia> RSC_SatFinite(16) === ProjSpec(StochasticC{16}(), SatFinite())
true
\end{minted}



\texttt{RNE\_SatNone} is the package{\textquotesingle}s \emph{initial} default spec. \texttt{default\_projspec} reads the session default (see \textbf{Session defaults} below), so every Base-register operator, the same-format convenience methods, and \texttt{T(x::Real)} construction follow \texttt{DefaultProjection()} — \texttt{RNE\_SatNone} until you change it.



Rounding chooses the neighbor; saturation decides what out-of-range results become. Watch all three interact on an overflowing product in \texttt{Binary8p4se} (\texttt{MaxFinite = 224}):




\begin{minted}{jlcon}
julia> w, two = Binary8p4se(200.0), Binary8p4se(2.0);

julia> Multiply(Binary8p4se, RNE_SatNone, w, two)       # overflow → ±Inf
Binary8p4se(Inf ≡ 0x7f)

julia> Multiply(Binary8p4se, RNE_SatFinite, w, two)     # overflow → MaxFinite
Binary8p4se(224.0 ≡ 0x7e)

julia> Multiply(Binary8p4se, RTZ_SatNone, w, two)
Binary8p4se(224.0 ≡ 0x7e)   # SatNone + directed-toward-zero clamps finite
\end{minted}



\begin{tcolorbox}[toptitle=-1mm,bottomtitle=1mm,colback=admonition-note!50!white,colframe=admonition-note,title=\textbf{Note}]
\texttt{SatNone} is not {\textquotedbl}no saturation handling{\textquotedbl}: it is the draft{\textquotesingle}s row set in which nearest modes overflow to ±Inf (or NaN for finite formats) while directed modes that point away from the overflow clamp to the extremal finite value.

\end{tcolorbox}


\subsection{Stochastic rounding}



\label{5064376698567748197}{}


Stochastic modes consume \texttt{N} random bits per projection. You control the source three ways — implicitly (task-local RNG), by passing an \texttt{rng}, or by fixing the draw \texttt{R} itself, which makes single projections exactly reproducible and is the right tool for tests:




\begin{minted}{jlcon}
julia> σ = RSA_SatNone();                 # ≡ ProjSpec(StochasticA{8}(), SatNone())

julia> Add(Binary8p4se, σ, Binary8p4se(2.0), Binary8p4se(0.03125); rng = Xoshiro(1))
Binary8p4se(2.0 ≡ 0x48)

julia> Add(Binary8p4se, σ, Binary8p4se(2.0), Binary8p4se(0.03125); R = 0)
Binary8p4se(2.0 ≡ 0x48)      # smallest draw: rounds down

julia> Add(Binary8p4se, σ, Binary8p4se(2.0), Binary8p4se(0.03125); R = 255)
Binary8p4se(2.25 ≡ 0x49)     # largest draw: rounds up
\end{minted}



The exact fraction here is 1/8 of an ulp, so over all 256 draws exactly 32 round up — \texttt{StochasticA} is unbiased in expectation. \texttt{R} must lie in \texttt{0:2{\textasciicircum}N-1}.



\subsection{Session defaults}



\label{3970345369238056532}{}


Six session-wide defaults are readable as \texttt{DefaultX()} and settable as \texttt{DefaultX!(v)}:




\begin{table}[h]
\centering
\begin{tabulary}{\linewidth}{R R R}
\toprule
default & initial value & setter accepts \\
\toprule
\texttt{DefaultType} & \texttt{Binary8p2se} & any fully-parameterized \texttt{Binary} type \\
\texttt{DefaultReturnType} & \texttt{Binary8p2se} & any fully-parameterized \texttt{Binary} type \\
\texttt{DefaultAccumulatorType} & \texttt{binary32} (\texttt{Float32}) & any \texttt{AbstractFloat} type \\
\texttt{DefaultRoundingMode} & \texttt{NearestTiesToEven()} & mode instance or type \\
\texttt{DefaultSaturationMode} & \texttt{SatNone()} & mode instance or type \\
\texttt{DefaultProjection} & \texttt{RNE\_SatNone} & a \texttt{ProjSpec}, or \texttt{(mode, sat)} \\
\texttt{DefaultRNG} & the \texttt{Xoshiro} type & RNG type or instance \\
\texttt{DefaultRbits} & \texttt{8} & \texttt{Int} in \texttt{1:60} \\
\bottomrule
\end{tabulary}

\end{table}



The projection default and its components are kept coherent in both directions: setting \texttt{DefaultRoundingMode!} or \texttt{DefaultSaturationMode!} rebuilds \texttt{DefaultProjection} around the change, and setting \texttt{DefaultProjection!} directly decomposes it back into both components — so \texttt{DefaultProjection() === ProjSpec(DefaultRoundingMode(), DefaultSaturationMode())} always holds.




\begin{minted}{jlcon}
julia> DefaultRoundingMode!(TowardZero())
TowardZero()

julia> DefaultProjection()               # followed the component
(TowardZero, SatNone)

julia> DefaultProjection!(RNA_SatFinite)
(NearestTiesToAway, SatFinite)

julia> DefaultRoundingMode(), DefaultSaturationMode()   # followed the projection
(NearestTiesToAway(), SatFinite())
\end{minted}



\texttt{DefaultProjection} is not merely advisory: \texttt{default\_projspec} reads it, so the same-format convenience methods (\texttt{a + b}, \texttt{Exp(x)}, …), the Base-register operators, and \texttt{T(x::Real)} construction all follow it.




\begin{minted}{jlcon}
julia> x, y = Binary8p4se(200.0), Binary8p4se(2.0);

julia> x * y                              # RNE_SatNone: overflow → +Inf
Binary8p4se(Inf ≡ 0x7f)

julia> DefaultProjection!(RTZ_SatFinite);

julia> x * y                              # now clamps to MaxFinite
Binary8p4se(224.0 ≡ 0x7e)

julia> Multiply(Binary8p4se, RNE_SatNone, x, y)   # explicit ρ is unaffected
Binary8p4se(Inf ≡ 0x7f)
\end{minted}



\begin{tcolorbox}[toptitle=-1mm,bottomtitle=1mm,colback=admonition-warning!50!white,colframe=admonition-warning,title=\textbf{Changing the default is a global semantic change}]
Every caller of the convenience forms — including code in other packages — sees it. Library code that needs a specific projection should name it explicitly rather than rely on the session default.

\end{tcolorbox}


Following the default costs nothing while it holds its initial value: the convenience forms consume it through the same speculation guard as the \texttt{with\_default\_*} combinators, so they compile against the constant and stay allocation-free with concretely inferred results (pinned in the test suite). Once you change the default they cross a function barrier instead — one dynamic dispatch per call, everything inside still specialized. The explicit forms \texttt{Add(T, ρ, x, y)} are unaffected either way.



To \emph{consume} a default in your own code without paying dynamic-dispatch costs, go through the \texttt{with\_default\_*} combinators — \texttt{with\_default\_type}, \texttt{with\_default\_returntype}, \texttt{with\_default\_accumulatortype}, \texttt{with\_default\_projection} — which call \texttt{f(default, args...)}:




\begin{minted}{jlcon}
julia> with_default_type((T, x) -> T(x), 1.5)
Binary8p2se(1.5 ≡ 0x41)
\end{minted}



While a default still holds its initial value, the combinator{\textquotesingle}s call is statically compiled against that constant — no dynamic dispatch. After the default is changed it crosses a function barrier: one dynamic dispatch at entry, everything inside fully specialized.



Allocation contract: when \texttt{f}{\textquotesingle}s result type does not depend on the default — \texttt{with\_default\_projection} with the formats fixed by the caller is the normal shape — the call is zero-allocation with a concretely inferred result (pinned in the test suite). When the result{\textquotesingle}s type \emph{is} the default (\texttt{with\_default\_type} used as a constructor), the value is computed on the specialized path but boxes once where it escapes — the irreducible cost of a runtime-chosen type.



\section{Scalar operations: the two registers}



\label{2538548316682979338}{}


\textbf{The spec-named register} exposes every draft operation under its draft name with an explicit result format and spec: \texttt{Op(fr, ρ, operands...)}. Operands may be any formats; the result format is the first argument. The catalog:



\begin{itemize}
\item \textbf{30 unary} — \texttt{Abs}, \texttt{Negate}, \texttt{Sqrt}, \texttt{RSqrt}, \texttt{Recip}, \texttt{Exp}, \texttt{Exp2}, \texttt{ExpMinusOne}, \texttt{Log}, \texttt{Log2}, \texttt{LogOnePlus}, \texttt{Softplus}, the trig/hyperbolic families, and the π-scaled families;


\item \textbf{18 binary} — \texttt{Add}, \texttt{Subtract}, \texttt{Multiply}, \texttt{Divide}, \texttt{CopySign}, \texttt{Hypot}, \texttt{ArcTan2}, \texttt{ArcTan2Pi}, and the eight Minimum/Maximum variants plus the Finite variants;


\item \textbf{3 ternary} — \texttt{FMA}, \texttt{FAA}, \texttt{Clamp};


\item \textbf{\texttt{Convert}} — the one operation that also accepts non-\texttt{Binary} operands.

\end{itemize}


\textbf{The Base register} makes each format an ordinary Julia number under its \emph{default} spec (\texttt{RNE\_SatNone}): \texttt{+ - * /}, \texttt{fma}, \texttt{exp}, \texttt{log}, \texttt{sqrt}, \texttt{min}, \texttt{max}, \texttt{abs}, \texttt{atan(y, x)}, \texttt{sinpi}, \texttt{inv}, comparisons, and friends — same-format operands only, no silent cross-format promotion (mixing formats promotes to \texttt{Float64} explicitly).




\begin{minted}{jlcon}
julia> exp(Binary8p4se(0.25))                       # Base register
Binary8p4se(1.25 ≡ 0x42)

julia> Exp(Binary8p4se, RNE_SatNone, Binary8p4se(0.25))   # identical, explicit
Binary8p4se(1.25 ≡ 0x42)
\end{minted}



Semantics highlights (all per the draft): \texttt{x/0 → NaN} for every \texttt{x} including ±Inf; \texttt{Recip(±Inf) = 0}; \texttt{0·∞ → NaN}; a NaN operand generally propagates, except the \texttt{*Number}/\texttt{*Finite} extremum variants which prefer non-NaN (and finite) operands; trig of ±Inf is NaN; the π-scaled family reduces exactly mod 2 first.



\section{Arrays, kernels, and sorting}



\label{12068187713603161220}{}


Every operation has array methods — \texttt{Op(fr, ρ, A)}, \texttt{Op(fr, ρ, A, B)}, \texttt{Op(fr, ρ, A, B, C)} — plus the generic \texttt{vmap} / \texttt{vmap!}:




\begin{minted}{jlcon}
julia> A = Binary8p4se.(randn(Xoshiro(7), 4) .* 2)
4-element Vector{Binary8p4se}:
 Binary8p4se(-0.875 ≡ 0xbe)
 Binary8p4se(4.0 ≡ 0x50)
 Binary8p4se(-3.25 ≡ 0xcd)
 Binary8p4se(2.25 ≡ 0x49)

julia> Exp(Binary8p4se, RNE_SatNone, A)
4-element Vector{Binary8p4se}:
 Binary8p4se(0.40625 ≡ 0x35)
 Binary8p4se(56.0 ≡ 0x6e)
 Binary8p4se(0.0390625 ≡ 0x1a)
 Binary8p4se(9.0 ≡ 0x59)
\end{minted}



For pure (non-stochastic) specs, unary and binary array calls run as \textbf{table gathers}: the first call builds and caches a 256-byte (unary) or 64 KiB (binary) result table for that exact \texttt{(op, formats, ρ)} specialization, and every later element costs a single lookup — measured at 0.27 ns/element unary, 0.5 ns/element binary.



Ternary operations (\texttt{FMA}, \texttt{FAA}, \texttt{Clamp}) ride the same gather whenever the three operand bitwidths keep the table affordable, tiered by \texttt{K1 + K2 + K3}:



\begin{itemize}
\item \textbf{Eager} (up to 256 KiB; every all-\texttt{K ≤ 6} signature): built on the first array call, like unary/binary.


\item \textbf{Adaptive} (up to 2 MiB; the \texttt{K = 7} band): built only once a signature has processed enough elements to earn its build, in a byte-bounded, LRU-evicted cache.


\item \textbf{Compute} (\texttt{K = 8}; a 16 MiB table stops being a cache win): the scalar pipeline runs per element, optionally threaded for long arrays.

\end{itemize}


Every table entry — eager or adaptive — is built through the scalar path, so it is bit-identical by construction. Stochastic calls always run the scalar pipeline per element (each element draws its own random bits).



Inspect or reset the cache — unary/binary and ternary alike — with \texttt{table\_bytes()} and \texttt{empty\_tables!()}. The ternary policy thresholds and the threading cutoff are internal \texttt{Ref}s (\texttt{SmallFloats.TERNARY\_EAGER\_BITS} and neighbors in \texttt{tables.jl} and \texttt{kernels.jl}) for the rare case you need to tune them.



Sorting is special-cased: values compare through integer order keys, and vectors of \texttt{Binary} sort with an \textbf{O(n) counting sort} installed as the default algorithm — about 8× the stock comparison sort at 64 K elements, \texttt{rev=true} included. \texttt{sort(A)} just works; NaN sorts last (first under \texttt{rev=true}), matching Base{\textquotesingle}s conventions. \texttt{TotalOrder(x, y)} exposes the draft{\textquotesingle}s total order directly (single NaN largest).



\section{Blocks and scaled operations}



\label{17754733778172498420}{}


A \texttt{Block\{B,FS,FE\}} is a scale in format \texttt{FS} and \texttt{B} elements in format \texttt{FE}; its represented values are \texttt{scale × element} (the MX-style scheme). Construct directly or quantize a tuple with the draft{\textquotesingle}s algorithm, which picks the scale from the maximum finite |element| and projects each element against it:




\begin{minted}{jlcon}
julia> b = Block(Binary8p1uf(4.0), Binary8p4se(1.5), Binary8p4se(-0.75),
                 Binary8p4se(2.0), Binary8p4se(0.5))
Block{4, Binary8p1uf, Binary8p4se}(Binary8p1uf(4.0 ≡ 0x82), (…))

julia> ConvertFromBlock(Binary8p3se, RNE_SatNone, b)   # decode scale×elem, project
(Binary8p3se(6.0 ≡ 0x4a), Binary8p3se(-3.0 ≡ 0xc6), Binary8p3se(8.0 ≡ 0x4c), Binary8p3se(2.0 ≡ 0x44))

julia> BlockReduceAdd(Binary8p4se, RNE_SatNone, b)     # exact Σ scale·xᵢ, then project
Binary8p4se(13.0 ≡ 0x5d)

julia> ConvertToBlockMaxAbsFinite(Binary8p1uf, Binary8p4se, RNE_SatNone, RNE_SatNone,
           (Binary8p4se(100.0), Binary8p4se(-12.0), Binary8p4se(0.5), Binary8p4se(3.0)))
Block{4, Binary8p1uf, Binary8p4se}(Binary8p1uf(64.0 ≡ 0x86), (…))
\end{minted}



Every scalar operation lifts to blocks (\texttt{BlockAdd}, \texttt{BlockExp}, …, taking blocks and a result scale) and to the scaled form \texttt{ScaledOp} (\texttt{ScaledAdd(fr, ρ, s1, x1, s2, x2)} — block size 1). Reductions: \texttt{BlockReduceAdd}, \texttt{BlockReduceMultiply}, and \texttt{BlockDotProduct(fr, ρ, bx, by)}, whose lane products and accumulation are \textbf{exact} before the single final projection — there is no hidden intermediate rounding. \texttt{ConvertToBlock(fs, fr, ρ, xs, s)} quantizes against a scale you supply. \texttt{BlockVector} stores many same-shape blocks in structure-of-arrays layout.



\section{Packed storage}



\label{8567984038512272605}{}


For memory-bound work, \texttt{PackedVector\{F\}} stores code points at \texttt{bitwidth(F)} bits per element (a \texttt{Binary5p2se} vector at 5 bits/element instead of 8):




\begin{minted}{jlcon}
julia> v = Binary5p2se.(rand(Xoshiro(3), 6) .* 4);

julia> pv = PackedVector(v); sizeof(pv.data)
8                     # 6 × 5 bits = 30 bits → one 64-bit word

julia> pv[3] == v[3]
true
\end{minted}



It is a full \texttt{AbstractVector\{F\}} (indexing, \texttt{collect}, iteration), and \texttt{vmap} accepts it directly, unpacking cache-friendly tiles internally. The rule is \emph{store packed, compute unpacked}: there is deliberately no in-place packed arithmetic.



\section{Conformance and κ-approximate implementations}



\label{16045636492794422690}{}


\texttt{conformance()} returns the live declaration — formats, operations, mode vocabulary, the table specializations actually instantiated this session, and every registered approximation. \texttt{conformance\_report()} prints it; \texttt{conformance\_dict()} returns a plain nested \texttt{Dict} for JSON/TOML serialization.



The default API is bit-exact, always. If you \emph{want} a faster inexact kernel, register it — and the registry measures your honesty:




\begin{minted}{jlcon}
julia> ftz = ftz_variant(:Exp, Binary8p4se, Binary8p4se, RNE_SatFinite);  # Annex example

julia> impl = register_approx!(:my_fast_exp, :Exp, Binary8p4se, (Binary8p4se,),
                               RNE_SatFinite, ftz);

julia> kappa(:my_fast_exp)      # measured max code-point deviation, verified exhaustively
4.0
\end{minted}



Declaring \texttt{κ} smaller than the measured deviation throws; NaN-mismatching implementations must be acknowledged with an explicit \texttt{κ = NaN}. Retrieve with \texttt{approx(:my\_fast\_exp)}, list with \texttt{list\_approx()}, and measure anything yourself with \texttt{measure\_kappa} / \texttt{codedistance}.



\section{Performance guidance}



\label{8216999033093502309}{}


\begin{itemize}
\item \textbf{Pass format types statically.} Through \texttt{const} bindings, type parameters, or function arguments, every entry point fully specializes (scalar \texttt{Add} ≈ 18–26 ns, \texttt{project} ≈ 13 ns, zero allocations). A format type read from a \textbf{non-\texttt{const} global} forces Julia{\textquotesingle}s dynamic dispatch on every call ({\textasciitilde}1 µs for keyword calls); one function barrier \texttt{f(::Type\{T\}, …) where \{T\}} restores full speed.


\item \textbf{The convenience forms are free at the initial default.} \texttt{x + y}, \texttt{Exp(x)}, and \texttt{T(2.1)} read the session default through a speculation guard, so they are allocation-free and concretely inferred while it holds its initial value. After you change the default they cost one dynamic dispatch per call; name ρ explicitly (\texttt{Add(T, ρ, x, y)} with a \texttt{const} ρ) in hot code that must be insensitive to the session default.


\item \textbf{Bulk work belongs in array calls.} The table-gather kernels are {\textasciitilde}50× the scalar path; the first call per specialization pays a one-time build (≈ 0.4 ms unary, tens of ms for 8×8 binary tables, up to a few ms for a 2 MiB ternary table).


\item \textbf{Ternary array calls scale with bitwidth.} \texttt{K ≤ 6} operand formats table eagerly ({\textasciitilde}35× the scalar loop); \texttt{K = 7} tables adaptively once a signature is hot; \texttt{K = 8} runs the compute kernel, optionally threaded for long arrays ({\textasciitilde}4× at 4 threads) — no action needed, the array call picks the right path.


\item \textbf{Stochastic array calls} run the scalar pipeline per element; pass an explicit \texttt{rng} for reproducibility.


\item \textbf{Memory:} \texttt{PackedVector} for storage; \texttt{BlockVector} for many blocks.


\item The reproducible benchmark suite lives at \texttt{benchmarking/benchmarking.jl} and generates a full markdown report for your machine.

\end{itemize}


\section{Environment switches}



\label{2477270776453872570}{}


\texttt{ENV[{\textquotedbl}SmallFloats\_Float128{\textquotedbl}] = {\textquotedbl}disable{\textquotedbl}} (set before \texttt{using SmallFloats}) disables the internal \texttt{Float128} fast paths in favor of pure MPFR — results are bit-identical (tested); only build/oracle speed changes.



\part{User Examples}


\chapter{User Examples}



\label{7866567007412100191}{}


Worked, runnable examples. Every output shown was captured from a real session; seeds are fixed so you can reproduce them exactly.



\section{Basic}



\label{13551325986707881962}{}


\subsection{One value, every rounding mode}



\label{14660386703135810651}{}


The value 2.30078125 sits between the \texttt{Binary8p4se} grid points 2.25 and 2.5 (ulp = 0.25 in [2, 4)):




\begin{minted}{julia}
using SmallFloats

for μ in (NearestTiesToEven(), NearestTiesToAway(), TowardPositive(),
          TowardNegative(), TowardZero(), ToOdd())
    v = Convert(Binary8p4se, ProjSpec(μ, SatNone()), 2.30078125)
    println(rpad(string(typeof(μ)), 20), " → ", decode(v))
end
\end{minted}




\begin{minted}{text}
NearestTiesToEven    → 2.25
NearestTiesToAway    → 2.25
TowardPositive       → 2.5
TowardNegative       → 2.25
TowardZero           → 2.25
ToOdd                → 2.25
\end{minted}



(\texttt{ToOdd} keeps 2.25 because its significand is already odd; try 2.05 to see it move.)



\subsection{Enumerating a whole format}



\label{15296297375297823522}{}


Small formats are small enough to look at in full — 16 code points for \texttt{Binary4p2se}, sorted by the total order (single NaN last):




\begin{minted}{julia}
decode.(sort(Binary4p2se.(0x00:0x0f)))      # broadcast the code-point constructor
\end{minted}




\begin{minted}{text}
[-Inf, -2.0, -1.5, -1.0, -0.75, -0.5, -0.25, 0.0,
  0.25, 0.5, 0.75, 1.0, 1.5, 2.0, Inf, NaN]
\end{minted}



This is a good way to build intuition for a format before committing to it: you can \emph{see} the subnormal spacing, the binade structure, and the dynamic range.



\subsection{Saturation in one line each}



\label{18285742174011812433}{}



\begin{minted}{jlcon}
julia> w, two = Binary8p4se(200.0), Binary8p4se(2.0);   # MaxFinite is 224

julia> Multiply(Binary8p4se, RNE_SatNone, w, two)
Binary8p4se(Inf ≡ 0x7f)

julia> Multiply(Binary8p4se, RNE_SatFinite, w, two)
Binary8p4se(224.0 ≡ 0x7e)
\end{minted}



\subsection{Random values under a chosen projection}



\label{10520373459060902538}{}


\texttt{rand} and \texttt{randn} take a \texttt{projection} keyword on their scalar \texttt{::Type} forms. One underlying Float64 draw, landed on the format{\textquotesingle}s grid four different ways — the seed fixes the draw, the projection decides the landing:




\begin{minted}{jlcon}
julia> rand(Xoshiro(2), Float64)                # the underlying uniform draw
0.0022505868897625403

julia> rand(Xoshiro(2), Binary8p4se)            # default: floor (stays in [0,1))
Binary8p4se(0.001953125 ≡ 0x02)

julia> rand(Xoshiro(2), Binary8p4se; projection = RTP_SatNone)   # ceiling
Binary8p4se(0.0029296875 ≡ 0x03)

julia> rand(Xoshiro(2), Binary8p4se; projection = RNE_SatNone)   # nearest
Binary8p4se(0.001953125 ≡ 0x02)

julia> rand(Xoshiro(2), Binary8p4se; projection = RSA_SatNone(8))  # stochastic
Binary8p4se(0.001953125 ≡ 0x02)
\end{minted}



A stochastic projection draws its random bits from the \emph{same} rng as the uniform draw, so seeded streams stay reproducible. Note the default is the contract-keeper: a nearest or upward projection can return exactly \texttt{1.0} (mass near the top rounds up), which is why floor is the default.



The same keyword on \texttt{randn} — here the draw is \texttt{-1.9757…}:




\begin{minted}{jlcon}
julia> randn(Xoshiro(6), Binary8p4se)           # default: nearest + SatFinite
Binary8p4se(-2.0 ≡ 0xc8)

julia> randn(Xoshiro(6), Binary8p4se; projection = RTZ_SatFinite)  # toward zero
Binary8p4se(-1.875 ≡ 0xc7)
\end{minted}



Where the saturation half of the default earns its keep: \texttt{Binary3p1se} has \texttt{MaxFinite = 1.0}, so normal tail draws overflow constantly. Seed 360 draws \texttt{z = 3.326…}:




\begin{minted}{jlcon}
julia> randn(Xoshiro(360), Binary3p1se)         # SatFinite clamps the tail
Binary3p1se(1.0 ≡ 0x02)

julia> randn(Xoshiro(360), Binary3p1se; projection = RNE_SatNone)
Binary3p1se(Inf ≡ 0x03)                         # opt out, get the overflow
\end{minted}



The array and \texttt{!} forms (\texttt{rand(T, dims)}, \texttt{randn!(A)}, …) always use the defaults; for arrays under another projection, draw scalars: \texttt{[rand(rng, T; projection = ρ) for \_ in 1:n]}.



\section{General AI}



\label{12995498770613901834}{}


\subsection{Log-odds belief updating — and where tiny formats bite reasoning}



\label{12572314235496160742}{}


Classic probabilistic reasoning: accumulate independent evidence as log-likelihood ratios. A sensor with hit rate 0.85 and false-alarm rate 0.30 contributes \texttt{log(0.85/0.30) ≈ 1.0415} per alarm. In \texttt{Binary8p3se} that datum is 1.0 — and the running belief stalls exactly the way the gradient accumulator does in the Deep Learning section:




\begin{minted}{julia}
using SmallFloats

function logodds_demo(nalarms)
    llr = Binary8p3se(log(0.85 / 0.30))      # quantized evidence weight
    acc = Binary8p3se(0.0)
    for _ in 1:nalarms
        acc = Add(Binary8p3se, RNE_SatNone, acc, llr)
    end
    post(l) = 1 / (1 + exp(-l))
    exact = nalarms * log(0.85 / 0.30)
    (decode(llr), round(exact; digits=2), decode(acc),
     round(post(exact); digits=5), round(post(decode(acc)); digits=5))
end
logodds_demo(12)
\end{minted}




\begin{minted}{text}
(1.0, 12.5, 8.0, 1.0, 0.99966)
# llr datum, exact log-odds, quantized log-odds, exact posterior, quantized posterior
\end{minted}



The accumulator is wrong by a third (8.0 vs 12.5 — near 8 the ulp is 2.0, so \texttt{8 + 1} rounds straight back to 8), yet the \emph{decision} is untouched: both posteriors are ≈ 1. Threshold decisions need order, not magnitude — which is why coarse formats work in reasoning pipelines far past the point where their arithmetic looks broken. When magnitude does matter, the cures are the same as for gradients: a wider accumulator format, or stochastic rounding.



\subsection{Fuzzy inference in an unsigned format}



\label{10838530668520339695}{}


Membership degrees live in [0, 1] and are never negative — a job for an unsigned finite format. \texttt{Binary8p6uf} gives a [0, 15.5] range with 1/32 resolution at 1; the classic fuzzy connectives are registry operations (\texttt{Minimum} = Gödel t-norm, \texttt{Maximum} = s-norm, \texttt{Multiply} = product t-norm):




\begin{minted}{julia}
F = Binary8p6uf
trap(x, a, b, c, d) = clamp(min((x - a) / (b - a), 1.0, (d - x) / (d - c)), 0.0, 1.0)
warm = F(trap(26.0, 15, 20, 24, 28))         # membership of 26 °C in "warm"
hot  = F(trap(26.0, 24, 30, 100, 101))       # … and in "hot"

(decode(warm), decode(hot),
 decode(Minimum(F, RNE_SatNone, warm, hot)),      # warm AND hot
 decode(Maximum(F, RNE_SatNone, warm, hot)),      # warm OR hot
 decode(Multiply(F, RNE_SatFinite, warm, hot)))   # product t-norm
\end{minted}




\begin{minted}{text}
(0.5, 0.3359375, 0.3359375, 0.5, 0.16796875)
\end{minted}



Because every connective is a bit-exact registry op, a fuzzy rule base evaluated in \texttt{Binary8p6uf} is \emph{reproducible across machines to the code point} — a property sampled float pipelines cannot promise.



\subsection{Search heuristics: what quantized scores keep is the ordering}



\label{18015521726916863751}{}


Game-tree and beam search consume evaluation scores only through comparisons. Quantizing 200 heuristic scores to 8 bits collapses them onto 78 distinct code points — yet the argmax survives, and sorting is exact (integer order keys, O(n) counting sort, single NaN ordered last deterministically):




\begin{minted}{julia}
using Random
rng = Xoshiro(21)
scores = randn(rng, 200) .* 3
q = Binary8p4se.(scores)

(argmax(scores), argmax(decode.(q)), argmax(scores) == argmax(decode.(q)),
 decode.(sort(q; rev=true)[1:5]), length(unique(codepoint.(q))))
\end{minted}




\begin{minted}{text}
(140, 140, true, [8.0, 7.5, 6.5, 6.0, 6.0], 78)
\end{minted}



The two 6.0s are the caveat: quantization creates \emph{ties} the exact scores did not have. \texttt{TotalOrder}{\textquotesingle}s deterministic tie behavior means the search expands the same node on every run — ties change which answer you get, not whether you can reproduce it.



\section{Machine Learning}



\label{11381515430689351323}{}


\subsection{Quantizing a weight tensor and measuring the damage}



\label{12845276900436285715}{}



\begin{minted}{julia}
using SmallFloats, Random, Statistics

rng = Xoshiro(42)
w = randn(rng, 10_000) .* 0.25          # typical trained-weight scale
q = Binary8p4se.(w)                     # project every weight
back = decode.(q)

mean((w .- back).^2), maximum(abs.(w .- back)), count(isinf, back) / length(back)
\end{minted}




\begin{minted}{text}
(4.41e-5, 0.0312, 0.0)                  # MSE, max |error|, overflow fraction
\end{minted}



The max error is exactly half an ulp of the largest binade the data reached — the worst case nearest rounding permits.



\subsection{Picking a format: precision vs range}



\label{9815972904090952286}{}


For the same Gaussian tensor, \texttt{K = 8} formats trade significand bits against exponent range. RMSE doubles per lost significand bit; MaxFinite grows explosively:




\begin{minted}{julia}
rng = Xoshiro(7); x = randn(rng, 50_000)
for F in (Binary8p5se, Binary8p4se, Binary8p3se, Binary8p2se)
    back = [decode(Convert(F, RNE_SatFinite, xi)) for xi in x]
    println(rpad(formatname(F), 12), " rmse = ", round(sqrt(mean((x .- back).^2)); sigdigits=3),
            "   maxfinite = ", decode(MaxFiniteOf(F)))
end
\end{minted}




\begin{minted}{text}
Binary8p5se  rmse = 0.0133   maxfinite = 15.0
Binary8p4se  rmse = 0.0266   maxfinite = 224.0
Binary8p3se  rmse = 0.0528   maxfinite = 49152.0
Binary8p2se  rmse = 0.103    maxfinite = 2.147483648e9
\end{minted}



For unit-scale data, spend bits on precision; buy range only when your data needs it (or get it from a block scale — see below).



\subsection{Stochastic rounding is unbiased where nearest is not}



\label{4196957658022129611}{}


Nearest rounding maps 0.30078125 to 0.3125 every single time — a systematic +0.0117 bias. Stochastic rounding is right \emph{on average}:




\begin{minted}{julia}
target = 0.30078125                       # ν = 3/16 of an ulp above 0.296875
σ = RSA_SatNone(16)                       # StochasticA{16} with SatNone
rng = Xoshiro(1)
m = mean(decode(Convert(Binary8p4se, σ, target; rng)) for _ in 1:200_000)
(decode(Binary8p4se(target)), m)
\end{minted}




\begin{minted}{text}
(0.3125, 0.300765)                        # RNE result vs stochastic mean ≈ target
\end{minted}



This is the property that makes stochastic rounding matter for accumulating many small contributions (gradients, activations statistics) in a low-precision format.



\section{Deep Learning}



\label{15209394077164581735}{}


\subsection{MX-style block quantization — and the staging pitfall}



\label{18175628149299408468}{}


Block formats pair a shared scale with narrow elements, tracking dynamic range that a bare element format cannot. Here each row of \texttt{W} lives at a different scale, spanning {\textasciitilde}2¹⁶ across the matrix:




\begin{minted}{julia}
using SmallFloats, Random, Statistics

rng = Xoshiro(3)
W = randn(rng, 64, 64) .* 2 .* (2.0 .^ (collect(0:63) ./ 4))   # per-row scales

function mx_rows(W, ::Type{FST}, ::Type{FS}, ::Type{FE}, ::Val{B}) where {FST,FS,FE,B}
    n, m = size(W)
    [begin
        # stage through a WIDE-RANGE format so nothing clamps before scaling
        seg = ntuple(k -> Convert(FST, RNE_SatFinite, W[i, (j-1)*B + k]), Val(B))
        ConvertToBlockMaxAbsFinite(FS, FE, RNE_SatNone, RNE_SatNone, seg)
    end for i in 1:n, j in 1:m ÷ B]
end

blocks = mx_rows(W, Binary8p2se, Binary8p1uf, Binary8p4se, Val(32))
recon = [let b = blocks[i, (j-1) ÷ 32 + 1]
             decode(b.s) * decode(b.x[(j-1) % 32 + 1])
         end for i in 1:64, j in 1:64]
plain = [decode(Convert(Binary8p4se, RNE_SatFinite, w)) for w in W]

relerr(A) = sqrt(mean(((W .- A) ./ max.(abs.(W), 1e-9)).^2))
relerr(plain), relerr(recon)
\end{minted}




\begin{minted}{text}
(0.616, 0.109)      # plain 8p4se clamps the large rows; MX tracks them
\end{minted}



\begin{tcolorbox}[toptitle=-1mm,bottomtitle=1mm,colback=admonition-warning!50!white,colframe=admonition-warning,title=\textbf{Stage wide, then block-quantize}]
\texttt{ConvertToBlockMaxAbsFinite} takes already-\texttt{Binary} elements. If you stage your \texttt{Float64} data through the \emph{element} format first, \texttt{SatFinite} clamps the large values \textbf{before} the block scale can absorb them, and blocks buy you nothing — we measured exactly that (0.617 vs 0.616) with \texttt{Binary8p4se} staging. Stage through a wide-range format (\texttt{Binary8p2se} above); the residual 0.109 here is the staging format{\textquotesingle}s own precision, which bounds what any downstream scheme can keep.

\end{tcolorbox}


At \texttt{B = 32} with an 8-bit scale the storage cost is 8.25 bits/value.



\subsection{Quantized dot products with one final rounding}



\label{1256464788428311555}{}


\texttt{BlockDotProduct} computes every lane product and the accumulation \emph{exactly}, then projects once:




\begin{minted}{julia}
rng = Xoshiro(9)
a64, b64 = randn(rng, 32), randn(rng, 32)
qb(v) = ConvertToBlockMaxAbsFinite(Binary8p1uf, Binary8p4se, RNE_SatNone, RNE_SatNone,
            ntuple(i -> Convert(Binary8p4se, RNE_SatFinite, v[i]), Val(32)))
dq = BlockDotProduct(Binary8p4se, RNE_SatNone, qb(a64), qb(b64))
(a64'b64, decode(dq))
\end{minted}




\begin{minted}{text}
(5.0634, 4.5)       # difference is input quantization only, never accumulation error
\end{minted}



\subsection{Why training loops like stochastic rounding: the swamping demo}



\label{3016318026063494718}{}


Accumulate 400 gradient steps of 0.011 into a \texttt{Binary8p3se} accumulator. Under nearest rounding, the moment the accumulator dwarfs the increment, every add rounds back to the accumulator — it \textbf{stalls}. Stochastic rounding keeps absorbing the increments in expectation:




\begin{minted}{julia}
function accumulate_demo(nsteps, g)
    σ = RSA_SatNone(16)
    rng = Xoshiro(11)
    acc_rne = Binary8p3se(0.0); acc_sto = Binary8p3se(0.0)
    for _ in 1:nsteps
        acc_rne = Add(Binary8p3se, RNE_SatNone, acc_rne, Convert(Binary8p3se, RNE_SatNone, g))
        acc_sto = Add(Binary8p3se, σ, acc_sto, Convert(Binary8p3se, σ, g; rng); rng)
    end
    (nsteps * g, decode(acc_rne), decode(acc_sto))
end
accumulate_demo(400, 0.011)
\end{minted}




\begin{minted}{text}
(4.4, 0.125, 5.0)   # exact sum, RNE accumulator (stalled!), stochastic accumulator
\end{minted}



The stochastic result is noisy (5.0 vs 4.4 on this seed) but unbiased; the RNE result is \emph{wrong by 35×} and no amount of steps will fix it. In practice you keep a wider accumulator when you can — and use stochastic rounding when you can{\textquotesingle}t.



\subsection{Activation functions are 256-byte lookup tables}



\label{12405290678054954100}{}


For pure specs, a unary activation over an 8-bit format is a \emph{finite function} — the whole nonlinearity is one 256-byte table, built bit-exactly through the scalar path and then gathered per element. This is precisely the activation LUT you would burn into an accelerator:




\begin{minted}{julia}
using SmallFloats, Random
empty_tables!()
rng = Xoshiro(4)
pre = Binary8p4se.(randn(rng, 4096) .* 2.5)       # pre-activations, ±2.5σ
act = Tanh(Binary8p4se, RNE_SatNone, pre)          # table-gather kernel

(table_bytes(),
 round(count(v -> abs(decode(v)) == 1.0, act) / 4096; digits=3),
 length(unique(codepoint.(act))),
 decode(NextLessThan(one(Binary8p4se))))
\end{minted}




\begin{minted}{text}
(256, 0.384, 100, 0.9375)
# LUT size; fraction saturated to ±1; distinct output codes; last value below 1
\end{minted}



Two things worth staring at. First, \textbf{saturation is severe at this input scale}: 38\% of pre-activations land exactly on ±1 — under \texttt{RNE}, \texttt{tanh} reaches 1 as soon as the true value rounds past the midpoint of the last gap, and the entire saturated tail becomes indistinguishable to the next layer. (Under \texttt{TowardNegative} the asymptote is honored instead: \texttt{tanh} of any finite positive input projects to 0.9375, never 1 — the projection mode is part of the activation design space.) Second, the choice of nonlinearity changes \emph{information retention} in code points, not just shape:




\begin{minted}{julia}
actS = Softplus(Binary8p4se, RNE_SatNone, pre)
(length(unique(codepoint.(actS))), table_bytes())
\end{minted}




\begin{minted}{text}
(61, 512)          # softplus keeps 61 distinct codes here; two LUTs now cached
\end{minted}



Auditing an activation for a format is this cheap: enumerate all 256 inputs, look at the output histogram, count what survives.



\subsection{Packing a quantized model}



\label{856710218023185633}{}



\begin{minted}{julia}
n = 1_000_000
model = [rawvalue(Binary5p2se, UInt8(rand(0:31))) for _ in 1:n]
pv = PackedVector(model)
(sizeof(model), sizeof(pv.data))
\end{minted}




\begin{minted}{text}
(1000000, 625000)   # 8 bits → 5 bits per value; indexing and vmap work directly on pv
\end{minted}



\part{Julia Compatibility}


\chapter{Julia Compatibility}



\label{10344462847155476790}{}


\texttt{Binary} values speak ordinary Julia. The \textbf{Base register} (\texttt{src/juliacompat.jl}) maps every draft operation that has a Base counterpart onto its Base name, so generic Julia code — and your fingers — can use the spellings they already know. Every method there is exactly one same-format spec-register call under the session default projection; there is no third semantics: \texttt{x + y} \emph{is} \texttt{Add(x, y)}, which \emph{is} \texttt{Add(T, DefaultProjection(), x, y)}.



\section{Using the overloads}



\label{5063837776835138997}{}



\begin{minted}{jlcon}
julia> x, y, z = Binary8p4se(1.6), Binary8p4se(0.25), Binary8p4se(2.0);

julia> x + y                       # Add
Binary8p4se(1.875 ≡ 0x47)

julia> exp(y)                      # Exp — a 256-byte table lookup once warm
Binary8p4se(1.25 ≡ 0x42)

julia> atan(y, x)                  # ArcTan2, in Base's (y, x) argument order
Binary8p4se(0.15625 ≡ 0x2a)

julia> fma(x, y, z)                # FMA — one rounding; muladd is the same op
Binary8p4se(2.5 ≡ 0x4a)

julia> sincos(y)                   # composite: (Sin(y), Cos(y)) componentwise
(Binary8p4se(0.25 ≡ 0x30), Binary8p4se(1.0 ≡ 0x40))
\end{minted}



Because the veneers follow the session default projection, changing it changes them (see the session-defaults section of the \hyperlinkref{10272851679537904033}{User Guide}):




\begin{minted}{jlcon}
julia> Binary8p4se(200.0) * Binary8p4se(2.0)     # RNE_SatNone: overflow → Inf
Binary8p4se(Inf ≡ 0x7f)

julia> DefaultProjection!(RTZ_SatFinite);

julia> Binary8p4se(200.0) * Binary8p4se(2.0)     # now clamps to MaxFinite
Binary8p4se(224.0 ≡ 0x7e)
\end{minted}



Code that must be insensitive to the session default names its projection: \texttt{Multiply(T, RNE\_SatNone, x, y)}.



\section{The mapping}



\label{12864072587101288386}{}



\begin{table}[h]
\centering
\begin{tabulary}{\linewidth}{R R}
\toprule
Base spelling & draft operation \\
\toprule
\texttt{+} \texttt{-} \texttt{*} \texttt{/}, unary \texttt{-} & \texttt{Add}, \texttt{Subtract}, \texttt{Multiply}, \texttt{Divide}, \texttt{Negate} \\
\texttt{abs}, \texttt{inv}, \texttt{sqrt} & \texttt{Abs}, \texttt{Recip}, \texttt{Sqrt} \\
\texttt{exp}, \texttt{exp2}, \texttt{expm1}, \texttt{log}, \texttt{log2}, \texttt{log1p} & \texttt{Exp}, \texttt{Exp2}, \texttt{ExpMinusOne}, \texttt{Log}, \texttt{Log2}, \texttt{LogOnePlus} \\
\texttt{sin}, \texttt{cos}, \texttt{tan}, \texttt{asin}, \texttt{acos}, \texttt{atan} & \texttt{Sin}, \texttt{Cos}, \texttt{Tan}, \texttt{ArcSin}, \texttt{ArcCos}, \texttt{ArcTan} \\
\texttt{sinh}, \texttt{cosh}, \texttt{tanh}, \texttt{asinh}, \texttt{acosh}, \texttt{atanh} & \texttt{Sinh}, \texttt{Cosh}, \texttt{Tanh}, \texttt{ArcSinh}, \texttt{ArcCosh}, \texttt{ArcTanh} \\
\texttt{sinpi}, \texttt{cospi}, \texttt{tanpi} & \texttt{SinPi}, \texttt{CosPi}, \texttt{TanPi} \\
\texttt{atan(y, x)} & \texttt{ArcTan2} (Base{\textquotesingle}s argument order) \\
\texttt{copysign}, \texttt{hypot} & \texttt{CopySign}, \texttt{Hypot} \\
\texttt{min}, \texttt{max} & \texttt{Minimum}, \texttt{Maximum} (NaN-propagating, exactly Base{\textquotesingle}s float semantics) \\
\texttt{fma}, \texttt{muladd} & \texttt{FMA} (both: one rounding) \\
\texttt{clamp} & \texttt{Clamp} \\
\texttt{sincos}, \texttt{sincospi}, \texttt{minmax} & componentwise composites of the draft ops \\
\bottomrule
\end{tabulary}

\end{table}



The mapping is a declarative partition over the op lists, and the test suite asserts it is exhaustive: every operation is either mapped above or listed in \texttt{\_NO\_BASE\_COUNTERPART}. Adding an operation to the registry forces an explicit decision here.



\section{What is \emph{not} mapped, and why}



\label{2392765366633338434}{}


\textbf{Draft operations with no Base spelling} — call them by their draft names: \texttt{RSqrt}, \texttt{Softplus}, \texttt{ArcSinPi}/\texttt{ArcCosPi}/\texttt{ArcTanPi}/\texttt{ArcTan2Pi} (Base has only the \texttt{sinpi} family), the NaN-ignoring/magnitude/finite extremum families (\texttt{MinimumNumber}, \texttt{MaximumMagnitude}, \texttt{MinimumFinite}, …; Base has no NaN-ignoring pair), and \texttt{FAA} (no Base fused add-add).




\begin{minted}{jlcon}
julia> min(Binary8p4se(NaN), x)          # Base semantics: NaN propagates
Binary8p4se(NaN ≡ 0x80)

julia> MinimumNumber(Binary8p4se(NaN), x)   # the NaN-ignoring draft op
Binary8p4se(1.625 ≡ 0x45)
\end{minted}



\textbf{Mixed \texttt{Binary} formats} — deliberately a \texttt{promotion ... failed to change} error, never a silent widening. Mixing formats is an explicit \texttt{Convert}:




\begin{minted}{jlcon}
julia> x + Convert(Binary8p4se, RNE_SatNone, Binary5p3sf(1.0))
Binary8p4se(2.5 ≡ 0x4a)
\end{minted}



\textbf{\texttt{Binary} with ordinary numbers} — promotes to \texttt{Float64} (the exact carrier) through the rules in \texttt{formats.jl}, so the result is a \texttt{Float64}, not a re-projected \texttt{Binary}:




\begin{minted}{jlcon}
julia> x + 2.0
3.625
\end{minted}



Project it back explicitly when that is what you mean: \texttt{Convert(Binary8p4se, RNE\_SatNone, decode(x) + 2.0)}.



\section{The rest of the Base surface}



\label{15708086746009923714}{}


Beyond arithmetic, \texttt{Binary} integrates where the other layers provide it: predicates and constants (\texttt{isnan}, \texttt{isfinite}, \texttt{signbit}, \texttt{zero}, \texttt{one}, \texttt{eps}, \texttt{typemin}/\texttt{typemax}, \texttt{floatmin}/\texttt{floatmax}), comparisons and \texttt{isless} on integer order keys, \texttt{sort} via an O(n) counting sort, \texttt{nextfloat}/\texttt{prevfloat} as the draft Next operations, \texttt{codepoint}, and \texttt{rand}/\texttt{randn} (see the \hyperlinkref{10272851679537904033}{User Guide}). All of it goes through the same projection engine and decode tables as the draft-named API.



\part{Technical Guide}


\chapter{Technical Guide}



\label{4463396622184713734}{}


How SmallFloats.jl works inside: the encoding, the projection engine, the oracle and its correctness protocol, the performance layers, and the verification doctrine that holds it all together. Read the \hyperlinkref{10272851679537904033}{User Guide} first; this page assumes its vocabulary.



\section{Layer map}



\label{9668567837796074949}{}


Source files load in dependency order, each layer speaking only downward:




\begin{table}[h]
\centering
\begin{tabulary}{\linewidth}{R R R}
\toprule
layer & file & provides \\
\toprule
formats & \texttt{formats.jl} & \texttt{Binary\{K,P,SGN,EXT\}}, the 120 named aliases, Group M queries \\
specs & \texttt{projspec.jl} & rounding/saturation singletons, \texttt{ProjSpec\{R,S\}}, the predefined spec grid \\
defaults & \texttt{defaults.jl} & settable session defaults (\texttt{DefaultType}, \texttt{DefaultProjection}, …) behind \texttt{Ref}s; consumed via the speculation guard \\
codec & \texttt{decode\_encode.jl} & decode (generated tables + bit-composed compute), encode, order keys, counting sort, \texttt{Class}, Next ops \\
engine & \texttt{project.jl} & \texttt{round\_to\_precision} (mask-based Float64 core + generic core), \texttt{saturate}, \texttt{project}, \texttt{project\_interval} \\
ops & \texttt{ops\_scalar.jl} & result-kind protocol, \texttt{apply\_op}, the operation registry, the spec register \\
compat & \texttt{juliacompat.jl} & the Base register: declarative op ⇒ Base-function map \\
oracle & \texttt{oracle.jl} & ω-semantics for all 52 operations \\
tables & \texttt{tables.jl} & the pure-ρ result-table cache (unary/binary + bitwidth-gated ternary) \\
kernels & \texttt{kernels.jl} & Shape-A gathers, Shape-B scalar/threaded loops, \texttt{vmap} \\
blocks & \texttt{blocks.jl} & \texttt{Block}, block/scaled ops, exact reductions \\
packed & \texttt{packed.jl} & sub-byte \texttt{PackedVector} storage \\
approx & \texttt{approx.jl} & κ measurement/registry, conformance declaration \\
rand & \texttt{rand.jl} & Random-API hooks: uniform-[0,1) floor projection, clamped normal \\
\bottomrule
\end{tabulary}

\end{table}



\section{Encoding and decoding}



\label{8606053261248653135}{}


A value is its code point (\texttt{UInt8}). The bit layout is the draft{\textquotesingle}s: sign (signed formats), biased exponent, trailing significand; one NaN at the negative-zero slot; no −0; ±Inf adjacent to the extremes in extended formats.



\texttt{decode} is a \textbf{\texttt{@generated} constant-tuple lookup}: per format, a \texttt{2{\textasciicircum}K}-tuple of \texttt{Float64} datums built once from the computational decode (so table and computation are correct by construction and asserted equivalent exhaustively). Constant inputs still fold — \texttt{maxfinite\_datum(T)} is a compile-time constant — while runtime decode is a single indexed load (≈ 0.7 ns). The computational decode assembles the Float64 \textbf{by bits} (normalize with \texttt{leading\_zeros}, place exponent and mantissa fields, \texttt{reinterpret}), valid because every datum{\textquotesingle}s exponent sits deep inside Float64{\textquotesingle}s normal range; \texttt{Float64} is thereby the \emph{exact} carrier for all datums, an invariant the suite checks against an independent big-float transliteration of the draft.



Ordering runs on \textbf{integer order keys}: a sign-magnitude fold into \texttt{UInt16}, monotone with the total order, NaN mapped to \texttt{typemax}. Same-format \texttt{TotalOrder}, \texttt{isless}, and the numeric comparisons are key comparisons ({\textasciitilde}1 ns); since a format has at most \texttt{2{\textasciicircum}K + 1} distinct keys, vectors sort with an \textbf{O(n) counting sort} installed via \texttt{Base.Sort.defalg} (forward and reverse orderings; anything exotic falls back to the stock algorithm).



\section{The projection engine}



\label{8339828346285081570}{}


\texttt{project(fr, ρ, X; R, sticky)} is the single write path into a code point:




\begin{minted}{text}
RoundToPrecision  →  Saturate  →  Encode
\end{minted}



\textbf{RoundToPrecision} produces a \texttt{Rounded(kind, sign, S, Q)} — an exact scaled significand \texttt{S} and exponent \texttt{Q} per the draft{\textquotesingle}s \texttt{Q = max(⌊log₂|X|⌋, 1−B) − P + 1}. Two implementations, proven equivalent exhaustively:



\begin{itemize}
\item the \textbf{generic core} (\texttt{\_rtp\_core}) works on any carrier (\texttt{Float64}, \texttt{Float128}, \texttt{BigFloat}) via exact power-of-two scaling and a fraction ν compared against the mode{\textquotesingle}s decision points;


\item the \textbf{mask-based Float64 core} (\texttt{\_rtp\_f64}) extracts sign/exponent/significand fields directly, represents ν as a 128-bit fixed-point integer with an OR-mask sticky for bits shifted out, and evaluates every mode — including the stochastic \texttt{⌊ν·2{\textasciicircum}N⌋} comparisons and \texttt{RNITE} ties — as integer field tests. Specials, zeros, and (Convert-only-reachable) subnormal Float64 inputs bail to the generic core.

\end{itemize}


\textbf{Symbolic sticky} is how enclosures talk to the engine: \texttt{sticky ∈ \{−1, 0, +1\}} declares the true value to be the carried value plus an infinitesimal of that sign. The engine folds it into every comparison; the delicate case — true value just \emph{below} a representable dyadic — decrements into the previous binade and sets ν = 1⁻ (encoded in the mask core as an all-ones fixed-point fraction, which reproduces every predicate including the RNITE tie behavior). This is what lets directed modes land exactly on asymptotes: \texttt{tanh → 1⁻} under \texttt{TowardNegative} projects to \texttt{NextLessThan(1)}, automatically.



\textbf{Saturate} classifies against the format{\textquotesingle}s range with two integer comparisons per side, then maps the classification through the draft{\textquotesingle}s twenty saturation rows to \texttt{as-is / MaxFinite / MinFinite / ±Inf / NaN}. The comparisons stay cheap because:



\begin{itemize}
\item the extremal magnitude in canonical \texttt{(S, Q)} form is a constant-folded function of the type parameters;


\item the rounded value and the extremum share one lexicographic \texttt{(Q, S)} order (subnormals and the lowest normal binade share the same \texttt{Q});


\item signed formats use sign–magnitude symmetry, and unsigned underflow is just the sign bit;


\item an internal \texttt{HUGEQ} sentinel represents {\textquotedbl}finite but astronomically large{\textquotedbl} (the ν = 1⁻ image of an infinite endpoint), so directed \texttt{SatNone} clamps it to MaxFinite correctly.

\end{itemize}


\textbf{Encode} is pure integer bit assembly, including the significand-carry renormalization and the subnormal/normal field split.



\section{The oracle and the result-kind protocol}



\label{9066011901059263514}{}


Every operation{\textquotesingle}s defined result is computed by \texttt{ωeval}, which returns one of five result kinds; \texttt{apply\_op} fast-splits the common one and finishes the rest:




\begin{table}[h]
\centering
\begin{tabulary}{\linewidth}{R R R}
\toprule
kind & meaning & finished by \\
\toprule
\texttt{Float64} & exact (specials; representable arithmetic) & direct \texttt{project} \\
\texttt{Float128} & exact by \textbf{width analysis} & direct \texttt{project} (Float128 carrier) \\
\texttt{StickyF} & wide-spread \texttt{FMA}/\texttt{FAA} tail: head value (\texttt{Float64} or \texttt{Float128}) + tail sign & direct \texttt{project} with \texttt{sticky} set — no allocation \\
\texttt{BigExactF} & exact at 2200-bit precision (wide-spread tail for \texttt{Add}; the near-impossible \texttt{FAA} distillation miss) & \texttt{project} on \texttt{BigFloat} \\
\texttt{Enclose128F} & correctly-rounded Float128 \textbf{bracket} & sticky agreement, MPFR fallback \\
\texttt{EncloseF} & MPFR directed enclosure \texttt{f(prec)}, optional Float128 pre-filter \texttt{fq}, optional eager Float64 estimate \texttt{yd} & three-stage: \texttt{yd} → \texttt{fq} → interval protocol \\
\bottomrule
\end{tabulary}

\end{table}



Two \textbf{rigor classes} govern every non-\texttt{Float64} path, and their arguments are never mixed:



\textbf{Class R (unconditional)} rests on two facts that hold without any accuracy assumption:



\begin{itemize}
\item \textbf{Width analysis.} Sums of decoded datums are \emph{exactly representable} in \texttt{Float128} whenever operand bits + exponent spread fit 113 bits — checked in advance by integer exponent arithmetic (\texttt{Add} at ΔE ≤ 100, \texttt{FMA} ≤ 92, \texttt{FAA} span ≤ 98). Beyond the threshold, \texttt{Add} escalates to the 2200-bit path, while \texttt{FMA}/\texttt{FAA} take a non-allocating \textbf{sticky-head} shortcut: past that spread the smaller term is provably too small (by construction of the threshold) to affect anything but the tail direction of the larger one, so the result is \texttt{StickyF(head, sign)} — no \texttt{BigFloat} involved. \texttt{FAA}{\textquotesingle}s three-term case runs a bounded Float128 2sum-distillation (Priest-style, ≤ 6 sweeps) to find that head/tail split directly; the residual \texttt{BigExactF} MPFR fallback exists only for the near-impossible case the distillation doesn{\textquotesingle}t converge.


\item \textbf{IEEE correct rounding.} Division and \texttt{sqrt} are correctly rounded at \emph{every} binary width, by mandate. \texttt{Divide}/\texttt{Recip} therefore use the \textbf{Float64} quotient directly: it is within half an ulp of truth, so it serves as \texttt{EncloseF}{\textquotesingle}s eager \texttt{yd} estimate with no Float128 arithmetic on the path at all. \texttt{Sqrt}/\texttt{RSqrt} use the \textbf{Float128} CR result: an inexact nearest-CR value \texttt{q} brackets the truth in the \emph{open} interval \texttt{(prevfloat(q), nextfloat(q))}. Exactness itself is detected by an \texttt{fma} residual test.

\end{itemize}


\textbf{Class E (envelope-conditional).} Faithful-but-not-CR evaluations stand in for an enclosure only inside a generous envelope, resolved by the two-sided sticky gate. Two stages, cheapest first:



\begin{itemize}
\item \emph{Float64 stage:} for operations whose Float64 libm is faithful (≤ 1 ulp ≈ 2⁻⁵² relative — the exp/log families, the hyperbolics, forward trig inside the |x| ≤ 10¹⁵ reduction window, \texttt{atan}, \texttt{asinh}, the softplus composition), the eager estimate \texttt{yd} carries envelope \texttt{E = |yd|·2⁻⁴⁵}, ≥ 2⁷ slacker than any faithful-libm error (measured margin on this machine: ≥ 2⁶·⁶).


\item \emph{Float128 stage:} libquadmath{\textquotesingle}s estimate \texttt{y = fq()} carries \texttt{E = |y|·2⁻⁹⁰}, at least 2¹⁸ slacker than any published libquadmath bound.

\end{itemize}


In both stages, if the sticky-projected endpoints agree, that code point is the answer; if not, the next stage (ultimately the MPFR ladder) decides. An envelope failure can therefore only cost speed, never correctness — unless both endpoints agree \emph{and} the envelope is wrong, which the differential-build tests rule out empirically by building every standard table twice (Float128-first and pure-MPFR) and byte-diffing, and which the exhaustive oracle cross-check re-verifies per operation against the 3072-bit protocol.



\textbf{The interval protocol} (\texttt{project\_interval}) is the termination backbone: evaluate the MPFR enclosure at precision \texttt{p}; if \texttt{lo == hi} the value is exactly representable — project it; otherwise project \texttt{lo} with \texttt{sticky = +1} and \texttt{hi} with \texttt{sticky = −1}; agreement (projection is monotone at fixed \texttt{R}) yields the answer; disagreement doubles \texttt{p}, clamped to the \texttt{maxprec} ceiling (default 4096 bits, honored exactly even when it is not a power-of-two multiple of the 256-bit start). Termination requires that the enclosure not be chasing a value the grid can actually hit — which is why the π-scaled operations carry \textbf{Niven peels}: for dyadic arguments, \texttt{tan(πr)} takes rational values only at the quarter-integers (±1) and \texttt{atan2/π} only on the diagonals (±¼, ±¾); those cases are answered exactly before any enclosure is built, and Niven{\textquotesingle}s theorem proves the peel set complete.



\section{Tables and kernels}



\label{17163678065643276217}{}


For pure specs, unary and binary operations are \textbf{finite functions} — 256 or 65 536 entries — so the kernel layer materializes them once per \texttt{(op, formats, ρ)} into a locked cache (\texttt{Dict\{TableKey, Memory\{UInt8\}\}}, double-checked locking, builds outside the lock) and serves every later array call as a gather: Shape-A, one load per element, measured 0.27 ns/elem unary and 0.5 ns/elem binary. Tables are built \emph{through the scalar path}, so they inherit its bit-exactness; the suite asserts table ≡ scalar over every entry.



Ternary (\texttt{FMA}, \texttt{FAA}, \texttt{Clamp}) is a finite function too — 2{\textasciicircum}(K1+K2+K3) entries — but that count spans four orders of magnitude across the 3–8 bitwidth range (512 B at K=3, 16 MiB at K=8), so one policy doesn{\textquotesingle}t fit the whole range. A separate \texttt{TernaryKey → TernaryEntry} cache (\texttt{\_ternary\_table\_for}, in \texttt{tables.jl}) tiers by Σ bitwidth:



\begin{itemize}
\item \textbf{Eager} (≤ 18 bits, all \texttt{K ≤ 6} combinations, ≤ 256 KiB): builds and caches on the first array call.


\item \textbf{Adaptive} (≤ 21 bits, the \texttt{K = 7} band, up to 2 MiB): accumulates a per-signature element count across calls and builds only once a signature has processed enough elements to amortize the build; a byte-bounded LRU eviction (\texttt{TERNARY\_CACHE\_BYTES}) guards against many hot signatures coexisting.


\item \textbf{Compute} (\texttt{K = 8}, 16+ MiB — a table is never worth it): \texttt{vmap!} runs Shape-B, the fully specialized scalar pipeline per element, optionally split across \texttt{Threads.@threads} for long enough arrays (each ternary draw is independent under a fixed, non-stochastic ρ, so lanes cannot interact).

\end{itemize}


Every ternary table entry, eager or adaptive, is still built \emph{through the scalar path} — the tiering changes when/whether the cache exists, never what it contains. Stochastic calls of any arity always take Shape-B, with the RNG resolved once per array rather than per element.



\section{Blocks: exactness without a superaccumulator}



\label{10094224062570944436}{}


\texttt{blockdecode} produces each lane{\textquotesingle}s \texttt{scale × element} exactly (≤ 17-bit significands in Float64). Reductions then apply \textbf{span filters}: one integer pass over lane exponents decides whether the whole sum (or dot product, with ≤ 33-bit lane products) is exactly representable in \texttt{Float128}; if so, a plain \texttt{Float128} accumulation \emph{is} the exact answer; if not, an exact big-float accumulation takes over. Either way there is exactly one projection, at the end. \texttt{ωBlockProject} follows the draft{\textquotesingle}s special rows for the scale (NaN, 0, ±Inf) and divides each element result by the scale through its own cheapest-first CR-bracket / enclosure cascade (exact Float64 quotient → CR Float128 → bracket/pre-filter → MPFR interval), mirroring the scalar quotient group{\textquotesingle}s rigor arguments.



\section{The κ registry and conformance}



\label{9937209558637538978}{}


κ is the maximum code-point distance (along the total order) between an implementation{\textquotesingle}s result and the defined result, over inputs with finite defined results; any mismatch on non-finite defined results makes κ = NaN. Because inputs are enumerable, \texttt{register\_approx!} \emph{measures} κ exhaustively at registration and rejects understatement — a declared bound is a verified property, not a promise. \texttt{conformance()} assembles the declaration live from the operation registry, the table cache, and the approximation registry.



\section{Verification doctrine}



\label{12691031809522550080}{}


The value sets are small enough that sampling is never necessary, so the suite enumerates — ≈ 8.9 M assertions in all:



\begin{itemize}
\item formats against an independent draft transliteration (14 679);


\item ordering over all 2.5 M same-format pairs plus Next-op edge tables (7.6 M);


\item every unary operation on every input against a 3072-bit protocol run; divide and the ternaries exhaustively;


\item stochastic R-sweeps with directed-asymptote pins;


\item table ≡ scalar over every entry, including the ternary tiers (eager and adaptive) against the scalar path;


\item the sticky-head \texttt{FMA}/\texttt{FAA} escalation against the MPFR reference across every rounding-mode family and adversarial cancellation cases;


\item blocks against a from-scratch reference composition;


\item Float128 carrier ≡ Float64 carrier, and Float128-first ≡ MPFR differential builds;


\item the mask rounding core ≡ the generic core over datums and reachable sums/products at boundary stochastic budgets N ∈ \{45, 60\};


\item packed round-trips, and κ/conformance behavior.

\end{itemize}


Deterministic \textbf{specialization regressions} (concrete inferred return types at the public entry points, zero warm-path allocation) stand in for timing assertions.



\section{Benchmark doctrine}



\label{14987590650541321188}{}


Recorded after two measurement post-mortems: a benchmark closure over any non-\texttt{const} global measures Julia{\textquotesingle}s dispatch machinery, not the code under test — and it distorts \emph{ratios}, not just absolutes, because dispatch cost varies with call shape (a dynamic keyword call costs {\textasciitilde}1 µs; a dynamic positional call far less; six unresolved interior sites cost six times one).



The rules the shipped \texttt{benchmarking/benchmarking.jl} enforces structurally:



\begin{itemize}
\item format types enter as type parameters, never as globals;


\item operands come from untimed setup;


\item functions retrieved reflectively pass through argument barriers to specialize;


\item a preflight aborts the run if warm scalar paths allocate — including the wide-spread \texttt{FMA}/\texttt{FAA} sticky-head path.

\end{itemize}


The table-build section reports both the cold build (cache evicted per sample) and the steady-state warm cache hit, since callers amortize the former through the latter. Vary one binding per variant; verify specialization before believing a number.



\section{Deliberate limitations}



\label{16229184572553167566}{}


No implicit cross-format arithmetic (promotion is to \texttt{Float64}, explicitly). No in-place packed arithmetic. Threading is opt-in and narrow: only the untabled ternary (\texttt{K = 8}) compute kernel threads, and only above a size cutoff and when \texttt{Threads.nthreads() > 1}; every other kernel is single-threaded. \texttt{Irrational}/\texttt{Rational} inputs to \texttt{Convert} are rejected rather than double-rounded silently. The \texttt{Float128} machinery never changes results — disabling it (\texttt{SmallFloats\_Float128=disable}) is a tested no-op semantically.



\part{Technical Examples}


\chapter{Technical Examples}



\label{4476403263922386643}{}


Internals-level recipes: introspecting the pipeline, verifying custom code against the oracle, measuring κ, and benchmarking without fooling yourself. Outputs are captured from real sessions. Several examples use unexported internals (\texttt{SmallFloats.round\_to\_precision}, \texttt{SmallFloats.project}, …); those are stable enough to learn from but are not covered by semantic-versioning guarantees.



\section{Basic}



\label{13551325986707881962}{}


\subsection{Watching RoundToPrecision work}



\label{2755073532632396567}{}


\texttt{round\_to\_precision} returns the draft{\textquotesingle}s exact \texttt{(sign, S, Q)} decomposition before saturation ever sees it:




\begin{minted}{julia}
using SmallFloats
using SmallFloats: round_to_precision

r = round_to_precision(4, 8, NearestTiesToEven(), 2.30078125, 0, 0)   # P=4, bias=8
(r.sign, r.S, r.Q, r.S * 2.0^r.Q)
\end{minted}




\begin{minted}{text}
(1, 9, -2, 2.25)          # S·2^Q = 9/4: the nearest P=4 value
\end{minted}



\subsection{Symbolic sticky: projecting {\textquotedbl}just below{\textquotedbl} a number}



\label{1121493031502271068}{}


The engine accepts \texttt{sticky ∈ \{-1, 0, +1\}} meaning {\textquotedbl}the true value is the carrier plus an infinitesimal of this sign{\textquotedbl}. This is how enclosure endpoints and asymptotes (\texttt{tanh → 1⁻}) are projected exactly:




\begin{minted}{julia}
using SmallFloats: project
project(Binary8p4se, ProjSpec(TowardNegative(), SatNone()), 1.0; sticky=-1)
\end{minted}




\begin{minted}{text}
Binary8p4se(0.9375 ≡ 0x3f)     # == NextLessThan(1.0): the engine crossed the binade
\end{minted}



\subsection{Tables are the scalar path, memoized}



\label{17860062526638989794}{}



\begin{minted}{julia}
using SmallFloats: get_table
empty_tables!()
tbl = get_table(:Exp, Binary8p4se, Binary8p4se, RNE_SatNone)
(tbl[Int(0x45) + 1],
 codepoint(Exp(Binary8p4se, RNE_SatNone, rawvalue(Binary8p4se, 0x45))),
 table_bytes())
\end{minted}




\begin{minted}{text}
(0x52, 0x52, 256)              # table entry ≡ scalar result; one 256-byte table cached
\end{minted}



\subsection{Order keys}



\label{11135872524320346846}{}


Ordering is integer arithmetic: a sign-magnitude fold, monotone with the total order, NaN at the top. This is what comparisons and the O(n) counting sort run on:




\begin{minted}{julia}
using SmallFloats: order_key
[(v, order_key(v)) for v in (Binary8p4se(-1.0), Binary8p4se(0.0),
                             Binary8p4se(1.0), Binary8p4se(NaN))]
\end{minted}




\begin{minted}{text}
-1.0 → 64      0.0 → 129      1.0 → 193      NaN → 65535
\end{minted}



\section{General AI}



\label{12995498770613901834}{}


\subsection{Ranking safety: an exhaustive monotonicity audit of score conversion}



\label{8732453457535267162}{}


Decision and ranking systems consume scores through \emph{order}. If scores are produced in one format and compared in another, the conversion must be monotone — and formats are small enough to prove it by enumeration rather than trust it. Walk every \texttt{Binary8p3se} datum in total order and check that its \texttt{Binary8p4se} image never goes backward on the integer order keys:




\begin{minted}{julia}
using SmallFloats
using SmallFloats: order_key

function monotone_conversion(::Type{From}, ::Type{To}) where {From,To}
    prev = nothing
    for v in sort(From.(0x00:UInt8(2^bitwidth(From) - 1)))
        isnan(decode(v)) && continue
        g = Convert(To, RNE_SatNone, v)
        prev !== nothing && order_key(g) < order_key(prev) && return false
        prev = g
    end
    true
end
monotone_conversion(Binary8p3se, Binary8p4se)
\end{minted}




\begin{minted}{text}
true
\end{minted}



Run this for every \texttt{(From, To, ρ)} triple your system actually uses; it is a few hundred integer comparisons per triple. A ranking pipeline whose conversions all pass this audit cannot invert a preference by changing formats — a guarantee no amount of spot-testing provides.



\subsection{κ-safe decision margins for approximate evaluators}



\label{7954056064411327915}{}


Search under a compute budget often wants a cheap, approximate evaluation function. The κ registry turns {\textquotedbl}how approximate?{\textquotedbl} into a \emph{measured} code-point bound — and code-point bounds compose into a decision rule: \textbf{if two defined evaluations differ by more than 2κ code points along the total order, the approximate evaluator cannot invert their comparison.} Verify the rule exhaustively for a κ = 2 evaluator:




\begin{minted}{julia}
using SmallFloats: order_key

fast(x) = (r = Exp(Binary8p4se, RNE_SatNone, x);
           isfinite(decode(r)) ? NextGreaterThan(NextGreaterThan(r)) : r)
κ, exhaustive = measure_kappa(fast, :Exp, Binary8p4se, (Binary8p4se,), RNE_SatNone)

function margin_audit(fast, κ)
    codes = [rawvalue(Binary8p4se, UInt8(c)) for c in 0:255]
    safe = violations = 0
    for a in codes, b in codes
        da, db = Exp(Binary8p4se, RNE_SatNone, a), Exp(Binary8p4se, RNE_SatNone, b)
        (isnan(decode(da)) || isnan(decode(db))) && continue
        codedistance(da, db) > 2κ || continue
        safe += 1
        (order_key(fast(a)) < order_key(fast(b))) ==
            (order_key(da) < order_key(db)) || (violations += 1)
    end
    (safe, violations)
end
(κ, exhaustive, margin_audit(fast, κ)...)
\end{minted}




\begin{minted}{text}
(2.0, true, 49522, 0)
\end{minted}



49,522 operand pairs clear the 2κ margin, and the approximate evaluator agrees with the defined ordering on every one of them — zero violations, exhaustively. Inside the margin, comparisons are genuinely undecidable at this κ; a search can treat sub-margin comparisons as ties to expand, or escalate those few nodes to the exact evaluator. Either way the pruning is \emph{provably} sound, with κ measured at registration rather than promised.



\section{Machine Learning}



\label{11381515430689351323}{}


\subsection{Verifying a quantizer exhaustively against an independent reference}



\label{15080892262823953367}{}


Formats are small enough that {\textquotedbl}test on a few points{\textquotedbl} is never necessary. Here every in-range datum of \texttt{Binary8p3se} is projected into \texttt{Binary8p4se} and checked against a brute-force nearest search under 256-bit arithmetic (distance agreement; the draft{\textquotesingle}s tie rule is the projection engine{\textquotesingle}s job and is pinned elsewhere in the suite):




\begin{minted}{julia}
using SmallFloats

function ref_nearest_distance(::Type{T}, x) where {T}
    fins = [v for v in (rawvalue(T, UInt8(c)) for c in 0:255) if isfinite(decode(v))]
    minimum(abs(setprecision(() -> BigFloat(decode(v)) - BigFloat(x), BigFloat, 256))
            for v in fins)
end

function verify_quantizer()
    ok = true
    for c in 0x00:0xff
        x = decode(rawvalue(Binary8p3se, c))
        (isfinite(x) && abs(x) <= decode(MaxFiniteOf(Binary8p4se))) || continue
        got = Convert(Binary8p4se, RNE_SatNone, x)
        ok &= abs(decode(got) - x) == Float64(ref_nearest_distance(Binary8p4se, x))
    end
    ok
end
verify_quantizer()
\end{minted}




\begin{minted}{text}
true
\end{minted}



This pattern — enumerate the inputs, compare against an independently written reference — is how the entire package is tested, and it is available to \emph{your} quantization code at trivial cost.



\subsection{The κ workflow, including the part where it says no}



\label{64156751436703573}{}


Register an approximate kernel and the registry measures it; understate the bound and it refuses:




\begin{minted}{julia}
step2(x) = (r = Exp(Binary8p4se, RNE_SatNone, x);
            isfinite(decode(r)) ? NextGreaterThan(NextGreaterThan(r)) : r)

measure_kappa(step2, :Exp, Binary8p4se, (Binary8p4se,), RNE_SatNone)
\end{minted}




\begin{minted}{text}
(2.0, true)                # κ = 2 code points, verified over all 256 inputs
\end{minted}




\begin{minted}{julia}
register_approx!(:cheater, :Exp, Binary8p4se, (Binary8p4se,), RNE_SatNone, step2; κ=1)
\end{minted}




\begin{minted}{text}
ERROR: ArgumentError: declared κ = 1.0 understates measured κ = 2.0 — registration rejected
\end{minted}



\subsection{Exporting the conformance declaration}



\label{14596415327394333925}{}



\begin{minted}{julia}
d = conformance_dict()          # plain nested Dict{String,Any}
(d["package"], length(d["formats"]), length(d["operations"]),
 [s["op"] for s in d["cached_specializations"]])
\end{minted}



Serialize \texttt{d} with any JSON/TOML writer to attach a machine-readable conformance statement to experiment artifacts.



\section{Deep Learning}



\label{15209394077164581735}{}


\subsection{Proving your fused kernel exact: BlockDotProduct vs 512-bit truth}



\label{13954472669056568806}{}


The block layer promises \emph{one} projection with exact lane arithmetic. Trust, then verify — against big-float truth, over random blocks with mixed scales and an honest special-value mix:




\begin{minted}{julia}
using SmallFloats, Random

rng = Xoshiro(5)
mk() = Block(Binary8p1uf(2.0^rand(rng, -3:3)),
             ntuple(_ -> rawvalue(Binary8p4se, UInt8(rand(rng, 0:255))), 32))

function verify_dots(trials)
    for _ in 1:trials
        bx, by = mk(), mk()
        lx = [decode(bx.s) * decode(v) for v in bx.x]
        ly = [decode(by.s) * decode(v) for v in by.x]
        (any(!isfinite, lx) || any(!isfinite, ly)) && continue
        truth = setprecision(() -> sum(BigFloat(lx[i]) * BigFloat(ly[i]) for i in 1:32),
                             BigFloat, 512)
        got = BlockDotProduct(Binary8p4se, RNE_SatNone, bx, by)
        ref = setprecision(() -> SmallFloats.project(Binary8p4se, RNE_SatNone, truth), BigFloat, 512)
        codepoint(got) == codepoint(ref) || return false
    end
    true
end
verify_dots(200)
\end{minted}




\begin{minted}{text}
true
\end{minted}



The same shape verifies any custom fused kernel you write: compose the truth in \texttt{BigFloat}, project once, compare code points.



\subsection{Stochastic rounding, audited: the full-R sweep}



\label{6013932738708787830}{}


Every stochastic projection is a deterministic function of the draw \texttt{R}, so its \emph{distribution} is checkable exactly. For \texttt{x = 2 + 3/64} in \texttt{Binary8p4se} the fraction is ν = 3/16 of an ulp, so \texttt{StochasticA\{4\}} must round up for exactly 3 of the 16 draws:




\begin{minted}{julia}
σ4 = RSA_SatNone(4)                  # ≡ ProjSpec(StochasticA{4}(), SatNone())
x = 2.0 + 3/64
count(decode(SmallFloats.project(Binary8p4se, σ4, x; R)) == 2.25 for R in 0:15)
\end{minted}




\begin{minted}{text}
3
\end{minted}



Sweeping \texttt{R} like this turns {\textquotedbl}is my stochastic pipeline unbiased?{\textquotedbl} from a statistical question into an exhaustive one — the pattern the shipped test suite uses.



\subsection{An accelerator-style activation, κ-measured: hard-tanh vs Tanh}



\label{7978642076077323668}{}


Hardware activation units often ship piecewise approximations. The κ machinery makes the substitution honest: measure the approximation{\textquotesingle}s worst code-point deviation from the defined activation, exhaustively, and register it under that measured bound. Hard-tanh — \texttt{clamp(x, −1, 1)} — as a stand-in for \texttt{Tanh}:




\begin{minted}{julia}
one4 = Binary8p4se(1.0)
hardtanh(x) = Clamp(Binary8p4se, RNE_SatNone, x, Negate(one4), one4)

κ, exhaustive = measure_kappa(hardtanh, :Tanh, Binary8p4se, (Binary8p4se,), RNE_SatNone)
(κ, exhaustive)
\end{minted}




\begin{minted}{text}
(4.0, true)                # worst deviation: 4 code points, verified on all 256 inputs
\end{minted}



Where is it worst? At \texttt{x = 1.0}: hard-tanh returns 1.0 while the defined \texttt{tanh(1.0)} is 0.75 — four code points along the total order. The registration round-trip, including the conformance surface:




\begin{minted}{julia}
impl = register_approx!(:hardtanh_act, :Tanh, Binary8p4se, (Binary8p4se,),
                        RNE_SatNone, hardtanh; κ=4)
(kappa(:hardtanh_act), kappa_measured(impl), :hardtanh_act in list_approx())
\end{minted}




\begin{minted}{text}
(4.0, 4.0, true)           # declared = measured; conformance_report() now lists it
\end{minted}



(\texttt{unregister\_approx!(:hardtanh\_act)} removes it.) Declaring \texttt{κ=3} would be \emph{rejected} — understatement is impossible by construction. The decision of whether a κ = 4 activation is acceptable now belongs to your accuracy budget, not to hope: combined with the margin rule from the General AI section, any two pre-activations whose defined \texttt{Tanh} images sit more than 8 code points apart keep their order under hard-tanh, provably.



\subsection{Benchmarking without measuring the dispatcher}



\label{585007679975857547}{}


The package{\textquotesingle}s benchmark doctrine, in one snippet (needs the \texttt{benchmarking/} environment for Chairmarks). Format types enter as \textbf{type parameters}; operands come from Chairmarks{\textquotesingle} \emph{untimed} \texttt{setup}; and if you retrieve functions reflectively (\texttt{getfield(SmallFloats, op)}), pass them through an argument barrier so they specialize:




\begin{minted}{julia}
using Chairmarks, SmallFloats, Random
using Statistics: median

function bench_add(::Type{T}) where {T<:Binary}          # T: type parameter, not a global
    pool = [rawvalue(T, rand(UInt8)) for _ in 1:4096]
    @be (rand(pool), rand(pool)) (t -> Add(T, RNE_SatNone, t[1], t[2]))(_)
end

b = bench_add(Binary8p4se)
(round(median(b).time * 1e9; digits=1), median(b).allocs)
\end{minted}




\begin{minted}{text}
(16.3, 0.0)          # ns per full scalar Add, zero allocations
\end{minted}



\begin{tcolorbox}[toptitle=-1mm,bottomtitle=1mm,colback=admonition-warning!50!white,colframe=admonition-warning,title=\textbf{The 60× trap}]
The same call with \texttt{T} read from a non-\texttt{const} global measures {\textasciitilde}1 µs — Julia{\textquotesingle}s dynamic keyword dispatch, not this package. Two project post-mortems trace to exactly this mistake; the shipped \texttt{benchmarking/benchmarking.jl} asserts specialization (zero warm-path allocation) before it believes any number, and so should yours.

\end{tcolorbox}


\part{Benchmarks}


\chapter{SmallFloats.jl benchmark report}



\label{17395114360235324059}{}


Generated: 2026-07-24 22:01 UTC  ·  Julia 1.12.6  ·  alderlake (32 logical CPUs, 1 Julia thread)  ·  Float128 paths: enabled  ·  Chairmarks 1.3.1



Reference format for per-operation tables: \texttt{Binary8p4se} under \texttt{(NearestTiesToEven, SatNone)}. Every table names its operand class: the scalar-operation tables appear in four variants — all code points (NaN and ±Inf sampled), NaN excluded, finite-only, and per-operation in-domain — and every other sampled table uses the all-code-points pool, identified in its note. Times are per call; medians with minima alongside. Methodology per the recorded benchmark doctrine: type-parameterized barriers, untimed setup, specialization preflight.



\section{Core primitives}



\label{10766655646123147915}{}


The decode/compare/step/classify layer plus the projection engine. Operands: all code points — NaN and ±Inf sampled.




\begin{table}[h]
\centering
\begin{tabulary}{\linewidth}{R R R R}
\toprule
operation & median & min & allocs \\
\toprule
\texttt{decode} & 1.4 ns & 1.4 ns & 0 \\
\texttt{order\_key} & 1.6 ns & 1.4 ns & 0 \\
\texttt{x < y} & 1.6 ns & 1.6 ns & 0 \\
\texttt{TotalOrder} & 1.8 ns & 1.6 ns & 0 \\
\texttt{Class} & 1.8 ns & 1.2 ns & 0 \\
\texttt{NextGreaterThan} & 1.8 ns & 1.2 ns & 0 \\
\texttt{project (RNE·SatNone)} & 6.5 ns & 1.8 ns & 0 \\
\texttt{project (StochasticA[8], R drawn)} & 7.0 ns & 1.9 ns & 0 \\
\bottomrule
\end{tabulary}

\end{table}



\section{Scalar operations}



\label{11509067801894334317}{}


Each arity is measured under four operand classes — safe args, no NaN/Inf, no NaN, and all code points — so operand-class effects (NaN fast rows, ±Inf special rows) are separable.




\clearpage



\subsection{Unary (30)}



\label{2812026322430809105}{}


\subsubsection{Safe args}



\label{15024080724600605730}{}


Finite operands within each operation{\textquotesingle}s safe domain — the fully unmasked per-operation scalar cost. The argument-restricted ops (Sqrt, RSqrt, Log, Log2, LogOnePlus, Recip, Divide, ArcSin, ArcCos, ArcCosh, ArcTanh) draw from explicit per-argument safe-domain predicates; every other op uses finite operand tuples whose defined result is not NaN (oracle-derived). Sorted by median.




\begin{table}[h]
\centering
\begin{tabulary}{\linewidth}{R R R R}
\toprule
operation & median & min & allocs \\
\toprule
\texttt{Abs} & 6.7 ns & 4.3 ns & 0 \\
\texttt{Negate} & 6.8 ns & 4.7 ns & 0 \\
\texttt{TanPi} & 12.4 ns & 5.5 ns & 0 \\
\texttt{Recip} & 18.5 ns & 7.9 ns & 0 \\
\texttt{Sqrt} & 18.5 ns & 5.0 ns & 0 \\
\texttt{RSqrt} & 20.6 ns & 9.4 ns & 0 \\
\texttt{Cosh} & 20.9 ns & 7.2 ns & 0 \\
\texttt{Cos} & 21.6 ns & 7.3 ns & 0 \\
\texttt{Sin} & 21.7 ns & 5.2 ns & 0 \\
\texttt{Exp2} & 22.2 ns & 7.2 ns & 0 \\
\texttt{Exp} & 22.4 ns & 7.2 ns & 0 \\
\texttt{ExpMinusOne} & 23.0 ns & 5.3 ns & 0 \\
\texttt{Tanh} & 23.1 ns & 5.5 ns & 0 \\
\texttt{ArcSin} & 23.2 ns & 4.8 ns & 0 \\
\texttt{ArcCos} & 24.3 ns & 5.7 ns & 0 \\
\texttt{ArcSinPi} & 25.4 ns & 5.2 ns & 0 \\
\texttt{Sinh} & 25.6 ns & 6.2 ns & 0 \\
\texttt{Log} & 25.9 ns & 5.8 ns & 0 \\
\texttt{Tan} & 26.7 ns & 5.5 ns & 0 \\
\texttt{ArcCosPi} & 26.8 ns & 5.9 ns & 0 \\
\texttt{Log2} & 26.9 ns & 5.8 ns & 0 \\
\texttt{LogOnePlus} & 27.1 ns & 5.2 ns & 0 \\
\texttt{SinPi} & 28.0 ns & 6.1 ns & 0 \\
\texttt{CosPi} & 28.0 ns & 6.9 ns & 0 \\
\texttt{ArcTan} & 28.2 ns & 5.0 ns & 0 \\
\texttt{ArcTanh} & 30.8 ns & 6.1 ns & 0 \\
\texttt{ArcTanPi} & 30.9 ns & 5.0 ns & 0 \\
\texttt{Softplus} & 33.3 ns & 14.5 ns & 0 \\
\texttt{ArcCosh} & 33.4 ns & 5.2 ns & 0 \\
\texttt{ArcSinh} & 33.5 ns & 5.3 ns & 0 \\
\bottomrule
\end{tabulary}

\end{table}



\subsubsection{No NaN, Inf args}



\label{16036960291186107490}{}


Operands exclude NaN and ±Inf; finite datums only (zeros and subnormals kept). Domain-restricted ops still take NaN fast rows on out-of-domain finite operands. Sorted by median.




\begin{table}[h]
\centering
\begin{tabulary}{\linewidth}{R R R R}
\toprule
operation & median & min & allocs \\
\toprule
\texttt{Sqrt} & 1.7 ns & 1.6 ns & 0 \\
\texttt{ArcCosh} & 1.8 ns & 1.8 ns & 0 \\
\texttt{ArcCos} & 1.8 ns & 1.8 ns & 0 \\
\texttt{ArcSin} & 4.9 ns & 1.8 ns & 0 \\
\texttt{ArcTanh} & 6.0 ns & 1.8 ns & 0 \\
\texttt{ArcCosPi} & 6.1 ns & 1.6 ns & 0 \\
\texttt{Abs} & 6.7 ns & 4.4 ns & 0 \\
\texttt{Negate} & 6.8 ns & 4.5 ns & 0 \\
\texttt{Log2} & 8.5 ns & 1.9 ns & 0 \\
\texttt{RSqrt} & 9.4 ns & 1.6 ns & 0 \\
\texttt{ArcSinPi} & 16.9 ns & 1.8 ns & 0 \\
\texttt{Recip} & 17.9 ns & 1.6 ns & 0 \\
\texttt{Cos} & 21.4 ns & 7.0 ns & 0 \\
\texttt{Sin} & 21.5 ns & 4.6 ns & 0 \\
\texttt{Exp2} & 22.1 ns & 7.0 ns & 0 \\
\texttt{Exp} & 22.3 ns & 7.0 ns & 0 \\
\texttt{ExpMinusOne} & 23.0 ns & 4.6 ns & 0 \\
\texttt{Tanh} & 23.2 ns & 5.0 ns & 0 \\
\texttt{Cosh} & 25.1 ns & 7.0 ns & 0 \\
\texttt{Log} & 25.4 ns & 1.8 ns & 0 \\
\texttt{Sinh} & 25.7 ns & 4.7 ns & 0 \\
\texttt{LogOnePlus} & 26.2 ns & 1.4 ns & 0 \\
\texttt{Tan} & 27.1 ns & 5.5 ns & 0 \\
\texttt{CosPi} & 27.8 ns & 6.5 ns & 0 \\
\texttt{SinPi} & 27.8 ns & 5.7 ns & 0 \\
\texttt{ArcTan} & 28.3 ns & 5.7 ns & 0 \\
\texttt{TanPi} & 30.8 ns & 5.4 ns & 0 \\
\texttt{ArcTanPi} & 30.9 ns & 5.0 ns & 0 \\
\texttt{Softplus} & 33.4 ns & 14.2 ns & 0 \\
\texttt{ArcSinh} & 33.6 ns & 4.7 ns & 0 \\
\bottomrule
\end{tabulary}

\end{table}



\subsubsection{No NaN args}



\label{11020842090045916616}{}


Operands exclude the NaN code point; ±Inf and every finite datum are sampled. Sorted by median.




\begin{table}[h]
\centering
\begin{tabulary}{\linewidth}{R R R R}
\toprule
operation & median & min & allocs \\
\toprule
\texttt{RSqrt} & 1.6 ns & 1.6 ns & 0 \\
\texttt{ArcCosh} & 1.8 ns & 1.8 ns & 0 \\
\texttt{Log} & 1.8 ns & 1.6 ns & 0 \\
\texttt{Log2} & 2.0 ns & 1.8 ns & 0 \\
\texttt{Sqrt} & 5.9 ns & 1.6 ns & 0 \\
\texttt{ArcTanh} & 6.5 ns & 1.8 ns & 0 \\
\texttt{Abs} & 6.7 ns & 4.8 ns & 0 \\
\texttt{Negate} & 6.8 ns & 4.5 ns & 0 \\
\texttt{ArcSinPi} & 7.2 ns & 1.8 ns & 0 \\
\texttt{ArcCosPi} & 7.4 ns & 1.6 ns & 0 \\
\texttt{ArcCos} & 16.4 ns & 1.8 ns & 0 \\
\texttt{Recip} & 17.8 ns & 1.6 ns & 0 \\
\texttt{ArcSin} & 19.6 ns & 1.8 ns & 0 \\
\texttt{Cosh} & 21.4 ns & 4.9 ns & 0 \\
\texttt{Cos} & 21.4 ns & 1.8 ns & 0 \\
\texttt{Sin} & 21.6 ns & 1.8 ns & 0 \\
\texttt{Exp2} & 22.1 ns & 5.0 ns & 0 \\
\texttt{Exp} & 22.3 ns & 5.0 ns & 0 \\
\texttt{ExpMinusOne} & 23.1 ns & 4.5 ns & 0 \\
\texttt{Tanh} & 23.2 ns & 4.9 ns & 0 \\
\texttt{Sinh} & 25.4 ns & 4.6 ns & 0 \\
\texttt{LogOnePlus} & 26.2 ns & 1.4 ns & 0 \\
\texttt{Tan} & 26.6 ns & 1.8 ns & 0 \\
\texttt{CosPi} & 27.7 ns & 1.9 ns & 0 \\
\texttt{SinPi} & 27.8 ns & 1.9 ns & 0 \\
\texttt{ArcTan} & 28.3 ns & 5.4 ns & 0 \\
\texttt{TanPi} & 30.8 ns & 1.8 ns & 0 \\
\texttt{ArcTanPi} & 30.9 ns & 4.8 ns & 0 \\
\texttt{Softplus} & 33.4 ns & 5.4 ns & 0 \\
\texttt{ArcSinh} & 33.5 ns & 4.6 ns & 0 \\
\bottomrule
\end{tabulary}

\end{table}



\subsubsection{All code points}



\label{8134844262983830456}{}


Operands drawn uniformly over ALL code points — NaN and ±Inf are sampled; medians of domain-restricted ops are diluted by instant NaN rows. Sorted by median. Transcendental rows mix special-row fast returns with enclosure-path evaluations, so these are \emph{scalar-path} costs; bulk unary work routes through 256-byte tables (see Array kernels).




\begin{table}[h]
\centering
\begin{tabulary}{\linewidth}{R R R R}
\toprule
operation & median & min & allocs \\
\toprule
\texttt{RSqrt} & 1.6 ns & 1.6 ns & 0 \\
\texttt{ArcCosh} & 1.8 ns & 1.8 ns & 0 \\
\texttt{Log} & 1.8 ns & 1.6 ns & 0 \\
\texttt{Log2} & 2.0 ns & 1.8 ns & 0 \\
\texttt{ArcSin} & 2.1 ns & 1.6 ns & 0 \\
\texttt{Sqrt} & 5.3 ns & 1.6 ns & 0 \\
\texttt{ArcTanh} & 6.6 ns & 1.8 ns & 0 \\
\texttt{Abs} & 6.7 ns & 1.9 ns & 0 \\
\texttt{Negate} & 6.8 ns & 1.6 ns & 0 \\
\texttt{ArcCosPi} & 7.1 ns & 1.6 ns & 0 \\
\texttt{ArcSinPi} & 7.4 ns & 1.6 ns & 0 \\
\texttt{Recip} & 17.9 ns & 1.6 ns & 0 \\
\texttt{Cosh} & 21.1 ns & 1.6 ns & 0 \\
\texttt{Cos} & 21.4 ns & 1.8 ns & 0 \\
\texttt{Sin} & 21.4 ns & 1.6 ns & 0 \\
\texttt{Exp2} & 22.1 ns & 1.8 ns & 0 \\
\texttt{Exp} & 22.3 ns & 2.0 ns & 0 \\
\texttt{ExpMinusOne} & 23.0 ns & 1.6 ns & 0 \\
\texttt{Tanh} & 23.2 ns & 1.8 ns & 0 \\
\texttt{ArcCos} & 23.9 ns & 1.8 ns & 0 \\
\texttt{Sinh} & 25.4 ns & 1.8 ns & 0 \\
\texttt{LogOnePlus} & 26.2 ns & 1.4 ns & 0 \\
\texttt{Tan} & 26.5 ns & 1.8 ns & 0 \\
\texttt{CosPi} & 27.7 ns & 1.8 ns & 0 \\
\texttt{SinPi} & 27.8 ns & 1.8 ns & 0 \\
\texttt{ArcTan} & 28.3 ns & 1.9 ns & 0 \\
\texttt{ArcTanPi} & 30.8 ns & 1.8 ns & 0 \\
\texttt{TanPi} & 30.8 ns & 1.6 ns & 0 \\
\texttt{Softplus} & 33.4 ns & 1.9 ns & 0 \\
\texttt{ArcSinh} & 33.6 ns & 1.8 ns & 0 \\
\bottomrule
\end{tabulary}

\end{table}




\clearpage



\subsection{Binary (18)}



\label{9493958918153706001}{}


\subsubsection{Safe args}



\label{1931488528371625443}{}


Finite operands within each operation{\textquotesingle}s safe domain — the fully unmasked per-operation scalar cost. The argument-restricted ops (Sqrt, RSqrt, Log, Log2, LogOnePlus, Recip, Divide, ArcSin, ArcCos, ArcCosh, ArcTanh) draw from explicit per-argument safe-domain predicates; every other op uses finite operand tuples whose defined result is not NaN (oracle-derived). Sorted by median.




\begin{table}[h]
\centering
\begin{tabulary}{\linewidth}{R R R R}
\toprule
operation & median & min & allocs \\
\toprule
\texttt{MinimumMagnitude} & 7.0 ns & 4.8 ns & 0 \\
\texttt{MaximumMagnitude} & 7.1 ns & 6.9 ns & 0 \\
\texttt{CopySign} & 7.1 ns & 5.0 ns & 0 \\
\texttt{Minimum} & 7.2 ns & 5.0 ns & 0 \\
\texttt{Maximum} & 7.2 ns & 4.6 ns & 0 \\
\texttt{MinimumFinite} & 7.2 ns & 4.8 ns & 0 \\
\texttt{MaximumFinite} & 7.2 ns & 4.5 ns & 0 \\
\texttt{MinimumNumber} & 7.2 ns & 4.7 ns & 0 \\
\texttt{MinimumMagnitudeNumber} & 7.3 ns & 5.6 ns & 0 \\
\texttt{MaximumMagnitudeNumber} & 7.3 ns & 7.1 ns & 0 \\
\texttt{MaximumNumber} & 7.3 ns & 4.8 ns & 0 \\
\texttt{Multiply} & 7.4 ns & 5.1 ns & 0 \\
\texttt{Add} & 7.9 ns & 5.5 ns & 0 \\
\texttt{Subtract} & 9.3 ns & 6.7 ns & 0 \\
\texttt{Divide} & 18.2 ns & 5.7 ns & 0 \\
\texttt{Hypot} & 27.9 ns & 7.6 ns & 0 \\
\texttt{ArcTan2} & 32.6 ns & 8.2 ns & 0 \\
\texttt{ArcTan2Pi} & 35.2 ns & 6.9 ns & 0 \\
\bottomrule
\end{tabulary}

\end{table}



\subsubsection{No NaN, Inf args}



\label{9842352165407780921}{}


Operands exclude NaN and ±Inf; finite datums only (zeros and subnormals kept). Domain-restricted ops still take NaN fast rows on out-of-domain finite operands. Sorted by median.




\begin{table}[h]
\centering
\begin{tabulary}{\linewidth}{R R R R}
\toprule
operation & median & min & allocs \\
\toprule
\texttt{MinimumMagnitude} & 7.1 ns & 4.9 ns & 0 \\
\texttt{CopySign} & 7.1 ns & 5.0 ns & 0 \\
\texttt{Minimum} & 7.2 ns & 4.9 ns & 0 \\
\texttt{Maximum} & 7.2 ns & 5.0 ns & 0 \\
\texttt{MaximumMagnitude} & 7.2 ns & 6.9 ns & 0 \\
\texttt{MaximumFinite} & 7.2 ns & 4.9 ns & 0 \\
\texttt{MinimumFinite} & 7.2 ns & 4.6 ns & 0 \\
\texttt{MaximumMagnitudeNumber} & 7.3 ns & 7.1 ns & 0 \\
\texttt{MinimumMagnitudeNumber} & 7.3 ns & 5.0 ns & 0 \\
\texttt{MaximumNumber} & 7.3 ns & 4.8 ns & 0 \\
\texttt{MinimumNumber} & 7.3 ns & 4.8 ns & 0 \\
\texttt{Multiply} & 7.4 ns & 5.5 ns & 0 \\
\texttt{Add} & 7.9 ns & 5.7 ns & 0 \\
\texttt{Subtract} & 9.3 ns & 7.0 ns & 0 \\
\texttt{Divide} & 18.2 ns & 2.5 ns & 0 \\
\texttt{Hypot} & 27.8 ns & 7.6 ns & 0 \\
\texttt{ArcTan2} & 32.6 ns & 7.7 ns & 0 \\
\texttt{ArcTan2Pi} & 35.1 ns & 6.6 ns & 0 \\
\bottomrule
\end{tabulary}

\end{table}



\subsubsection{No NaN args}



\label{5394694367249928840}{}


Operands exclude the NaN code point; ±Inf and every finite datum are sampled. Sorted by median.




\begin{table}[h]
\centering
\begin{tabulary}{\linewidth}{R R R R}
\toprule
operation & median & min & allocs \\
\toprule
\texttt{MinimumMagnitude} & 7.1 ns & 4.8 ns & 0 \\
\texttt{CopySign} & 7.1 ns & 5.2 ns & 0 \\
\texttt{Maximum} & 7.2 ns & 4.6 ns & 0 \\
\texttt{Minimum} & 7.2 ns & 4.8 ns & 0 \\
\texttt{MaximumMagnitude} & 7.2 ns & 5.1 ns & 0 \\
\texttt{MinimumFinite} & 7.3 ns & 4.5 ns & 0 \\
\texttt{MaximumMagnitudeNumber} & 7.3 ns & 5.1 ns & 0 \\
\texttt{MinimumMagnitudeNumber} & 7.3 ns & 5.0 ns & 0 \\
\texttt{MaximumFinite} & 7.3 ns & 5.1 ns & 0 \\
\texttt{Multiply} & 7.3 ns & 2.1 ns & 0 \\
\texttt{MinimumNumber} & 7.4 ns & 4.9 ns & 0 \\
\texttt{MaximumNumber} & 7.4 ns & 4.8 ns & 0 \\
\texttt{Add} & 7.9 ns & 2.7 ns & 0 \\
\texttt{Subtract} & 9.3 ns & 6.4 ns & 0 \\
\texttt{Divide} & 18.2 ns & 2.5 ns & 0 \\
\texttt{Hypot} & 27.8 ns & 5.0 ns & 0 \\
\texttt{ArcTan2} & 32.6 ns & 7.1 ns & 0 \\
\texttt{ArcTan2Pi} & 35.1 ns & 6.8 ns & 0 \\
\bottomrule
\end{tabulary}

\end{table}



\subsubsection{All code points}



\label{11499910499181378231}{}


Operands drawn uniformly over ALL code points — NaN and ±Inf are sampled; medians of domain-restricted ops are diluted by instant NaN rows. Sorted by median.




\begin{table}[h]
\centering
\begin{tabulary}{\linewidth}{R R R R}
\toprule
operation & median & min & allocs \\
\toprule
\texttt{MinimumMagnitude} & 7.1 ns & 1.8 ns & 0 \\
\texttt{Minimum} & 7.2 ns & 1.8 ns & 0 \\
\texttt{Maximum} & 7.2 ns & 1.6 ns & 0 \\
\texttt{MaximumMagnitude} & 7.2 ns & 1.8 ns & 0 \\
\texttt{MinimumFinite} & 7.2 ns & 4.6 ns & 0 \\
\texttt{MaximumFinite} & 7.2 ns & 4.7 ns & 0 \\
\texttt{CopySign} & 7.2 ns & 1.6 ns & 0 \\
\texttt{MinimumMagnitudeNumber} & 7.3 ns & 5.3 ns & 0 \\
\texttt{MaximumMagnitudeNumber} & 7.3 ns & 5.3 ns & 0 \\
\texttt{MaximumNumber} & 7.3 ns & 4.8 ns & 0 \\
\texttt{MinimumNumber} & 7.3 ns & 5.0 ns & 0 \\
\texttt{Multiply} & 7.4 ns & 1.8 ns & 0 \\
\texttt{Add} & 7.9 ns & 2.1 ns & 0 \\
\texttt{Subtract} & 9.3 ns & 3.7 ns & 0 \\
\texttt{Divide} & 18.3 ns & 2.5 ns & 0 \\
\texttt{Hypot} & 27.8 ns & 2.3 ns & 0 \\
\texttt{ArcTan2} & 32.5 ns & 3.5 ns & 0 \\
\texttt{ArcTan2Pi} & 35.0 ns & 3.7 ns & 0 \\
\bottomrule
\end{tabulary}

\end{table}



\subsection{Ternary (3)}



\label{1836852574804533888}{}


\subsubsection{Safe args}



\label{9973129846987395370}{}


Finite operands within each operation{\textquotesingle}s safe domain — the fully unmasked per-operation scalar cost. The argument-restricted ops (Sqrt, RSqrt, Log, Log2, LogOnePlus, Recip, Divide, ArcSin, ArcCos, ArcCosh, ArcTanh) draw from explicit per-argument safe-domain predicates; every other op uses finite operand tuples whose defined result is not NaN (oracle-derived). Sorted by median.




\begin{table}[h]
\centering
\begin{tabulary}{\linewidth}{R R R R}
\toprule
operation & median & min & allocs \\
\toprule
\texttt{Clamp} & 8.0 ns & 4.8 ns & 0 \\
\texttt{FAA} & 9.8 ns & 7.4 ns & 0 \\
\texttt{FMA} & 9.9 ns & 8.8 ns & 0 \\
\bottomrule
\end{tabulary}

\end{table}



\subsubsection{No NaN, Inf args}



\label{15876135725619624151}{}


Operands exclude NaN and ±Inf; finite datums only (zeros and subnormals kept). Domain-restricted ops still take NaN fast rows on out-of-domain finite operands. Sorted by median.




\begin{table}[h]
\centering
\begin{tabulary}{\linewidth}{R R R R}
\toprule
operation & median & min & allocs \\
\toprule
\texttt{Clamp} & 8.0 ns & 4.8 ns & 0 \\
\texttt{FAA} & 9.8 ns & 7.6 ns & 0 \\
\texttt{FMA} & 10.0 ns & 8.0 ns & 0 \\
\bottomrule
\end{tabulary}

\end{table}



\subsubsection{No NaN args}



\label{14279209625200149733}{}


Operands exclude the NaN code point; ±Inf and every finite datum are sampled. Sorted by median.




\begin{table}[h]
\centering
\begin{tabulary}{\linewidth}{R R R R}
\toprule
operation & median & min & allocs \\
\toprule
\texttt{Clamp} & 8.0 ns & 4.8 ns & 0 \\
\texttt{FAA} & 9.8 ns & 7.2 ns & 0 \\
\texttt{FMA} & 10.0 ns & 7.4 ns & 0 \\
\bottomrule
\end{tabulary}

\end{table}



\subsubsection{All code points}



\label{689699942421014554}{}


Operands drawn uniformly over ALL code points — NaN and ±Inf are sampled; medians of domain-restricted ops are diluted by instant NaN rows. Sorted by median.




\begin{table}[h]
\centering
\begin{tabulary}{\linewidth}{R R R R}
\toprule
operation & median & min & allocs \\
\toprule
\texttt{Clamp} & 8.0 ns & 1.6 ns & 0 \\
\texttt{FAA} & 9.8 ns & 3.7 ns & 0 \\
\texttt{FMA} & 10.0 ns & 4.0 ns & 0 \\
\bottomrule
\end{tabulary}

\end{table}



\section{Sensitivity studies}



\label{14065249436593140616}{}


How the scalar costs move with the format and with the projection specification.



\subsection{Format sensitivity}



\label{18234318524667261557}{}


Same three binary ops across formats; \texttt{Binary8p1uf} exercises the wide-exponent-spread escalations, small-K formats the tiny-table regime. Operands: all code points — NaN and ±Inf sampled.




\begin{table}[h]
\centering
\begin{tabulary}{\linewidth}{R R R R}
\toprule
operation & median & min & allocs \\
\toprule
\texttt{Add⟨Binary8p4se⟩} & 7.9 ns & 2.1 ns & 0 \\
\texttt{Divide⟨Binary8p4se⟩} & 18.3 ns & 2.5 ns & 0 \\
\texttt{Multiply⟨Binary8p4se⟩} & 7.4 ns & 1.8 ns & 0 \\
\texttt{Add⟨Binary8p3sf⟩} & 7.4 ns & 2.1 ns & 0 \\
\texttt{Divide⟨Binary8p3sf⟩} & 16.0 ns & 2.1 ns & 0 \\
\texttt{Multiply⟨Binary8p3sf⟩} & 6.8 ns & 1.6 ns & 0 \\
\texttt{Add⟨Binary8p1uf⟩} & 72.8 ns & 7.4 ns & 0 \\
\texttt{Divide⟨Binary8p1uf⟩} & 8.3 ns & 2.5 ns & 0 \\
\texttt{Multiply⟨Binary8p1uf⟩} & 6.8 ns & 1.4 ns & 0 \\
\texttt{Add⟨Binary5p2se⟩} & 7.9 ns & 2.3 ns & 0 \\
\texttt{Divide⟨Binary5p2se⟩} & 9.0 ns & 2.5 ns & 0 \\
\texttt{Multiply⟨Binary5p2se⟩} & 7.3 ns & 1.8 ns & 0 \\
\texttt{Add⟨Binary3p1se⟩} & 6.7 ns & 2.3 ns & 0 \\
\texttt{Divide⟨Binary3p1se⟩} & 6.9 ns & 2.5 ns & 0 \\
\texttt{Multiply⟨Binary3p1se⟩} & 5.7 ns & 1.6 ns & 0 \\
\bottomrule
\end{tabulary}

\end{table}



\subsection{Projection by rounding/saturation mode}



\label{15662701981835679547}{}


\texttt{project(Binary8p4se, ρ, x)} over the mode vocabulary (stochastic budgets N = 8). Operands: all code points — NaN and ±Inf sampled.




\begin{table}[h]
\centering
\begin{tabulary}{\linewidth}{R R R R}
\toprule
operation & median & min & allocs \\
\toprule
\texttt{NearestTiesToEven} & 6.5 ns & 1.6 ns & 0 \\
\texttt{NearestTiesToAway} & 7.7 ns & 2.3 ns & 0 \\
\texttt{TowardPositive} & 6.0 ns & 2.1 ns & 0 \\
\texttt{TowardNegative} & 6.0 ns & 2.5 ns & 0 \\
\texttt{TowardZero} & 5.3 ns & 2.6 ns & 0 \\
\texttt{ToOdd} & 6.6 ns & 1.6 ns & 0 \\
\texttt{StochasticA[8]} & 7.1 ns & 1.8 ns & 0 \\
\texttt{StochasticB[8]} & 7.1 ns & 1.6 ns & 0 \\
\texttt{StochasticC[8]} & 7.1 ns & 1.6 ns & 0 \\
\texttt{RNE · SatFinite} & 7.0 ns & 1.6 ns & 0 \\
\texttt{RNE · SatPropagate} & 6.7 ns & 1.9 ns & 0 \\
\bottomrule
\end{tabulary}

\end{table}



\section{Kernels and storage}



\label{9956393270762268485}{}


Array-shaped work: the table-gather and compute kernels, the ternary bitwidth policy, sorting, table builds, blocks, and packed/converted storage.



\subsection{Array kernels (vmap)}



\label{14819160431636425252}{}


Warm caches: table specializations prebuilt, so table rows measure the gather; scalar-loop rows measure the full compute pipeline per element. The ternary row here is \texttt{Binary8p4se} (K=8, always the compute path); see the next section for how the ternary bitwidth policy behaves across K. Operands: all code points — NaN and ±Inf sampled.




\begin{table}[h]
\centering
\begin{tabulary}{\linewidth}{R R R R R}
\toprule
operation & median & min & allocs & per element \\
\toprule
\texttt{vmap unary (table gather), n=65536} & 8.71 μs & 8.59 μs & 0 & 0.13 ns/elem — 7.52 Gelem/s \\
\texttt{vmap binary (table gather), n=65536} & 16.78 μs & 16.48 μs & 0 & 0.26 ns/elem — 3.91 Gelem/s \\
\texttt{vmap ternary (scalar loop), n=65536} & 1.01 ms & 1.0 ms & 0 & 15.49 ns/elem — 0.06 Gelem/s \\
\texttt{vmap binary stochastic (scalar loop), n=65536} & 646.35 μs & 633.53 μs & 0 & 9.86 ns/elem — 0.1 Gelem/s \\
\texttt{vmap unary through PackedVector, n=65536} & 83.59 μs & 80.0 μs & 517 & 1.28 ns/elem — 0.78 Gelem/s \\
\bottomrule
\end{tabulary}

\end{table}



\subsection{Ternary bitwidth tiers (FMA/FAA)}



\label{14443973703327008650}{}


\texttt{FMA}/\texttt{FAA}/\texttt{Clamp} are total functions on \texttt{2{\textasciicircum}(K1+K2+K3)} code points, but that count spans 512 B (K=3) to 16 MiB (K=8), so the array kernel tables small operand formats eagerly, tables mid-size ones adaptively (after enough elements amortize the build; not shown here — see the adaptive-cache gate in \texttt{test/ternary\_opt.jl}), and always runs the scalar compute kernel at K=8, threaded above a size cutoff when \texttt{Threads.nthreads() > 1}. Each tier{\textquotesingle}s optimized row is paired with a scalar-loop baseline (policy Refs forced off around the measurement, restored after) so the win is visible per tier; this process has 1 Julia thread. No threaded/sequential comparison below — rerun with \texttt{julia -t N} (N > 1) to see it. Same reference format, ρ, and operand pool discipline as Array kernels above.




\begin{table}[h]
\centering
\begin{tabulary}{\linewidth}{R R R R R}
\toprule
operation & median & min & allocs & per element \\
\toprule
\texttt{FMA K=4 (eager table), n=65536} & 20.47 μs & 19.51 μs & 0 & 0.31 ns/elem — 3.2 Gelem/s \\
\texttt{FMA K=4 (eager table), scalar-loop baseline, n=65536} & 967.64 μs & 958.72 μs & 0 & 14.76 ns/elem — 0.07 Gelem/s \\
\texttt{FMA K=6 (eager table), n=65536} & 24.75 μs & 24.35 μs & 0 & 0.38 ns/elem — 2.65 Gelem/s \\
\texttt{FMA K=6 (eager table), scalar-loop baseline, n=65536} & 1.02 ms & 997.1 μs & 0 & 15.56 ns/elem — 0.06 Gelem/s \\
\texttt{FMA K=8 (compute), n=65536} & 1.0 ms & 986.96 μs & 0 & 15.28 ns/elem — 0.07 Gelem/s \\
\texttt{FMA K=8 (compute), scalar-loop baseline, n=65536} & 994.01 μs & 982.32 μs & 0 & 15.17 ns/elem — 0.07 Gelem/s \\
\bottomrule
\end{tabulary}

\end{table}



\subsection{Sorting (64 K values)}



\label{10624969220324633770}{}


Counting sort is installed as the default algorithm for \texttt{Binary} vectors. Operands: all code points — NaN and ±Inf sampled.




\begin{table}[h]
\centering
\begin{tabulary}{\linewidth}{R R R R}
\toprule
operation & median & min & allocs \\
\toprule
\texttt{sort! (counting sort via defalg), n=65536} & 156.82 μs & 154.99 μs & 2 \\
\texttt{sort! (stock comparison sort), n=65536} & 1.39 ms & 1.35 ms & 3 \\
\texttt{sort! rev=true (counting sort), n=65536} & 160.84 μs & 158.14 μs & 2 \\
\bottomrule
\end{tabulary}

\end{table}



\subsection{Table builds (oracle + projection, Float128-first)}



\label{17249747280115456051}{}


Cold cache per sample (\texttt{empty\_tables!} in untimed setup); JIT pre-warmed. The warm-hit column is the steady-state cost of \texttt{get\_table} when the specialization is already cached (median / min). Table entries enumerate every code point by construction (NaN and ±Inf included).




\begin{table}[h]
\centering
\begin{tabulary}{\linewidth}{R R R R R}
\toprule
operation & median & min & allocs & warm hit \\
\toprule
\texttt{Exp⟨8p4se⟩ (256 entries)} & 29.94 μs & 29.06 μs & 257 & 64.2 ns / 62.4 ns \\
\texttt{Tanh⟨8p4se⟩ (256 entries)} & 29.85 μs & 29.1 μs & 257 & 64.3 ns / 62.6 ns \\
\texttt{Add⟨8p4se×8p4se⟩ (64 K entries)} & 25.49 ms & 23.51 ms & 785668 & 64.4 ns / 62.7 ns \\
\texttt{Divide⟨8p4se×8p4se⟩ (64 K)} & 25.24 ms & 23.2 ms & 784906 & 63.3 ns / 61.7 ns \\
\texttt{Add⟨8p1uf×8p1uf⟩ (64 K, wide-spread)} & 42.65 ms & 42.06 ms & 1587030 & 64.3 ns / 62.7 ns \\
\bottomrule
\end{tabulary}

\end{table}



\subsection{Block and scaled operations}



\label{17033350743052450193}{}


Elements \texttt{Binary8p4se}, scales \texttt{Binary8p1uf}, B = 32. Operands: all code points — NaN and ±Inf sampled.




\begin{table}[h]
\centering
\begin{tabulary}{\linewidth}{R R R R R}
\toprule
operation & median & min & allocs & per lane \\
\toprule
\texttt{BlockAdd (B=32)} & 907.5 ns & 873.7 ns & 35 & 28.36 ns/lane \\
\texttt{BlockDotProduct → 8p4se (B=32)} & 294.0 ns & 220.0 ns & 3 & 9.19 ns/lane \\
\texttt{BlockReduceAdd → 8p4se (B=32)} & 457.8 ns & 21.8 ns & 0 & 14.3 ns/lane \\
\texttt{ConvertToBlockMaxAbsFinite (B=32)} & 1.84 μs & 1.75 μs & 75 & 57.55 ns/lane \\
\bottomrule
\end{tabulary}

\end{table}



\subsection{Conversions and packed storage}



\label{4461369721169500046}{}


Operands: all code points — NaN and ±Inf sampled.




\begin{table}[h]
\centering
\begin{tabulary}{\linewidth}{R R R R R}
\toprule
operation & median & min & allocs & per element \\
\toprule
\texttt{T(::UInt8) code-point constructor} & 1.2 ns & 1.2 ns & 0 &  \\
\texttt{rawvalue (unchecked kernel route)} & 1.2 ns & 1.2 ns & 0 &  \\
\texttt{T(::Float64) numeric constructor (projects)} & 7.6 ns & 2.8 ns & 0 &  \\
\texttt{Convert 8p4se → 8p3se (scalar)} & 6.6 ns & 1.8 ns & 0 &  \\
\texttt{Float64 → 8p4se (project)} & 6.5 ns & 1.8 ns & 0 &  \\
\texttt{PackedVector pack, n=65536} & 33.21 μs & 32.42 μs & 3 & 0.51 ns/elem \\
\texttt{PackedVector unpack (collect), n=65536} & 28.99 μs & 28.87 μs & 3 & 0.44 ns/elem \\
\bottomrule
\end{tabulary}

\end{table}



{\rule{\textwidth}{1pt}}


\emph{All numbers from this machine/run; absolute values vary by host. Regenerate with \texttt{julia --project=benchmark benchmark/benchmarking.jl}.}



\part{Adding Operations}


\chapter{Adding Operations}



\label{5861982930757625708}{}


How to add a new operation to SmallFloats.jl so that it inherits everything the existing catalog has: bit-exact defined results, the full projection-mode vocabulary, generated scalar/array/block surfaces, table caching, exhaustive test coverage, benchmark pickup, and a truthful conformance declaration. Read the \hyperlinkref{4463396622184713734}{Technical Guide} first — this page assumes its vocabulary (the result-kind protocol, the two rigor classes, the interval protocol, the \texttt{yd} envelope stage).



The architecture is registry-driven: for the supported arities, \textbf{most of the surface is generated}. Adding an operation is mostly (1) one registry line and (2) one rigorous \texttt{ωeval} method — plus the \emph{duties} that keep the package{\textquotesingle}s guarantees true. Every subsection below states the duties explicitly; they are not optional.



\section{Universal duties (every new operation)}



\label{15376602000632183864}{}


\begin{itemize}
\item[1. ] \textbf{Special rows are spec, not optimization.} Write out the NaN/±Inf/zero rows explicitly in \texttt{ωeval}, in draft style, before any general evaluation. Mark any behavior the draft text under-determines with \texttt{[interp]}.


\item[2. ] \textbf{Result-kind protocol.} \texttt{ωeval} must return one of \texttt{Float64 | Float128 | StickyF | BigExactF | EncloseF | Enclose128F}, with \texttt{Float64}/\texttt{Float128} reserved for \emph{provably exact} results and \texttt{StickyF} for exact-head-plus-tail-direction results whose soundness bound is discharged at the emitting site (see the \texttt{FMA}/\texttt{FAA} wide-spread escalations). Keep the return union narrow (ideally \texttt{\{Float64, EncloseF\}}) — a wide union defeats inference and costs allocations (measured: a 4-type union turned a 50 ns op into 240 ns with 2 allocations).


\item[3. ] \textbf{Termination (the Niven duty).} \texttt{project\_interval} terminates only if the enclosure is not chasing a value the projection grid can actually hit. Any input whose \emph{exact} result is a representable dyadic must be \textbf{peeled or detected} before an enclosure is built. For transcendental ops this is usually a finite set of special rows (prove completeness, as the π-family does with Niven{\textquotesingle}s theorem); for algebraic ops it is an exactness test.


\item[4. ] \textbf{Envelope justification.} If you supply an eager \texttt{yd} estimate, you owe an error bound: the true value must lie within \texttt{|yd|·2⁻⁴⁵} (\texttt{\_F64\_RELEXP}). Acceptable bases, strongest first: IEEE-CR hardware ops (÷, sqrt: ≤ ½ ulp, unconditional); short compositions of CR ops (≤ {\textasciitilde}1.5 ulp); faithful libm (≤ 1 ulp, an empirical-but-published claim); short faithful compositions with condition number ≤ 1 (≤ {\textasciitilde}3 ulp). \textbf{Measure it}: sample your estimator against ≥ 256-bit MPFR truth over the op{\textquotesingle}s domain, including the treacherous regions (near zeros of the result, domain boundaries), and demand ≥ 2⁶ slack under 2⁻⁴⁵. The same duty applies to a Float128 \texttt{fq} estimator against its 2⁻⁹⁰ envelope. Beware \textbf{argument-error amplification}: an estimate computed as \texttt{f(g(x))} where \texttt{g} rounds (e.g. \texttt{sin(π·r)} in Float64) can be catastrophically worse near zeros of \texttt{f} than a native result-domain evaluation (\texttt{sinpi(r)}) — measure the composition you actually ship.


\item[5. ] \textbf{Ladder validity.} The MPFR closure must return a genuine directed enclosure. Whole-expression \texttt{setrounding} is valid only when the composition is monotone in every intermediate (state the argument in a comment, as \texttt{Softplus} does); otherwise round each step in the correct direction explicitly (as \texttt{\_mpfr\_rsqrt} does for the anti-monotone \texttt{1/√x}).


\item[6. ] \textbf{Gates.} Run the full suite. The unary exhaustive ×hp gate and the coverage check pick up registry ops automatically; extend the hand-curated lists where applicable (monotonicity ops, binary exhaustive sweeps). Then regenerate the benchmark report — \texttt{bench\_scalar\_ops} iterates the registry lists, so your op appears in all four operand-class tables automatically; \textbf{add a \texttt{\_SAFE\_DOMAINS} entry} in \texttt{benchmarking/benchmarking.jl} if the op is argument-restricted, or the safe-args table will use the oracle-derived fallback pool.

\end{itemize}


\section{Scalar operations}



\label{11509067801894334317}{}


The registry lives in \texttt{src/ops\_scalar.jl}. Adding a name to \texttt{\_UNARY\_OPS} / \texttt{\_BINARY\_OPS} / \texttt{\_TERNARY\_OPS} causes, with no further work:



\begin{itemize}
\item registration in \texttt{OP\_REGISTRY} (group \texttt{:B}/\texttt{:C}/\texttt{:A} per the loop defaults — override with an explicit \texttt{register\_op!} if the default group is wrong);


\item the generated spec-register method \texttt{Op(fr, ρ, operands…; rng, R)} and the same-format convenience \texttt{Op(x…)};


\item generated array methods routing through \texttt{vmap} (Shape-A tables for pure ρ — always at arity ≤ 2, bitwidth-gated at arity 3 — Shape-B scalar/threaded loops otherwise);


\item generated \texttt{BlockOp} and \texttt{ScaledOp} variants (\texttt{src/blocks.jl});


\item the export (\texttt{src/SmallFloats.jl} loops the registry);


\item table cacheability, conformance listing, and test-harness coverage.

\end{itemize}


What is \emph{not} generated is the mathematics: \texttt{ωeval} in \texttt{src/oracle.jl}.



\subsection{Adding a unary operation}



\label{8519510493258928519}{}


Worked example: \textbf{\texttt{Cbrt}} (cube root), chosen because it exercises the hardest duty — exactness detection for an algebraic op.



\textbf{Step 1 — registry} (\texttt{src/ops\_scalar.jl}): append \texttt{:Cbrt} to \texttt{\_UNARY\_OPS}. The loop registers it as group \texttt{:B} with arity 1.



\textbf{Step 2 — ω-semantics} (\texttt{src/oracle.jl}), with every duty discharged:




\begin{minted}{julia}
# Exact-cube detection (termination duty): x = ±m·2^(3k) with m an odd perfect
# cube ⇒ cbrt(x) is a Float64-exact dyadic that project_interval must never
# chase. Oracle-path cost is irrelevant; correctness of the peel is everything.
function _exact_cbrt(x::Float64)
    fr, e = frexp(abs(x))               # |x| = fr·2^e, fr ∈ [0.5, 1)
    m = Int64(ldexp(fr, 53)); e -= 53   # |x| = m·2^e exactly
    tz = trailing_zeros(m); m >>= tz; e += tz
    e % 3 == 0 || return nothing
    r = round(Int64, cbrt(Float64(m)))
    for c in max(r - 1, 1):(r + 1)      # cbrt is faithful ⇒ true root ∈ r ± 1
        Int128(c)^3 == Int128(m) && return flipsign(ldexp(Float64(c), e ÷ 3), x)
    end
    nothing
end

function ωeval(::Val{:Cbrt}, x::Float64)
    isnan(x) && return NaN                      # special rows first: they are spec
    isinf(x) && return x                        # cbrt(±∞) = ±∞
    iszero(x) && return 0.0
    c = _exact_cbrt(x)
    c === nothing || return c                   # exact dyadic result — peeled
    # yd: Float64 cbrt is faithful (≤ 1 ulp ≈ 2⁻⁵²; measure it) ⇒ ≥ 2⁷ slack
    # under the 2⁻⁴⁵ envelope. fq: libquadmath cbrtq under the 2⁻⁹⁰ envelope.
    # Ladder: cbrt is monotone ⇒ whole-expression directed rounding encloses.
    _encl1(cbrt, x; fq=() -> cbrt(Float128(x)), yd=cbrt(x))
end
\end{minted}



\textbf{Step 3 — optional surface}: a Base veneer (\texttt{Base.cbrt(x::Binary) = Cbrt(x)} next to the others in \texttt{ops\_scalar.jl}), a docstring, and — since cbrt is monotone — its symbol in the test suite{\textquotesingle}s monotonicity list.



\textbf{Step 4 — validate}: envelope measurement for \texttt{cbrt} (include tiny and huge \texttt{x}; cbrt has no dangerous regions — the result is never near zero except at zero, which is peeled), then the full suite: the exhaustive ×hp gate now checks \texttt{Cbrt} on every code point of the reference format under four modes against the 3072-bit protocol, automatically.



\subsection{Adding a binary operation}



\label{7660721820420346737}{}


Worked example: \textbf{\texttt{LogAddExp}} — \texttt{log(eˣ + eʸ)} — chosen because it exercises composed estimators and the ladder-monotonicity argument.



\textbf{Step 1 — registry}: append \texttt{:LogAddExp} to \texttt{\_BINARY\_OPS} (default group \texttt{:C}; the special five arithmetic names get \texttt{:A}).



\textbf{Step 2 — ω-semantics}. One generic evaluator serves all three precisions — Float64 (\texttt{yd}), Float128 (\texttt{fq}), and BigFloat (the ladder):




\begin{minted}{julia}
# Stable form: h + log1p(exp(l − h)), h = max, l = min. Every step is monotone
# nondecreasing in (x, y), so whole-expression directed rounding in _mpfr2 is a
# valid enclosure (the same argument Softplus records).
_logaddexp(a, b) = (h = max(a, b); l = min(a, b); h + log1p(exp(l - h)))

function ωeval(::Val{:LogAddExp}, x::Float64, y::Float64)
    (isnan(x) | isnan(y)) && return NaN
    (x == Inf || y == Inf) && return Inf        # ∞ dominates
    x == -Inf && return y                       # log(0 + eʸ) = y, exactly
    y == -Inf && return x
    # yd: two faithful libm calls plus additions, condition ≤ 1 throughout
    # (result ≥ max(x, y)) ⇒ ≤ ~3 ulp ≈ 2⁻⁵⁰ — 2⁵ slack; measure it, including
    # x ≈ y (where log1p(exp(0)) dominates) and widely separated operands.
    _encl2(_logaddexp, x, y;
           fq=() -> _logaddexp(Float128(x), Float128(y)),
           yd=_logaddexp(x, y))
end
\end{minted}



\textbf{Termination note}: can \texttt{logaddexp(x, y)} be exactly dyadic for dyadic operands with the special rows peeled? \texttt{log(eˣ + eʸ) = d} requires \texttt{eˣ + eʸ = e{\textasciicircum}d} — by Lindemann–Weierstrass this forces relations that dyadic arguments cannot satisfy except on the peeled rows, so no further peel is needed; \textbf{record that argument in a comment}. When you cannot prove such a statement for your op, hunt for exact cases the way \texttt{\_exact\_cbrt} does.



\textbf{Step 3 — tests}: binary ops are not in the unary exhaustive gate; add a divide-style exhaustive sweep (256×256 operand pairs against a ladder-only reference, a handful of modes) to the suite, and an envelope-sanity loop entry if you introduced a new libm dependence.



\subsection{Adding a ternary operation}



\label{10687193682775408327}{}


Worked example: \textbf{\texttt{FMS}} — fused multiply-subtract, \texttt{x·y − z} — chosen because the cheapest correct implementation is \emph{delegation}, and delegation is a first-class technique here.



\textbf{Step 1 — registry}: append \texttt{:FMS} to \texttt{\_TERNARY\_OPS} (group \texttt{:A} by the loop). The generated array method routes through the same ternary bitwidth policy every other ternary op uses (\texttt{tables.jl}{\textquotesingle}s \texttt{\_ternary\_table\_for}): small operand formats table (eagerly or adaptively), \texttt{K = 8} runs the — optionally threaded — Shape-B scalar loop, and stochastic ρ always draws per element. No op-specific work needed; the policy keys on formats and ρ, not on which op it is.



\textbf{Step 2 — ω-semantics} by delegation to \texttt{FMA}{\textquotesingle}s already-verified analysis (width thresholds, \texttt{\_twosum} exactness, wide-spread fallback):




\begin{minted}{julia}
# x·y − z ≡ FMA(x, y, −z). Negation of z follows the Subtract convention:
# preserve NaN, and map −0 to +0 so the single-zero datum model is respected.
ωeval(::Val{:FMS}, x::Float64, y::Float64, z::Float64) =
    ωeval(Val(:FMA), x, y, isnan(z) ? z : (iszero(z) ? 0.0 : -z))
\end{minted}



Delegation inherits FMA{\textquotesingle}s result kinds, its exactness thresholds, \emph{and} its test pedigree — but verify the sign algebra of your reduction on the special rows (here: \texttt{FMS(∞, 1, ∞)} must be NaN, which the delegation produces because \texttt{FMA(∞, 1, −∞)} is \texttt{∞ − ∞}). Add the delegated op to the ternary exhaustive sweep in the suite.



\subsection{Adding a quaternary operation}



\label{1028380857650662652}{}


Arity 4 is \textbf{not yet plumbed} — the generation loops, kernels, and block layer handle arities 1–3. Adding a quaternary op is therefore two tasks: extend the plumbing (once), then add the op. Worked example: \textbf{\texttt{FMMA}} — \texttt{x·y + z·w}, a two-term dot product.



\textbf{Plumbing (one-time), four sites:}



\begin{itemize}
\item[1. ] \texttt{src/ops\_scalar.jl} — register manually after the arity loops (\texttt{register\_op!(:FMMA, 4, :A)}), and add an \texttt{op.arity == 4} branch to the spec-register generation loop, mirroring the ternary branch with a fourth operand:


\begin{minted}{julia}
else # arity 4
    @eval begin
        @inline function $name(fr::Type{<:Binary}, ρ::ProjSpec,
                               x::Binary, y::Binary, z::Binary, w::Binary;
                               rng::MaybeRNG=nothing, R::Union{Nothing,Int}=nothing)
            apply_op($V(), fr, ρ, _drawR(ρ, rng, R),
                     decode(x), decode(y), decode(z), decode(w))
        end
        @inline $name(x::T, y::T, z::T, w::T; kw...) where {T<:Binary} =
            $name(T, default_projspec(T), x, y, z, w; kw...)
    end
end
\end{minted}

(\texttt{apply\_op} and \texttt{ωeval} are already variadic — no change needed there.)


\item[2. ] \texttt{src/kernels.jl} — a four-array \texttt{vmap!} method. The plain Shape-B loop (copy the \emph{body} of the ternary method{\textquotesingle}s compute branch, add operand \texttt{D}) is enough to get correct results; the ternary method{\textquotesingle}s table-tiering and threading are a bitwidth-driven optimization layered on top (\texttt{tables.jl}{\textquotesingle}s \texttt{\_ternary\_table\_for} and its LRU cache), not required plumbing — extend to arity 4 only if a table win at that arity is worth the 2{\textasciicircum}(4K) growth. Also add an arity-4 branch in the registry-generated array-surface loop and the matching rng-threading overload (mirroring the three-array one) so stochastic ρ dispatches correctly.


\item[3. ] \texttt{src/blocks.jl} — either an arity-4 branch in the Block/Scaled generation loop (four operand blocks, four \texttt{blockdecode}s) or an explicit skip (\texttt{op.arity > 3 \&\& continue}) if a block form is not wanted; be deliberate.


\item[4. ] \texttt{src/approx.jl} — \texttt{conformance\_report} prints arities with \texttt{for a in 1:3}; widen to \texttt{1:4}.

\end{itemize}


\textbf{The op itself} follows the \texttt{FAA}/\texttt{FMA} width-analysis playbook exactly:




\begin{minted}{julia}
const _DE_FMMA = 92     # two exact ≤17-bit products: 18 + ΔE ≤ 113, margin 3

function ωeval(::Val{:FMMA}, x::Float64, y::Float64, z::Float64, w::Float64)
    (isnan(x) | isnan(y) | isnan(z) | isnan(w)) && return NaN
    (((iszero(x) && isinf(y)) || (isinf(x) && iszero(y))) ||
     ((iszero(z) && isinf(w)) || (isinf(z) && iszero(w)))) && return NaN
    p = x * y; q = z * w                          # exact: ≤8-bit significands
    if isinf(p) || isinf(q)
        (isinf(p) && isinf(q) && p != q) && return NaN
        return isinf(p) ? p : q
    end
    s, e = _twosum(p, q)
    e == 0.0 && return iszero(s) ? 0.0 : s        # exact in Float64
    (_f128() && !iszero(p) && !iszero(q) && _expdiff(p, q) <= _DE_FMMA) &&
        return Float128(p) + Float128(q)          # exact by width
    BigExactF(() -> setprecision(() ->
        BigFloat(x) * BigFloat(y) + BigFloat(z) * BigFloat(w), BigFloat, _BIGP))
end
\end{minted}



\textbf{Testing duty is heavier at arity 4}: the full cross-product is 2³² tuples, so the suite entry must be \emph{sampled} (seeded, ladder-referenced, ≥ 10⁵ tuples across several modes) plus hand-picked exhaustion of the special-row algebra and the width-threshold boundary (\texttt{ΔE ∈ \{91, 92, 93\}} constructions). Document the sampling in the test, since it breaks the suite{\textquotesingle}s {\textquotedbl}enumerate, never sample{\textquotedbl} norm — that exception must be visible, not silent.



\section{Block operations}



\label{6513894782918662055}{}


The block layer (\texttt{src/blocks.jl}) implements the draft{\textquotesingle}s elementwise schema \textbf{once} — \texttt{blockdecode → ω-op lanewise → blockproject} — and generates every registry op{\textquotesingle}s \texttt{BlockOp}/\texttt{ScaledOp} from it. The reductions are hand-written on a second shared pattern. Which subsection you need depends on whether your operation fits one of those two patterns.



\subsection{Adding a scaled block operation}



\label{10512613773019865025}{}


If the scalar operation is (or can be) a \textbf{registry op, the scaled form is free}: registering \texttt{:Cbrt} above already produced \texttt{ScaledCbrt(fr, ρ, s, x)} and \texttt{BlockCbrt(fr, ρ, b, sr)} with the draft{\textquotesingle}s §5.4 semantics — decode \texttt{scale × element} exactly per lane, apply the ω-op, \texttt{BlockProject} against the result scale through the division cascade. \textbf{Prefer this route}; it is the package{\textquotesingle}s non-divergence mechanism.



Write a scaled op by hand only when it is \emph{not} an elementwise image of a registry op. The primitive to compose with is \texttt{\_bp\_element(fr, ρ, R, res, Sdat)} — it owns the draft{\textquotesingle}s scale-special rows (NaN/0/±Inf scale) and the entire division-by-scale rigor cascade. Example, a scaled affine \texttt{s₁·x + s₂·y} projected against a \emph{unit} result scale:




\begin{minted}{julia}
function ScaledAffine(fr::Type{<:Binary}, ρ::ProjSpec,
                      s1::Binary, x::Binary, s2::Binary, y::Binary;
                      rng::MaybeRNG=nothing)
    a = ωeval(Val(:Multiply), decode(s1), decode(x))::Float64   # exact lanes
    b = ωeval(Val(:Multiply), decode(s2), decode(y))::Float64
    _bp_element(fr, ρ, _drawR(ρ, rng, nothing), ωeval(Val(:Add), a, b), 1.0)
end
\end{minted}



The \texttt{::Float64} assertions are load-bearing: lane decodes are exact by width (≤ 17-bit significands), and the assertion turns that analysis into a runtime check. The \texttt{Add} result may be any kind in the protocol — \texttt{\_bp\_element} finishes all of them. Dividing by a scale of \texttt{1.0} is the identity fast path.



\subsection{Adding an elementwise block operation}



\label{15527859041760837550}{}


Same schema, whole-block granularity. The three obligations: decode operand blocks with \texttt{blockdecode} (exact), evaluate lanewise with \texttt{ωeval} (never with ad-hoc Float64 math), and finish with \texttt{blockproject} against the supplied result scale (never with per-lane \texttt{project} — the scale-division rows belong to \texttt{blockproject}). Example — \texttt{BlockRelu}, an elementwise \texttt{max(x, 0)} that reuses the verified \texttt{MaximumNumber} semantics:




\begin{minted}{julia}
"""BlockRelu(fr, ρ, b, sr): lanewise max(xᵢ, 0) — NaN lanes propagate NaN per
MaximumNumber's number-preferring rows only when both operands are NaN, i.e.
never here; state your lane semantics this explicitly for any new op."""
function BlockRelu(fr::Type{<:Binary}, ρ::ProjSpec, b::Block{B}, sr::Binary;
                   rng::MaybeRNG=nothing) where {B}
    X = blockdecode(b)
    Z = ntuple(i -> ωeval(Val(:MaximumNumber), X[i], 0.0), Val(B))
    blockproject(fr, ρ, sr, Z; rng)
end
\end{minted}



Add the name to the block-surface export list and to \texttt{conformance()}{\textquotesingle}s \texttt{blocknames} extension vector so the declaration stays truthful. Test against a from-scratch reference composition (the suite{\textquotesingle}s existing block tests show the pattern) — element format × scale format × B, exhaustively where feasible.



\subsection{Adding a reductive block operation}



\label{9686810033048460081}{}


Reductions do \textbf{not} fit the elementwise schema; they follow the second pattern, visible in \texttt{BlockReduceAdd}/\texttt{BlockDotProduct}: (1) resolve the ∞/NaN \textbf{fold algebra} on Float64 classifications first; (2) an integer \textbf{span filter} decides whether the whole reduction is exactly representable in \texttt{Float128} — if so, plain \texttt{Float128} accumulation \emph{is} the exact answer; (3) otherwise a \texttt{BigExactF} exact accumulator at provably sufficient precision; (4) exactly \textbf{one} projection, at the very end, via \texttt{\_finish}. Never round intermediates.



Worked example — \textbf{\texttt{BlockSumOfSquares}}, \texttt{Σ xᵢ²}:




\begin{minted}{julia}
"""BlockSumOfSquares(fr, ρ, b): project(Σ decode-laneᵢ²) — one rounding total.
Fold algebra: any NaN lane ⇒ NaN; else any ±Inf lane ⇒ +Inf (squares cannot
cancel — no opposite-infinity NaN case exists, unlike ReduceAdd); else exact."""
function BlockSumOfSquares(fr::Type{<:Binary}, ρ::ProjSpec, b::Block{B};
                           rng::MaybeRNG=nothing, R::Union{Nothing,Int}=nothing) where {B}
    X = blockdecode(b)
    res = if any(isnan, X)
        NaN
    elseif any(isinf, X)
        Inf
    elseif all(iszero, X)
        0.0
    else
        # span filter: lanes carry ≤17-bit significands ⇒ squares ≤34 bits and
        # a square's exponent is 2·(lane exponent); the sum is exact in Float128
        # when 34 + 2·span + ⌈log₂B⌉ + 1 ≤ 113  ⇒  2·span + ⌈log₂B⌉ ≤ 78.
        if _f128() && 2 * _expspan(X) + _log2ceil(B) <= 78
            acc = Float128(0)
            for v in X
                q = Float128(v)
                acc += q * q                      # exact squares, exact sum
            end
            acc
        else
            BigExactF(() -> setprecision(BigFloat, _REDPREC) do
                acc = BigFloat(0)
                for v in X
                    acc += BigFloat(v)^2          # exact at _REDPREC
                end
                acc
            end)
        end
    end
    _finish(fr, ρ, _drawR(ρ, rng, R), res)
end
\end{minted}



The two lines that demand your care in any new reduction are the \textbf{fold algebra} (derive it from the draft{\textquotesingle}s reduce definition on the extended reals — here the absence of an opposite-infinities NaN case is a \emph{theorem about squares}, stated in the docstring) and the \textbf{span-filter inequality} (derive it from operand widths, write the derivation in the comment, and add boundary-condition tests at the threshold, exactly as the suite does for the existing \texttt{\_DE\_*} constants). Everything else — the exact accumulator, the single \texttt{\_finish}, stochastic \texttt{R} plumbing — is pattern.



\section{After any addition: the closing checklist}



\label{16308398306918412518}{}


\begin{itemize}
\item Full test suite green, including your new sweep entries.


\item Envelope measurements recorded for any new \texttt{yd}/\texttt{fq} estimator.


\item Benchmark report regenerated; \texttt{\_SAFE\_DOMAINS} entry added if the op is argument-restricted; sanity-check the op{\textquotesingle}s row in the safe-args table.


\item Docstring on the public name; \texttt{[interp]} markers where you interpreted.


\item \texttt{conformance\_report()} shows the op (automatic for registry ops; manual block additions extend \texttt{blocknames}).

\end{itemize}


\part{External Reference}


\chapter{External Reference}



\label{5597696625397351643}{}


Docstrings for the documented \textbf{public} surface — the exported names. The full export list is far larger (400+ names, including the 120 format aliases, the predefined projection-spec grid, and the generated operation and block registers); names without docstrings are covered descriptively in the \hyperlinkref{10272851679537904033}{User Guide}; unexported machinery is in the \hyperlinkref{2804742620609579687}{Internal Reference}.


\hypertarget{14453628373362608485}{\texttt{SmallFloats.SmallFloats}}  -- {Module.}

\begin{adjustwidth}{2em}{0pt}


\begin{minted}{julia}
SmallFloats
\end{minted}

A conforming, performance-oriented Julia implementation of the IEEE P3109 draft standard for arithmetic formats for machine learning (bitwidths 3–8).

Bit-exact defined results on every default path; the projection engine (\texttt{RoundToPrecision → Saturate → Encode}) is the single write path into a code point; approximate fast paths exist only behind the explicit κ registry (\texttt{register\_approx!} / \texttt{approx}), never substituted silently. See \texttt{conformance()} for the live declaration and \texttt{SmallFloats.draft\_revision()} for the implemented draft.

\textbf{Performance note}

Pass format types through \texttt{const} bindings, type parameters, or function arguments. All exported entry points fully specialize when the format type is statically known (measured: a complete scalar \texttt{Add} ≈ 26 ns, \texttt{project} ≈ 13 ns, zero allocations). Calling through a \textbf{non-\texttt{const} global} format type forces dynamic dispatch on every call — measured at {\textasciitilde}1 µs per scalar keyword call — which is Julia semantics, not package overhead; a single function barrier (\texttt{f(::Type\{T\}, args...) where \{T\}}) restores full speed. The test suite pins the specialization properties (concrete return types, zero warm-path allocation) as deterministic regressions.



\end{adjustwidth}
\hypertarget{8725867475956632017}{\texttt{SmallFloats.RNA\_SatFinite}}  -- {Constant.}

\begin{adjustwidth}{2em}{0pt}

(NearestTiesToAway, SatFinite).



\end{adjustwidth}
\hypertarget{16426343133939897049}{\texttt{SmallFloats.RNA\_SatNone}}  -- {Constant.}

\begin{adjustwidth}{2em}{0pt}

(NearestTiesToAway, SatNone).



\end{adjustwidth}
\hypertarget{9948955282984860181}{\texttt{SmallFloats.RNA\_SatPropagate}}  -- {Constant.}

\begin{adjustwidth}{2em}{0pt}

(NearestTiesToAway, SatPropagate).



\end{adjustwidth}
\hypertarget{3359638186361138749}{\texttt{SmallFloats.RNE\_SatFinite}}  -- {Constant.}

\begin{adjustwidth}{2em}{0pt}

(NearestTiesToEven, SatFinite).



\end{adjustwidth}
\hypertarget{6254846941289803954}{\texttt{SmallFloats.RNE\_SatNone}}  -- {Constant.}

\begin{adjustwidth}{2em}{0pt}

(NearestTiesToEven, SatNone) — the package-wide default ρ.



\end{adjustwidth}
\hypertarget{15613002957007008711}{\texttt{SmallFloats.RNE\_SatPropagate}}  -- {Constant.}

\begin{adjustwidth}{2em}{0pt}

(NearestTiesToEven, SatPropagate).



\end{adjustwidth}
\hypertarget{6431776125964272887}{\texttt{SmallFloats.RTN\_SatFinite}}  -- {Constant.}

\begin{adjustwidth}{2em}{0pt}

(TowardNegative, SatFinite).



\end{adjustwidth}
\hypertarget{18308597707086069238}{\texttt{SmallFloats.RTN\_SatNone}}  -- {Constant.}

\begin{adjustwidth}{2em}{0pt}

(TowardNegative, SatNone).



\end{adjustwidth}
\hypertarget{15860643992407021701}{\texttt{SmallFloats.RTN\_SatPropagate}}  -- {Constant.}

\begin{adjustwidth}{2em}{0pt}

(TowardNegative, SatPropagate).



\end{adjustwidth}
\hypertarget{1023017805428823485}{\texttt{SmallFloats.RTO\_SatFinite}}  -- {Constant.}

\begin{adjustwidth}{2em}{0pt}

(ToOdd, SatFinite).



\end{adjustwidth}
\hypertarget{16442577991271067062}{\texttt{SmallFloats.RTO\_SatNone}}  -- {Constant.}

\begin{adjustwidth}{2em}{0pt}

(ToOdd, SatNone).



\end{adjustwidth}
\hypertarget{15307190884920783579}{\texttt{SmallFloats.RTO\_SatPropagate}}  -- {Constant.}

\begin{adjustwidth}{2em}{0pt}

(ToOdd, SatPropagate).



\end{adjustwidth}
\hypertarget{7512419803180778431}{\texttt{SmallFloats.RTP\_SatFinite}}  -- {Constant.}

\begin{adjustwidth}{2em}{0pt}

(TowardPositive, SatFinite).



\end{adjustwidth}
\hypertarget{11481809423421058656}{\texttt{SmallFloats.RTP\_SatNone}}  -- {Constant.}

\begin{adjustwidth}{2em}{0pt}

(TowardPositive, SatNone).



\end{adjustwidth}
\hypertarget{5630752953060625376}{\texttt{SmallFloats.RTP\_SatPropagate}}  -- {Constant.}

\begin{adjustwidth}{2em}{0pt}

(TowardPositive, SatPropagate).



\end{adjustwidth}
\hypertarget{8937139283316265274}{\texttt{SmallFloats.RTZ\_SatFinite}}  -- {Constant.}

\begin{adjustwidth}{2em}{0pt}

(TowardZero, SatFinite).



\end{adjustwidth}
\hypertarget{569534270503185777}{\texttt{SmallFloats.RTZ\_SatNone}}  -- {Constant.}

\begin{adjustwidth}{2em}{0pt}

(TowardZero, SatNone).



\end{adjustwidth}
\hypertarget{11314891540121714893}{\texttt{SmallFloats.RTZ\_SatPropagate}}  -- {Constant.}

\begin{adjustwidth}{2em}{0pt}

(TowardZero, SatPropagate).



\end{adjustwidth}
\hypertarget{9201328443408064130}{\texttt{SmallFloats.ApproxImpl}}  -- {Type.}

\begin{adjustwidth}{2em}{0pt}

A registered κ-approximate implementation of one operation specialization.



\end{adjustwidth}
\hypertarget{8158499324480878091}{\texttt{SmallFloats.Binary}}  -- {Type.}

\begin{adjustwidth}{2em}{0pt}


\begin{minted}{julia}
Binary{K,P,SGN,EXT} <: AbstractFloat
\end{minted}

A P3109 floating-point value: a code point in the format with bitwidth \texttt{K ∈ 3:8}, precision \texttt{P}, signedness \texttt{SGN::Bool} (\texttt{true} = Signed), domain \texttt{EXT::Bool} (\texttt{true} = Extended, i.e. the datum set includes infinities).

The code point occupies the low \texttt{K} bits of the payload byte; the high \texttt{8-K} bits are maintained zero as a representation invariant. Construct raw code points with \texttt{T(c::UInt8)} (validated code-point construction), \texttt{rawvalue(T, c)} (unchecked kernel route), or \texttt{Binary\{...\}(Val(:code), x::UInt8)}; construct from numeric values with \texttt{T(x::Real)} (default projection spec) or \texttt{Convert}. UInt8 is the one argument type meaning \emph{code point}; all other Reals mean \emph{value}.



\end{adjustwidth}
\hypertarget{16566511736104831917}{\texttt{SmallFloats.Block}}  -- {Type.}

\begin{adjustwidth}{2em}{0pt}


\begin{minted}{julia}
Block{B, FS<:Binary, FE<:Binary}
\end{minted}

A P3109 block \texttt{(s, [x₁ … x\_B])}: scale factor in format \texttt{FS}, \texttt{B ≥ 1} elements in format \texttt{FE} (draft §5). \texttt{isbits}; blocks live in registers/stack.



\end{adjustwidth}
\hypertarget{1380584316956099505}{\texttt{SmallFloats.BlockVector}}  -- {Type.}

\begin{adjustwidth}{2em}{0pt}


\begin{minted}{julia}
BlockVector{B,FS,FE} <: AbstractVector{Block{B,FS,FE}}
\end{minted}

Structure-of-arrays storage: \texttt{scales::Vector\{FS\}}, \texttt{elems::Matrix\{FE\}} (B × n, column-major so each block{\textquotesingle}s elements are one contiguous column).



\end{adjustwidth}
\hypertarget{3620020081592949828}{\texttt{SmallFloats.FPClass}}  -- {Type.}

\begin{adjustwidth}{2em}{0pt}

Eight-way datum classification returned by \texttt{Class} (draft §4.13.1).



\end{adjustwidth}
\hypertarget{10417286468662382093}{\texttt{SmallFloats.NearestTiesToAway}}  -- {Type.}

\begin{adjustwidth}{2em}{0pt}

Round to nearest; ties away from zero (draft §4.7.1).



\end{adjustwidth}
\hypertarget{1826807226245170700}{\texttt{SmallFloats.NearestTiesToEven}}  -- {Type.}

\begin{adjustwidth}{2em}{0pt}

Round to nearest; ties to the even code point (IEEE default; draft §4.7.1).



\end{adjustwidth}
\hypertarget{4445951075310959743}{\texttt{SmallFloats.PackedVector}}  -- {Type.}

\begin{adjustwidth}{2em}{0pt}


\begin{minted}{julia}
PackedVector{F<:Binary} <: AbstractVector{F}
\end{minted}

Bit-packed storage of \texttt{F} code points at \texttt{bitwidth(F)} bits per element. Construct with \texttt{PackedVector(A::AbstractVector\{F\})}; recover bytes with \texttt{Vector(pv)} or \texttt{collect(pv)}. Memory: ⌈n·K/64⌉ words vs n bytes unpacked.



\end{adjustwidth}
\hypertarget{3971740801855465015}{\texttt{SmallFloats.ProjSpec}}  -- {Type.}

\begin{adjustwidth}{2em}{0pt}


\begin{minted}{julia}
ProjSpec{R<:RoundingMode3109, S<:SaturationMode}
\end{minted}

A projection specification ρ = (rounding mode, saturation mode), draft §4.2. Zero-size; construct as \texttt{ProjSpec(NearestTiesToEven(), SatNone())} or via the exported constants. Kernels specialize on its type.



\end{adjustwidth}
\hypertarget{5642129125365414505}{\texttt{SmallFloats.StochasticA}}  -- {Type.}

\begin{adjustwidth}{2em}{0pt}

Stochastic rounding, variant A of draft §4.7.4: RoundAway ⟺ ⌊ν·2{\textasciicircum}N⌋ + R ≥ 2{\textasciicircum}N.



\end{adjustwidth}
\hypertarget{15839936377740889884}{\texttt{SmallFloats.StochasticB}}  -- {Type.}

\begin{adjustwidth}{2em}{0pt}

Stochastic rounding, variant B: RoundAway ⟺ ⌊ν·2{\textasciicircum}(N+1)⌋ + (2R+1) ≥ 2{\textasciicircum}(N+1).



\end{adjustwidth}
\hypertarget{11520338016165756811}{\texttt{SmallFloats.StochasticC}}  -- {Type.}

\begin{adjustwidth}{2em}{0pt}

Stochastic rounding, variant C: RoundAway ⟺ RNITE(ν·2{\textasciicircum}N) + R ≥ 2{\textasciicircum}N.



\end{adjustwidth}
\hypertarget{11439938397858131910}{\texttt{SmallFloats.ToOdd}}  -- {Type.}

\begin{adjustwidth}{2em}{0pt}

Round inexact results to the nearest \emph{odd} code point (draft §4.7.3; sticky-friendly).



\end{adjustwidth}
\hypertarget{2540988212624594754}{\texttt{SmallFloats.TowardNegative}}  -- {Type.}

\begin{adjustwidth}{2em}{0pt}

Directed rounding toward −∞ (draft §4.7.2).



\end{adjustwidth}
\hypertarget{5586350886118351895}{\texttt{SmallFloats.TowardPositive}}  -- {Type.}

\begin{adjustwidth}{2em}{0pt}

Directed rounding toward +∞ (draft §4.7.2).



\end{adjustwidth}
\hypertarget{16105017792344936885}{\texttt{SmallFloats.TowardZero}}  -- {Type.}

\begin{adjustwidth}{2em}{0pt}

Directed rounding toward zero — truncation (draft §4.7.2).



\end{adjustwidth}
\hypertarget{14266569397601053665}{\texttt{SmallFloats.BlockDotProduct}}  -- {Method.}

\begin{adjustwidth}{2em}{0pt}

BlockDotProduct (draft §5.3.2). Lane products can carry 64 significant bits, so the products and their sum are formed in the exact accumulator, with the ∞/NaN fold algebra resolved on the Float64 classifications first.



\end{adjustwidth}
\hypertarget{8157221386497919512}{\texttt{SmallFloats.BlockReduceAdd}}  -- {Method.}

\begin{adjustwidth}{2em}{0pt}

BlockReduceAdd (draft §5.3.1): project(reduce(ωAdd, [0, X…])).



\end{adjustwidth}
\hypertarget{3734371764054498420}{\texttt{SmallFloats.BlockReduceMultiply}}  -- {Method.}

\begin{adjustwidth}{2em}{0pt}

BlockReduceMultiply (draft §5.3.1): project(reduce(ωMultiply, [1, X…])).



\end{adjustwidth}
\hypertarget{11726654467661344946}{\texttt{SmallFloats.Class}}  -- {Method.}

\begin{adjustwidth}{2em}{0pt}

Class(v) -> FPClass (draft §4.13.1): classify \texttt{v} as NaN, ±Inf, ±normal, ±subnormal, or zero.



\end{adjustwidth}
\hypertarget{5403593118059835048}{\texttt{SmallFloats.Convert}}  -- {Method.}

\begin{adjustwidth}{2em}{0pt}


\begin{minted}{julia}
Convert(fr, ρ, A::AbstractArray{<:Union{Float16,Float32,Float64,BFloat16}}; rng) -> Array{fr}
\end{minted}

Bulk ingestion: exact widening per element, then projection under ρ. Unlike the Binary-array Convert (a Shape-A table gather), external inputs are not enumerable, so this is a Shape-B loop; stochastic ρ draws per element.



\end{adjustwidth}
\hypertarget{2402979922820933143}{\texttt{SmallFloats.Convert}}  -- {Method.}

\begin{adjustwidth}{2em}{0pt}


\begin{minted}{julia}
Convert(fr, ρ, x) -> fr
\end{minted}

Draft §4.9 Convert⟨f\emph{x, f}r, ρ⟩. Accepts \texttt{Binary}, IEEE binary16/32/64 floats (widened exactly to the Float64 carrier), \texttt{Float128} (projected directly), \texttt{Integer} (exact via a sufficiently wide BigFloat), and \texttt{BigFloat} (projected directly; the caller warrants the value is exact).



\end{adjustwidth}
\hypertarget{6528265372875072184}{\texttt{SmallFloats.ConvertFromBlock}}  -- {Method.}

\begin{adjustwidth}{2em}{0pt}

ConvertFromBlock (§5.2.1): blockdecode, then project each element (no scale division).



\end{adjustwidth}
\hypertarget{15271042205516165074}{\texttt{SmallFloats.ConvertToBlock}}  -- {Method.}

\begin{adjustwidth}{2em}{0pt}

ConvertToBlock (§5.2.2): decode elements, BlockProject against the supplied scale.



\end{adjustwidth}
\hypertarget{11716888295752982041}{\texttt{SmallFloats.ConvertToBlockMaxAbsFinite}}  -- {Method.}

\begin{adjustwidth}{2em}{0pt}

ConvertToBlockMaxAbsFinite (§5.2.3): S = reduce(ωMaximumFinite, [NaN, |X₁|…]), project S under ρs, then BlockProject the elements under ρ.



\end{adjustwidth}
\hypertarget{16754626316686161819}{\texttt{SmallFloats.DefaultAccumulatorType}}  -- {Method.}

\begin{adjustwidth}{2em}{0pt}


\begin{minted}{julia}
DefaultAccumulatorType() -> Type{<:AbstractFloat}
DefaultAccumulatorType!(T::Type{<:AbstractFloat}) -> T
\end{minted}

The session{\textquotesingle}s default accumulator carrier for wide-precision work (reductions, dot products). Initialized to \texttt{binary32} (the exported IEEE 754 alias for \texttt{Float32}). Any \texttt{AbstractFloat} type is accepted (\texttt{binary64}, \texttt{Float128}, \texttt{BigFloat}, …).



\end{adjustwidth}
\hypertarget{1645330464161432694}{\texttt{SmallFloats.DefaultProjection}}  -- {Method.}

\begin{adjustwidth}{2em}{0pt}


\begin{minted}{julia}
DefaultProjection() -> ProjSpec
DefaultProjection!(ρ::ProjSpec) -> ρ
DefaultProjection!(m, s) -> ProjSpec
\end{minted}

The session{\textquotesingle}s default projection specification. Initialized to \texttt{(DefaultRoundingMode(), DefaultSaturationMode())} = \texttt{RNE\_SatNone}. Setting it directly decomposes ρ into \hyperlinkref{790339220741924172}{\texttt{DefaultRoundingMode}} and \hyperlinkref{11118232298850481653}{\texttt{DefaultSaturationMode}}, so the three stay coherent in both directions.



\end{adjustwidth}
\hypertarget{1239855964180895668}{\texttt{SmallFloats.DefaultRNG}}  -- {Method.}

\begin{adjustwidth}{2em}{0pt}


\begin{minted}{julia}
DefaultRNG() -> Type{<:AbstractRNG} | AbstractRNG
DefaultRNG!(rng) -> rng
\end{minted}

The session{\textquotesingle}s default random-number generator for stochastic rounding. Initialized to the \texttt{Xoshiro} \emph{type} (a fresh generator per use); may be set to an \texttt{AbstractRNG} instance for a reproducible stream.



\end{adjustwidth}
\hypertarget{16660646822053217478}{\texttt{SmallFloats.DefaultRbits}}  -- {Method.}

\begin{adjustwidth}{2em}{0pt}


\begin{minted}{julia}
DefaultRbits() -> Int
DefaultRbits!(n::Int) -> n
\end{minted}

The session{\textquotesingle}s default random-bit budget N for the stochastic rounding families (\texttt{StochasticA\{N\}}, \texttt{StochasticB\{N\}}, \texttt{StochasticC\{N\}}). Initialized to 8; must satisfy 1 ≤ N ≤ 60.



\end{adjustwidth}
\hypertarget{10309025597907569675}{\texttt{SmallFloats.DefaultReturnType}}  -- {Method.}

\begin{adjustwidth}{2em}{0pt}


\begin{minted}{julia}
DefaultReturnType() -> Type{<:Binary}
DefaultReturnType!(T::Type{<:Binary}) -> T
\end{minted}

The session{\textquotesingle}s default \emph{result} format — the format an operation projects into when the caller does not name one. Initialized to \texttt{Binary8p2se}. Independent of \texttt{DefaultType}, which is the default \emph{operand} format.



\end{adjustwidth}
\hypertarget{790339220741924172}{\texttt{SmallFloats.DefaultRoundingMode}}  -- {Method.}

\begin{adjustwidth}{2em}{0pt}


\begin{minted}{julia}
DefaultRoundingMode() -> RoundingMode3109
DefaultRoundingMode!(m) -> m
\end{minted}

The session{\textquotesingle}s default rounding mode. Initialized to \texttt{NearestTiesToEven()}. Setting it rebuilds \hyperlinkref{1645330464161432694}{\texttt{DefaultProjection}} from the new mode and the current \hyperlinkref{11118232298850481653}{\texttt{DefaultSaturationMode}}. Accepts an instance or a fully-parameterized type (\texttt{NearestTiesToAway}, \texttt{StochasticA\{8\}}, …).



\end{adjustwidth}
\hypertarget{11118232298850481653}{\texttt{SmallFloats.DefaultSaturationMode}}  -- {Method.}

\begin{adjustwidth}{2em}{0pt}


\begin{minted}{julia}
DefaultSaturationMode() -> SaturationMode
DefaultSaturationMode!(s) -> s
\end{minted}

The session{\textquotesingle}s default saturation mode. Initialized to \texttt{SatNone()}. Setting it rebuilds \hyperlinkref{1645330464161432694}{\texttt{DefaultProjection}} from the current \hyperlinkref{790339220741924172}{\texttt{DefaultRoundingMode}} and the new mode. Accepts an instance or a type.



\end{adjustwidth}
\hypertarget{4425608605544763458}{\texttt{SmallFloats.NextGreaterThan}}  -- {Method.}

\begin{adjustwidth}{2em}{0pt}

NextGreaterThan(v) (draft §4.16): the least datum greater than \texttt{v} in the total order — one step up the code lattice, with NaN → NaN and MaxFinite/+Inf → NaN at the top. \texttt{Base.nextfloat} on \texttt{Binary} is this operation.



\end{adjustwidth}
\hypertarget{6098566503846399249}{\texttt{SmallFloats.NextLessThan}}  -- {Method.}

\begin{adjustwidth}{2em}{0pt}

NextLessThan(v) (draft §4.16): the greatest datum less than \texttt{v} in the total order — one step down the code lattice, with NaN → NaN and MinFinite/−Inf → NaN at the bottom. \texttt{Base.prevfloat} on \texttt{Binary} is this operation.



\end{adjustwidth}
\hypertarget{10219551362409877369}{\texttt{SmallFloats.RoundOf}}  -- {Function.}

\begin{adjustwidth}{2em}{0pt}

Draft-named accessor: RoundOf(ρ) — the rounding mode of ρ (§4.2).



\end{adjustwidth}
\hypertarget{15323456454116582477}{\texttt{SmallFloats.SatOf}}  -- {Function.}

\begin{adjustwidth}{2em}{0pt}

Draft-named accessor: SatOf(ρ) — the saturation mode of ρ (§4.2).



\end{adjustwidth}
\hypertarget{11021615464506174358}{\texttt{SmallFloats.TotalOrder}}  -- {Method.}

\begin{adjustwidth}{2em}{0pt}

TotalOrder⟨fx,fy⟩ (draft §4.12.1): x ≤ y in the total order (single NaN largest). Same-format comparisons run on the integer order key (bitops plan K1) — proven equivalent to the decode order exhaustively over all pairs; cross-format keys are not comparable, so mixed formats keep the decode path.



\end{adjustwidth}
\hypertarget{8320965106983893823}{\texttt{SmallFloats.approx}}  -- {Method.}

\begin{adjustwidth}{2em}{0pt}

Retrieve a registered approximate implementation (callable via \texttt{.fn}).



\end{adjustwidth}
\hypertarget{10167532859324986487}{\texttt{SmallFloats.bitwidth}}  -- {Method.}

\begin{adjustwidth}{2em}{0pt}

Format bitwidth K (3–8).



\end{adjustwidth}
\hypertarget{12693397031431816949}{\texttt{SmallFloats.codedistance}}  -- {Method.}

\begin{adjustwidth}{2em}{0pt}

Code-point distance along the total order between two values of one format.



\end{adjustwidth}
\hypertarget{3926398212792611207}{\texttt{SmallFloats.conformance}}  -- {Method.}

\begin{adjustwidth}{2em}{0pt}


\begin{minted}{julia}
conformance() -> ConformanceDeclaration
\end{minted}

The package{\textquotesingle}s draft-§4.6-style conformance declaration, derived live from the operation registry, the table cache (the specializations actually instantiated), and the approximate-implementation registry. Serialize with \texttt{conformance\_dict} or render with \texttt{conformance\_report}.



\end{adjustwidth}
\hypertarget{590311556486812861}{\texttt{SmallFloats.conformance\_dict}}  -- {Function.}

\begin{adjustwidth}{2em}{0pt}

Nested \texttt{Dict\{String,Any\}} form of the declaration — serializable by any JSON/TOML writer.



\end{adjustwidth}
\hypertarget{18172896664359196039}{\texttt{SmallFloats.conformance\_report}}  -- {Function.}

\begin{adjustwidth}{2em}{0pt}

Human-readable conformance report.



\end{adjustwidth}
\hypertarget{17071822498540248224}{\texttt{SmallFloats.decode!}}  -- {Method.}

\begin{adjustwidth}{2em}{0pt}


\begin{minted}{julia}
decode!(dest::AbstractArray{Float32}, A::AbstractArray{<:Binary}) -> dest
decode!(dest::AbstractArray{Float64}, A::AbstractArray{<:Binary}) -> dest
\end{minted}

Exact bulk ωDecode: a Shape-A gather of the (narrowed) decode table into an external float array. Both carriers are exact for every K ≤ 8 datum.



\end{adjustwidth}
\hypertarget{15004588707537249672}{\texttt{SmallFloats.decode}}  -- {Method.}

\begin{adjustwidth}{2em}{0pt}


\begin{minted}{julia}
decode(v::Binary) -> Float64
\end{minted}

ωDecode (draft §4.7.2). Float64 is the universal exact carrier for K ≤ 8 datums (≤ 8-bit significands, |exponent| ≲ 2{\textasciicircum}7); exactness is asserted by exhaustive test.

Implemented as a constant-tuple lookup (bitops plan K2): the per-format table is generated once from \texttt{\_decode\_compute} above, so the two are correct by construction and asserted equivalent exhaustively; constant inputs still fold.



\end{adjustwidth}
\hypertarget{14130930211625684600}{\texttt{SmallFloats.empty\_tables!}}  -- {Method.}

\begin{adjustwidth}{2em}{0pt}

Drop every cached table and adaptive counter (they rebuild lazily on next use).



\end{adjustwidth}
\hypertarget{9308606959501922909}{\texttt{SmallFloats.expbias}}  -- {Method.}

\begin{adjustwidth}{2em}{0pt}

Exponent bias: 2{\textasciicircum}(K−P−1) signed, 2{\textasciicircum}(K−P) unsigned.



\end{adjustwidth}
\hypertarget{13856278338529021160}{\texttt{SmallFloats.expbitwidth}}  -- {Method.}

\begin{adjustwidth}{2em}{0pt}

Width of the exponent field in bits: (K − signbit) − (P − 1).



\end{adjustwidth}
\hypertarget{17243710630823570067}{\texttt{SmallFloats.f32\_exact}}  -- {Method.}

\begin{adjustwidth}{2em}{0pt}


\begin{minted}{julia}
f32_exact(op, f1, f2) -> Bool
\end{minted}

True iff \texttt{op ∈ (:Add, :Subtract, :Multiply)} on Float32-decoded operands is exact for every finite pair of (f1, f2) datums — checked once by exhaustive enumeration against a 300-bit BigFloat oracle, then cached. When true, a Float32 intermediate is an exact carrier under every projection mode.



\end{adjustwidth}
\hypertarget{9290299928136503291}{\texttt{SmallFloats.f32\_impl}}  -- {Method.}

\begin{adjustwidth}{2em}{0pt}


\begin{minted}{julia}
f32_impl(op, fr, ρ) -> fn
\end{minted}

The decode32 → guarded-op32 → project kernel for \texttt{op⟨… → fr, ρ⟩}, shaped for \texttt{register\_approx!}. Nearest-ties ρ only: directed modes after a nearest-rounded Float32 compute measure κ ≥ 1 (and κ = NaN at saturation boundaries), and stochastic ρ is rejected by κ measurement itself.



\end{adjustwidth}
\hypertarget{2286900636783861378}{\texttt{SmallFloats.formatname}}  -- {Method.}

\begin{adjustwidth}{2em}{0pt}

\texttt{formatname(T)} — the draft §3.2 name of a format type.



\end{adjustwidth}
\hypertarget{17132358657286871632}{\texttt{SmallFloats.ftz\_variant}}  -- {Method.}

\begin{adjustwidth}{2em}{0pt}


\begin{minted}{julia}
ftz_variant(op, fr, f1, ρ) -> fn
\end{minted}

Build the Annex-style approximate unary implementation: compute the defined result, then flush subnormal results to zero or \texttt{MinNormalOf(fr)}, whichever is nearer, ties toward zero (magnitude-symmetric for signed formats). Returned \texttt{fn} is suitable for \texttt{register\_approx!}; for \texttt{fr} with precision P its κ is 2{\textasciicircum}(P-2) when P ≥ 2 (largest flushed-to-zero subnormal) and 0 when P = 1 (no subnormals exist).



\end{adjustwidth}
\hypertarget{4250219520413861193}{\texttt{SmallFloats.isextended}}  -- {Method.}

\begin{adjustwidth}{2em}{0pt}

Whether the format{\textquotesingle}s domain is Extended (datum set includes infinities).



\end{adjustwidth}
\hypertarget{11031203868920563628}{\texttt{SmallFloats.issigned}}  -- {Method.}

\begin{adjustwidth}{2em}{0pt}

Whether the format is Signed (has a sign bit and negative datums).



\end{adjustwidth}
\hypertarget{2984505828320663690}{\texttt{SmallFloats.list\_approx}}  -- {Method.}

\begin{adjustwidth}{2em}{0pt}

Sorted names of all registered κ-approximate implementations.



\end{adjustwidth}
\hypertarget{14020185040969508771}{\texttt{SmallFloats.measure\_kappa}}  -- {Method.}

\begin{adjustwidth}{2em}{0pt}


\begin{minted}{julia}
measure_kappa(fn, op, fr, argformats, ρ; max_exhaustive=2^22, rng, samples=2^20)
    -> (κ::Float64, exhaustive::Bool)
\end{minted}

Measure κ of \texttt{fn(x::argformats[1], …)::fr} against the defined results of draft operation \texttt{op} under ρ. Enumerates the full input cross-product when it has at most \texttt{max\_exhaustive} points (always true for arity ≤ 2); otherwise measures on \texttt{samples} uniform draws and reports \texttt{exhaustive = false}.



\end{adjustwidth}
\hypertarget{12405409248851917562}{\texttt{SmallFloats.nrandbits}}  -- {Method.}

\begin{adjustwidth}{2em}{0pt}

Number of random bits N consumed per projected element; 0 for pure modes.



\end{adjustwidth}
\hypertarget{5547067338070232361}{\texttt{SmallFloats.projmode}}  -- {Method.}

\begin{adjustwidth}{2em}{0pt}

Translate a \texttt{Base.RoundingMode} to its \texttt{RoundingMode3109} counterpart (identity on the latter). Used only at API boundaries. Internals always carry RoundingMode3109.



\end{adjustwidth}
\hypertarget{16755457755945711135}{\texttt{SmallFloats.rawvalue}}  -- {Method.}

\begin{adjustwidth}{2em}{0pt}

Unsafe raw constructor used by kernels after invariants are established.



\end{adjustwidth}
\hypertarget{3878213895374713375}{\texttt{SmallFloats.register\_approx!}}  -- {Method.}

\begin{adjustwidth}{2em}{0pt}


\begin{minted}{julia}
register_approx!(name, op, fr, argformats, ρ, fn; κ=nothing, kwargs...) -> ApproxImpl
\end{minted}

Register \texttt{fn} as a named κ-approximate implementation of \texttt{op⟨argformats → fr, ρ⟩}. κ is measured by enumeration (see \texttt{measure\_kappa}); a declared \texttt{κ} smaller than the measured value is \textbf{rejected}. Omitting \texttt{κ} declares the measured value. NaN-κ implementations (non-finite mismatches) may be registered only by declaring \texttt{κ=NaN}.



\end{adjustwidth}
\hypertarget{11323396924847339952}{\texttt{SmallFloats.register\_f32!}}  -- {Method.}

\begin{adjustwidth}{2em}{0pt}


\begin{minted}{julia}
register_f32!(op, fr, argformats, ρ; κ=nothing) -> ApproxImpl
\end{minted}

Build, measure, and register the Float32 kernel for \texttt{op⟨argformats → fr, ρ⟩} under the canonical name \texttt{f32/op/args→fr/rm\_sm}. Registration measures κ by exhaustive enumeration and fails loudly on understatement — that is the point. Measured baseline (2026-07-24): κ = 0 for all 120 same-format signatures × \{Add, Subtract, Multiply, Divide, Sqrt\} × RNE\_\{SatNone, SatFinite, SatPropagate\}.



\end{adjustwidth}
\hypertarget{10680873580113776002}{\texttt{SmallFloats.table\_bytes}}  -- {Method.}

\begin{adjustwidth}{2em}{0pt}

Total bytes currently held by the table caches (unary/binary + ternary).



\end{adjustwidth}
\hypertarget{13844074395020045017}{\texttt{SmallFloats.trailingsigbits}}  -- {Method.}

\begin{adjustwidth}{2em}{0pt}

Trailing-significand width P − 1 (the stored fraction bits).



\end{adjustwidth}
\hypertarget{12365808433851041199}{\texttt{SmallFloats.unregister\_approx!}}  -- {Method.}

\begin{adjustwidth}{2em}{0pt}

Remove a registered κ-approximate implementation (no-op if absent).



\end{adjustwidth}
\hypertarget{12243573033237441134}{\texttt{SmallFloats.vmap!}}  -- {Function.}

\begin{adjustwidth}{2em}{0pt}


\begin{minted}{julia}
vmap!(dest, Val(op), fr, ρ, A[, B[, C]]; rng) -> dest
vmap(op, fr, ρ, A...; rng)
\end{minted}

Elementwise draft operation over arrays of \texttt{Binary} values, projecting into \texttt{fr} under ρ. Pure ρ runs the Shape-A gather whenever a table exists or the ternary bitwidth policy grants one; stochastic specializations (and untabled ternary signatures) run the scalar path per element (each element drawing its own R).



\end{adjustwidth}
\hypertarget{14797466715659470848}{\texttt{SmallFloats.vmap}}  -- {Method.}

\begin{adjustwidth}{2em}{0pt}


\begin{minted}{julia}
vmap(op, fr, ρ, pv::PackedVector; rng=nothing) -> Vector{fr}
\end{minted}

Elementwise \texttt{op} over packed storage: unpack a \texttt{256}-element tile into a byte scratch, run the ordinary byte kernel (\texttt{vmap!}), emit, repeat. Never computes on packed words directly — the deliberate compute-unpacked/store-packed boundary.



\end{adjustwidth}
\hypertarget{10552114894679459838}{\texttt{SmallFloats.with\_default\_accumulatortype}}  -- {Function.}

\begin{adjustwidth}{2em}{0pt}


\begin{minted}{julia}
with_default_type(f, args...)          -> f(DefaultType(), args...)
with_default_returntype(f, args...)    -> f(DefaultReturnType(), args...)
with_default_accumulatortype(f, args...) -> f(DefaultAccumulatorType(), args...)
with_default_projection(f, args...)    -> f(DefaultProjection(), args...)
\end{minted}

Call \texttt{f} with the named session default as its first argument — the supported way to \emph{consume} a default. While the default still holds its initial value (the overwhelmingly common case) the call is statically compiled against that constant: no dynamic dispatch, \texttt{f} fully specialized. After the default is changed, the call crosses a function barrier instead: one dynamic dispatch, everything inside fully specialized.

Allocation contract: the combinator{\textquotesingle}s return type unions both paths. When \texttt{f}{\textquotesingle}s result type does not depend on the default — e.g. \texttt{with\_default\_projection} with the formats fixed by the caller — the union is concrete and the call is \textbf{zero-allocation with a concretely inferred result}. When the result{\textquotesingle}s type \emph{is} the default (\texttt{with\_default\_type((T, x) -> T(x), …)}), the value is computed on the specialized path but boxes once at escape — the irreducible cost of a runtime-chosen type.


\begin{minted}{jlcon}
julia> with_default_type((T, x) -> T(x), 1.5)
Binary8p2se(1.5 ≡ 0x41)

julia> with_default_projection((ρ, x, y) -> Add(Binary8p4se, ρ, x, y),
                               Binary8p4se(1.5), Binary8p4se(0.25))
Binary8p4se(1.75 ≡ 0x46)
\end{minted}



\end{adjustwidth}
\hypertarget{16240126272133078948}{\texttt{SmallFloats.with\_default\_projection}}  -- {Function.}

\begin{adjustwidth}{2em}{0pt}


\begin{minted}{julia}
with_default_type(f, args...)          -> f(DefaultType(), args...)
with_default_returntype(f, args...)    -> f(DefaultReturnType(), args...)
with_default_accumulatortype(f, args...) -> f(DefaultAccumulatorType(), args...)
with_default_projection(f, args...)    -> f(DefaultProjection(), args...)
\end{minted}

Call \texttt{f} with the named session default as its first argument — the supported way to \emph{consume} a default. While the default still holds its initial value (the overwhelmingly common case) the call is statically compiled against that constant: no dynamic dispatch, \texttt{f} fully specialized. After the default is changed, the call crosses a function barrier instead: one dynamic dispatch, everything inside fully specialized.

Allocation contract: the combinator{\textquotesingle}s return type unions both paths. When \texttt{f}{\textquotesingle}s result type does not depend on the default — e.g. \texttt{with\_default\_projection} with the formats fixed by the caller — the union is concrete and the call is \textbf{zero-allocation with a concretely inferred result}. When the result{\textquotesingle}s type \emph{is} the default (\texttt{with\_default\_type((T, x) -> T(x), …)}), the value is computed on the specialized path but boxes once at escape — the irreducible cost of a runtime-chosen type.


\begin{minted}{jlcon}
julia> with_default_type((T, x) -> T(x), 1.5)
Binary8p2se(1.5 ≡ 0x41)

julia> with_default_projection((ρ, x, y) -> Add(Binary8p4se, ρ, x, y),
                               Binary8p4se(1.5), Binary8p4se(0.25))
Binary8p4se(1.75 ≡ 0x46)
\end{minted}



\end{adjustwidth}
\hypertarget{5780687441881078027}{\texttt{SmallFloats.with\_default\_returntype}}  -- {Function.}

\begin{adjustwidth}{2em}{0pt}


\begin{minted}{julia}
with_default_type(f, args...)          -> f(DefaultType(), args...)
with_default_returntype(f, args...)    -> f(DefaultReturnType(), args...)
with_default_accumulatortype(f, args...) -> f(DefaultAccumulatorType(), args...)
with_default_projection(f, args...)    -> f(DefaultProjection(), args...)
\end{minted}

Call \texttt{f} with the named session default as its first argument — the supported way to \emph{consume} a default. While the default still holds its initial value (the overwhelmingly common case) the call is statically compiled against that constant: no dynamic dispatch, \texttt{f} fully specialized. After the default is changed, the call crosses a function barrier instead: one dynamic dispatch, everything inside fully specialized.

Allocation contract: the combinator{\textquotesingle}s return type unions both paths. When \texttt{f}{\textquotesingle}s result type does not depend on the default — e.g. \texttt{with\_default\_projection} with the formats fixed by the caller — the union is concrete and the call is \textbf{zero-allocation with a concretely inferred result}. When the result{\textquotesingle}s type \emph{is} the default (\texttt{with\_default\_type((T, x) -> T(x), …)}), the value is computed on the specialized path but boxes once at escape — the irreducible cost of a runtime-chosen type.


\begin{minted}{jlcon}
julia> with_default_type((T, x) -> T(x), 1.5)
Binary8p2se(1.5 ≡ 0x41)

julia> with_default_projection((ρ, x, y) -> Add(Binary8p4se, ρ, x, y),
                               Binary8p4se(1.5), Binary8p4se(0.25))
Binary8p4se(1.75 ≡ 0x46)
\end{minted}



\end{adjustwidth}
\hypertarget{7019104179409499666}{\texttt{SmallFloats.with\_default\_type}}  -- {Method.}

\begin{adjustwidth}{2em}{0pt}


\begin{minted}{julia}
with_default_type(f, args...)          -> f(DefaultType(), args...)
with_default_returntype(f, args...)    -> f(DefaultReturnType(), args...)
with_default_accumulatortype(f, args...) -> f(DefaultAccumulatorType(), args...)
with_default_projection(f, args...)    -> f(DefaultProjection(), args...)
\end{minted}

Call \texttt{f} with the named session default as its first argument — the supported way to \emph{consume} a default. While the default still holds its initial value (the overwhelmingly common case) the call is statically compiled against that constant: no dynamic dispatch, \texttt{f} fully specialized. After the default is changed, the call crosses a function barrier instead: one dynamic dispatch, everything inside fully specialized.

Allocation contract: the combinator{\textquotesingle}s return type unions both paths. When \texttt{f}{\textquotesingle}s result type does not depend on the default — e.g. \texttt{with\_default\_projection} with the formats fixed by the caller — the union is concrete and the call is \textbf{zero-allocation with a concretely inferred result}. When the result{\textquotesingle}s type \emph{is} the default (\texttt{with\_default\_type((T, x) -> T(x), …)}), the value is computed on the specialized path but boxes once at escape — the irreducible cost of a runtime-chosen type.


\begin{minted}{jlcon}
julia> with_default_type((T, x) -> T(x), 1.5)
Binary8p2se(1.5 ≡ 0x41)

julia> with_default_projection((ρ, x, y) -> Add(Binary8p4se, ρ, x, y),
                               Binary8p4se(1.5), Binary8p4se(0.25))
Binary8p4se(1.75 ≡ 0x46)
\end{minted}



\end{adjustwidth}

\part{Internal Reference}


\chapter{Internal Reference}



\label{2804742620609579687}{}


Unexported machinery documented in-source and referenced by the \hyperlinkref{4463396622184713734}{Technical Guide} and \hyperlinkref{4476403263922386643}{Technical Examples}. Public exported names are in the \hyperlinkref{5597696625397351643}{External Reference}.


\hypertarget{307234055475539643}{\texttt{SmallFloats.DEFAULT\_RBITS}}  -- {Constant.}

\begin{adjustwidth}{2em}{0pt}

Default random-bit budget used by no-argument stochastic projection constructors.



\end{adjustwidth}
\hypertarget{2700147405095606771}{\texttt{SmallFloats.ExactExternalFloat}}  -- {Type.}

\begin{adjustwidth}{2em}{0pt}

External float element types whose widening to the Float64 carrier is exact.



\end{adjustwidth}
\hypertarget{5600219088821016819}{\texttt{SmallFloats.MaybeR}}  -- {Type.}

\begin{adjustwidth}{2em}{0pt}

An explicit stochastic draw \texttt{R}, or \texttt{nothing} to draw one from the rng.



\end{adjustwidth}
\hypertarget{6952275749280347714}{\texttt{SmallFloats.NAN\_ORDER\_KEY}}  -- {Constant.}

\begin{adjustwidth}{2em}{0pt}

The order key reserved for the single NaN — above every finite key and ±Inf.



\end{adjustwidth}
\hypertarget{364714985955516448}{\texttt{SmallFloats.TERNARY\_ADAPTIVE\_BITS}}  -- {Constant.}

\begin{adjustwidth}{2em}{0pt}

K1+K2+K3 up to which a ternary table may be built adaptively once the signature has processed \texttt{TERNARY\_BUILD\_ELEMS[]} elements (default 21 bits = 2 MiB — the K=7 band). Above this, ternary ops always run the compute kernel.



\end{adjustwidth}
\hypertarget{3202729790529991172}{\texttt{SmallFloats.TERNARY\_BUILD\_ELEMS}}  -- {Constant.}

\begin{adjustwidth}{2em}{0pt}

Cumulative element count at which an adaptive-band signature earns its table.



\end{adjustwidth}
\hypertarget{7367898474622902585}{\texttt{SmallFloats.TERNARY\_CACHE\_BYTES}}  -- {Constant.}

\begin{adjustwidth}{2em}{0pt}

Byte budget for the ternary cache; least-recently-used tables evict first.



\end{adjustwidth}
\hypertarget{13776331714046560855}{\texttt{SmallFloats.TERNARY\_EAGER\_BITS}}  -- {Constant.}

\begin{adjustwidth}{2em}{0pt}

K1+K2+K3 up to which a ternary table is built eagerly on first array call (default 18 bits = 256 KiB — covers every all-K≤6 signature).



\end{adjustwidth}
\hypertarget{16271046097526486057}{\texttt{SmallFloats.THREADED\_KERNELS}}  -- {Constant.}

\begin{adjustwidth}{2em}{0pt}

Master switch for threaded ternary compute loops.



\end{adjustwidth}
\hypertarget{8230965136284310108}{\texttt{SmallFloats.THREAD\_MIN\_ELEMS}}  -- {Constant.}

\begin{adjustwidth}{2em}{0pt}

Minimum element count before the pure-ρ ternary compute loop threads.



\end{adjustwidth}
\hypertarget{8770213630193056212}{\texttt{SmallFloats.BigExactF}}  -- {Type.}

\begin{adjustwidth}{2em}{0pt}

Deferred exact result: \texttt{f()} returns the value as an exact \texttt{BigFloat}.



\end{adjustwidth}
\hypertarget{12947903703664656909}{\texttt{SmallFloats.Enclose128F}}  -- {Type.}

\begin{adjustwidth}{2em}{0pt}

Correctly-rounded Float128 bracket (plan Site C, Class R): IEEE mandates correct rounding for Float128 \texttt{/}, \texttt{sqrt}, \texttt{fma}, so a nearest-CR result of a value known inexact brackets the truth in the \emph{open} interval \texttt{(lo, hi) = (prevfloat(q), nextfloat(q))} with no envelope assumption. \texttt{f} is the MPFR fallback for (near-impossible at P ≤ 8) grid-straddling brackets.



\end{adjustwidth}
\hypertarget{8610954342460218668}{\texttt{SmallFloats.EncloseF}}  -- {Type.}

\begin{adjustwidth}{2em}{0pt}

Deferred enclosure, resolved by \texttt{\_finish} through up to three stages, cheapest first:

\begin{itemize}
\item[1. ] \texttt{yd} — an \emph{eager} Float64 estimate (NaN ⇒ absent) whose libm evaluation is faithful (≤ 1 ulp); the true value is taken to lie within \texttt{|yd|·2{\textasciicircum}-45} (\texttt{\_F64\_RELEXP}), ≥ 2{\textasciicircum}7 slacker than any faithful-libm error.


\item[2. ] \texttt{fq} — a zero-argument Float128 estimator (plan Site D, Class E; \texttt{nothing} ⇒ absent): true value within \texttt{|y|·2{\textasciicircum}-90} of \texttt{y = fq()} (\texttt{\_F128\_RELEXP}), a bound ≥ 2{\textasciicircum}18 slacker than any published libquadmath transcendental error, discharged empirically by the differential-build tests.


\item[3. ] \texttt{f(prec)} — the rigorous MPFR ladder: directed \texttt{(lo, hi)::NTuple\{2,BigFloat\}} strictly bracketing the true value.

\end{itemize}
Each estimate stage returns only when the two-sided sticky projection of its envelope agrees on a single code point; any disagreement (or an absent/degenerate estimate) falls through to the next stage. Stage 3 always decides.



\end{adjustwidth}
\hypertarget{11170301116400441493}{\texttt{SmallFloats.Rounded}}  -- {Type.}

\begin{adjustwidth}{2em}{0pt}

Result of ωRoundToPrecision in canonical integer form: a kind tag plus \texttt{sign · S · 2{\textasciicircum}Q} for finite results (design §5.2). The saturation stage consumes this without ever re-touching the carrier value.



\end{adjustwidth}
\hypertarget{16066611238377189300}{\texttt{SmallFloats.StickyF}}  -- {Type.}

\begin{adjustwidth}{2em}{0pt}

Sticky-head exact result (wide-spread tail, Float128 revision follow-on): the true value is \texttt{v + sgn·ε} for an infinitesimal \texttt{ε > 0} — \texttt{v} carries every bit the projection can consume and \texttt{sgn ∈ \{-1,+1\}} the direction of the neglected tail. Sound because \texttt{v} is either exactly on a rounding threshold of the target grid (where \texttt{sticky} decides, for every mode including stochastic sub-grids) or strictly farther from the nearest threshold than the tail magnitude; the emitting sites in oracle.jl discharge that bound (operand significands ≤ 17 bits, target grids ≤ P−1+N ≤ 67 fractional bits, spreads > the \emph{DE}* thresholds). Non-allocating: replaces the BigFloat escalation for FMA/FAA.



\end{adjustwidth}
\hypertarget{261438848462006089}{\texttt{SmallFloats.TableKey}}  -- {Type.}

\begin{adjustwidth}{2em}{0pt}

Value key for the table cache: op + format parameters + ρ parameters (all values, so the cache itself is type-stable and non-allocating to query).



\end{adjustwidth}
\hypertarget{3589516491715718078}{\texttt{SmallFloats.TernaryKey}}  -- {Type.}

\begin{adjustwidth}{2em}{0pt}

Value key for a ternary table: op + result + three operand formats + ρ.



\end{adjustwidth}
\hypertarget{3431767106840558017}{\texttt{Core.Type}}  -- {Method.}

\begin{adjustwidth}{2em}{0pt}


\begin{minted}{julia}
(T::Type{<:Binary})(c::UInt8) -> T
\end{minted}

Construct a value from its \textbf{code point} (validated: for K < 8, \texttt{c} must be \texttt{< 2{\textasciicircum}K} or an \texttt{ArgumentError} is thrown; the range check folds away at K = 8 and for constant codes). \texttt{Binary5p3sf(0x08) == Binary5p3sf(1.0)}.

\texttt{UInt8} is the \emph{only} argument type with code-point semantics — every other \texttt{Real} (including other \texttt{Integer}s) constructs by projecting the numeric value, and \texttt{Convert} is numeric for all integers: \texttt{Binary8p4se(0x02)} is code point 2 (the second-smallest subnormal), while \texttt{Binary8p4se(2)} is the number 2.0. \texttt{rawvalue(T, c)} remains the unchecked kernel-internal route.



\end{adjustwidth}
\hypertarget{17928462762263157549}{\texttt{SmallFloats.\_bp\_element}}  -- {Method.}

\begin{adjustwidth}{2em}{0pt}

One element of ωBlockProject: the draft{\textquotesingle}s S-special rows, then ωDivide + ωProject.

Fast-path cascade after the special rows, ordered cheapest-first by result kind: exact Float64 quotient → CR Float128 quotient (exact or one-ulp bracket) → CR-divided Enclose128F bracket → fq-composition envelope (2{\textasciicircum}-89) — each resolved by the two-sided sticky gate; any miss falls through to the rigorous MPFR interval at the bottom.



\end{adjustwidth}
\hypertarget{15953021384962181472}{\texttt{SmallFloats.\_check\_tabulable}}  -- {Method.}

\begin{adjustwidth}{2em}{0pt}

Stochastic ρ is a distribution over R, never a table (design §5.4).



\end{adjustwidth}
\hypertarget{12008461486952181445}{\texttt{SmallFloats.\_comparable}}  -- {Method.}

\begin{adjustwidth}{2em}{0pt}

Numeric comparison is defined only when neither operand is NaN — every \texttt{==}/\texttt{<}/\texttt{<=} below returns \texttt{false} otherwise, per IEEE unorderedness.



\end{adjustwidth}
\hypertarget{7363307359868208211}{\texttt{SmallFloats.\_resolve\_rng}}  -- {Method.}

\begin{adjustwidth}{2em}{0pt}

Resolve a caller-supplied rng, falling back to the task-local default. Call sites hoist this out of loops so array kernels resolve once per call, not per element.



\end{adjustwidth}
\hypertarget{13712544213067904585}{\texttt{SmallFloats.\_rng\_for}}  -- {Method.}

\begin{adjustwidth}{2em}{0pt}

The rng an operation under ρ will actually draw from — \texttt{nothing} for pure ρ, which must never touch RNG state.



\end{adjustwidth}
\hypertarget{17911291590751857121}{\texttt{SmallFloats.\_ternary\_table\_for}}  -- {Method.}

\begin{adjustwidth}{2em}{0pt}


\begin{minted}{julia}
_ternary_table_for(op, fr, f1, f2, f3, ρ, nelems) -> Union{Nothing,Memory{UInt8}}
\end{minted}

The kernel-facing policy gate, called once per array operation with the call{\textquotesingle}s element count. Eager band (Σ Kᵢ ≤ \texttt{TERNARY\_EAGER\_BITS[]}): fetch/build now. Adaptive band (≤ \texttt{TERNARY\_ADAPTIVE\_BITS[]}): return the cached table if present, otherwise accumulate \texttt{nelems} against the signature and build only once it has earned \texttt{TERNARY\_BUILD\_ELEMS[]}. Beyond the adaptive band (the K=8 territory): always \texttt{nothing} — the compute kernel is the right tradeoff there.



\end{adjustwidth}
\hypertarget{8539869125382303808}{\texttt{SmallFloats.apply\_op}}  -- {Method.}

\begin{adjustwidth}{2em}{0pt}

apply\_op(Val(op), fr, ρ, R, xs...) — evaluate \texttt{op}{\textquotesingle}s ω-semantics on decoded Float64 operands and project the result into \texttt{fr} under ρ (with random bits R).



\end{adjustwidth}
\hypertarget{994737666889680542}{\texttt{SmallFloats.blockdecode}}  -- {Method.}

\begin{adjustwidth}{2em}{0pt}

ωBlockDecode (draft §5.1.1): per lane ωMultiply(decode(s), decode(xᵢ)) — exact Float64.



\end{adjustwidth}
\hypertarget{445456126373273252}{\texttt{SmallFloats.blockproject}}  -- {Method.}

\begin{adjustwidth}{2em}{0pt}

ωBlockProject (draft §5.1.2) over a tuple of result kinds, scale supplied as a value.



\end{adjustwidth}
\hypertarget{9707793631963955151}{\texttt{SmallFloats.encode}}  -- {Method.}

\begin{adjustwidth}{2em}{0pt}


\begin{minted}{julia}
encode(T, sign, S, Q) -> UInt8   (private; design §3.3)
\end{minted}

ωEncode from canonical integer form: value = sign · S · 2{\textasciicircum}Q, with S ∈ 0:2{\textasciicircum}P (S == 2{\textasciicircum}P is the next-binade carry, draft §4.7.4 NOTE 4). Precondition: the value is in the datum set of \texttt{T} (guaranteed by RoundToPrecision ∘ Saturate).



\end{adjustwidth}
\hypertarget{10413485981529569906}{\texttt{SmallFloats.get\_table}}  -- {Function.}

\begin{adjustwidth}{2em}{0pt}


\begin{minted}{julia}
get_table(op, fr, f1, [f2,] ρ) -> Memory{UInt8}
\end{minted}

Fetch (building and caching on first use) the complete result table for the pure-ρ specialization \texttt{op⟨fr; f1[, f2]⟩ under ρ}: entry \texttt{c + 1} (unary) or \texttt{(c1 << K2) + c2 + 1} (binary) holds the result code point for those operand code points. Throws for stochastic ρ. \texttt{@noinline} by design: kernels call this once per array operation and index the returned table in their hot loop.



\end{adjustwidth}
\hypertarget{7588185958709694166}{\texttt{SmallFloats.get\_table}}  -- {Method.}

\begin{adjustwidth}{2em}{0pt}


\begin{minted}{julia}
get_table(op, fr, f1, f2, f3, ρ) -> Memory{UInt8}
\end{minted}

Ternary form: entry \texttt{((c1 << K2 | c2) << K3) + c3 + 1} holds the result code point. Builds unconditionally (policy lives in \texttt{\_ternary\_table\_for}); pure ρ only.



\end{adjustwidth}
\hypertarget{16932623480385077633}{\texttt{SmallFloats.project}}  -- {Method.}

\begin{adjustwidth}{2em}{0pt}


\begin{minted}{julia}
project(T, ρ, X; R=0, sticky=0) -> T
\end{minted}

Project a closed-extended-real carrier value into format \texttt{T} under ρ. \texttt{X::Float64} must be \emph{exact} (the Float64 fast path); \texttt{X::BigFloat} likewise (or an interval endpoint with \texttt{sticky ∈ \{-1,+1\}}). \texttt{R} is the stochastic draw.



\end{adjustwidth}
\hypertarget{1151381039512186337}{\texttt{SmallFloats.project\_interval}}  -- {Method.}

\begin{adjustwidth}{2em}{0pt}


\begin{minted}{julia}
project_interval(T, ρ, f; R=0, maxprec=4096) -> T
\end{minted}

Project the true value of \texttt{f}, where \texttt{f(prec)} returns an enclosure \texttt{(d::BigFloat, u::BigFloat)} with the true value in \texttt{[d, u]} (endpoints from MPFR directed rounding). Rigor (design §9.2): if d == u the value is exact — project it. Otherwise the true value lies in the \emph{open} interval (d, u) (directed correct rounding of a non-representable value moves strictly); projection is monotone in the value for every mode at fixed R, so if project(d, sticky=+1) == project(u, sticky=-1) that common code is the answer; otherwise a projection grid point sits inside the interval and precision escalates.



\end{adjustwidth}
\hypertarget{835860055659755424}{\texttt{SmallFloats.unpack\_tile!}}  -- {Method.}

\begin{adjustwidth}{2em}{0pt}

Unpack \texttt{pv[first:first+len-1]} into \texttt{buf} (byte scratch); tile-granular kernel entry.



\end{adjustwidth}

\end{document}
